\documentclass[12pt,letterpaper]{article}

% Documento de trabajo de grado compatible con Overleaf y compilación local.
% Compilación recomendada: pdflatex -> bibtex -> pdflatex -> pdflatex.
\usepackage[spanish,es-nodecimaldot]{babel}
\usepackage[utf8]{inputenc}
\usepackage[T1]{fontenc}
\usepackage{lmodern}
\usepackage[margin=1in]{geometry}
\usepackage{setspace}
\usepackage{amsmath}
\usepackage{amssymb}
\usepackage{booktabs}
\usepackage{longtable}
\usepackage{array}
\usepackage{tabularx}
\usepackage{enumitem}
\usepackage{float}
\usepackage{pdflscape}
\usepackage{afterpage}
\usepackage{adjustbox}
\usepackage{tikz}
\usetikzlibrary{arrows.meta,positioning,shapes.geometric}
\usepackage[hidelinks,pdfpagelabels,plainpages=false]{hyperref}
\usepackage[natbibapa]{apacite}
\bibliographystyle{apacite}
\providecommand{\parencite}{\citep}
\providecommand{\textcite}{\citet}

\setcounter{secnumdepth}{3}
\setcounter{tocdepth}{3}
\setlength{\parindent}{0.5in}
\setlength{\parskip}{0pt}
\doublespacing
\sloppy
\urlstyle{same}

\renewcommand{\contentsname}{TABLA DE CONTENIDO}
% Coherencia terminologica (17-08-2026). babel-spanish reasigna
% \tablename a "Cuadro" y \listtablename a "Indice de cuadros" DESPUES
% del preambulo, de modo que un \renewcommand suelto queda sobrescrito: el
% indice general decia "LISTA DE TABLAS" pero la pagina de listado decia
% "Indice de cuadros". Las tres redefiniciones van dentro de
% \addto\captionsspanish, que es donde babel las respeta, y se unifica
% todo en "Tabla".
\addto\captionsspanish{%
  \renewcommand{\tablename}{Tabla}%
  \renewcommand{\listtablename}{LISTA DE TABLAS}%
  \renewcommand{\listfigurename}{LISTA DE FIGURAS}%
}
\newcolumntype{Y}{>{\raggedright\arraybackslash}X}

\newcommand{\ICOCIV}{ICOCIV}
\newcommand{\ICCP}{ICCP}
\newcommand{\DANE}{DANE}
\newcommand{\IDU}{IDU}
\newcommand{\INVIAS}{INVIAS}

\title{Desarrollo de una herramienta de software para el análisis de variaciones y la proyección de precios en obras civiles en la Secretaría de Infraestructura de la Gobernación del Cauca}
\author{Carlos Esteban Ojeda Calvache}
\date{Universidad del Cauca\\Facultad de Ingeniería Civil\\Popayán, Colombia\\2026}

\begin{document}
\hypersetup{pageanchor=false}

\begin{titlepage}
    \centering
    \vspace*{1.2cm}
    {\Large\bfseries Desarrollo de una herramienta de software para el análisis de variaciones y la proyección de precios en obras civiles en la Secretaría de Infraestructura de la Gobernación del Cauca\par}
    \vspace{1.2cm}
    {\large Trabajo de grado modalidad investigación\par}
    \vspace{1.4cm}
    {\large Carlos Esteban Ojeda Calvache\par}
    \vspace{0.3cm}
    {\large Código: 100416021068\par}
    \vspace{1cm}
    {\large Director: PhD. Lucio Gerardo Cruz Velasco\par}
    \vspace{1.2cm}
    {\large Universidad del Cauca\par}
    {\large Facultad de Ingeniería Civil\par}
    {\large Programa de Ingeniería Civil\par}
    \vspace{0.8cm}
    {\large Popayán, Colombia\par}
    {\large 2026\par}
    \vfill
\end{titlepage}
\hypersetup{pageanchor=true}

\pagenumbering{roman}
\tableofcontents
\newpage
\phantomsection
\addcontentsline{toc}{section}{LISTA DE TABLAS}
\listoftables
\newpage
\phantomsection
\addcontentsline{toc}{section}{LISTA DE FIGURAS}
\listoffigures
\newpage

\section*{LISTA DE ACRÓNIMOS}
\addcontentsline{toc}{section}{LISTA DE ACRÓNIMOS}

\noindent\begin{tabularx}{\textwidth}{@{}lY@{}}
\toprule
\textbf{Sigla} & \textbf{Descripción} \\
\midrule
SAVIP & Sistema de Análisis de Variaciones de Precios (aplicación desarrollada en este trabajo). \\
DANE & Departamento Administrativo Nacional de Estadística. \\
ICOCIV & Índice de Costos de la Construcción de Obras Civiles. \\
ICCP & Índice de Costos de la Construcción Pesada. \\
IDU & Instituto de Desarrollo Urbano. \\
INVIAS & Instituto Nacional de Vías. \\
APU & Análisis de Precios Unitarios. \\
MAE & Error absoluto medio. \\
RMSE & Raíz del error cuadrático medio. \\
MAPE & Error porcentual absoluto medio. \\
sMAPE & Error porcentual absoluto medio simétrico. \\
MASE & Error absoluto medio escalado. \\
ETS & Familia de modelos de suavizamiento exponencial formulada en espacio de estados con componentes de error, tendencia y estacionalidad. \\
DOCX & Formato de documento de Microsoft Word. \\
\bottomrule
\end{tabularx}

\newpage
\pagenumbering{arabic}

\section{INTRODUCCIÓN}

\subsection{Contexto}

La gestión de variaciones de precios en contratos de obra civil exige procedimientos técnicos verificables, especialmente cuando los análisis se sustentan en índices oficiales de costos. En Colombia, el \DANE{} produce información estadística asociada a los costos de construcción y, para el caso de obras civiles, publica el \ICOCIV{} como operación estadística orientada a medir la variación promedio de precios de bienes y servicios usados en este sector \parencite{dane_ficha_icociv,dane_metodologia_icociv}. Antes de la entrada en operación del \ICOCIV, el \ICCP{} cumplió un papel relevante en el seguimiento de costos de construcción pesada.

En la Secretaría de Infraestructura de la Gobernación del Cauca, el análisis de variaciones y proyecciones de precios requiere consultar fuentes oficiales, organizar series mensuales, seleccionar rutas jerárquicas de índices, aplicar criterios de transición entre \ICCP{} e \ICOCIV{} y documentar los resultados de forma reproducible. El proyecto desarrolla una aplicación de apoyo técnico, denominada \textbf{SAVIP} (Sistema de Análisis de Variaciones de Precios), que integra estas tareas en un flujo único: carga de anexos \ICOCIV, selección jerárquica, análisis histórico, proyección estadística, cálculo de empalme \ICCP--\ICOCIV, generación de reportes y exportación de resultados. A lo largo del documento, SAVIP designa la aplicación desarrollada, mientras que \ICOCIV{} e \ICCP{} designan los índices oficiales del \DANE{} que la aplicación analiza.

La herramienta no reemplaza el criterio técnico, jurídico ni contractual del profesional responsable. Su finalidad es ordenar datos, reducir manipulación manual, dejar trazabilidad de índices \(I\) e \(I_0\), fórmulas, rutas, fechas y resultados, y facilitar la revisión de cálculos asociados a variaciones de precios.

\subsection{Planteamiento del problema}

Los procesos de análisis de variaciones de precios en obras civiles pueden apoyarse en hojas de cálculo, anexos oficiales, documentos técnicos y criterios profesionales. Sin embargo, cuando la información se organiza manualmente, aumenta el riesgo de errores por digitación, selección incorrecta de rangos, arrastre de fórmulas, actualización parcial de índices o pérdida de trazabilidad sobre los valores empleados. Esta situación se vuelve más sensible cuando el análisis cruza dos fuentes históricas: el \ICCP, usado hasta el periodo de transición, y el \ICOCIV, definido como referente posterior para obras civiles.

La necesidad técnica identificada se relaciona con tres problemas principales. Primero, la consulta de índices oficiales requiere organizar datos mensuales y niveles jerárquicos de forma consistente. Segundo, la proyección de índices demanda criterios estadísticos documentados y validación de resultados. Tercero, el empalme \ICCP--\ICOCIV{} exige aplicar fórmulas de transición y conservar evidencia de las decisiones tomadas, en especial cuando se seleccionan equivalencias entre series.

Por tanto, el problema abordado consiste en diseñar e implementar una aplicación que permita consultar, analizar, proyectar y documentar variaciones de precios con base en información oficial, reduciendo tareas manuales y mejorando la reproducibilidad del cálculo sin sustituir la revisión técnica del profesional responsable.

\subsection{Objetivos}

\subsubsection{Objetivo general}

Desarrollar una herramienta de software para el análisis de variaciones y proyecciones de precios en obras civiles, incluyendo aquellas gestionadas por la Secretaría de Infraestructura de la Gobernación del Cauca, mediante el uso de información histórica oficial del ICCP e ICOCIV reportada por el DANE, con el propósito de generar resultados técnicamente consistentes de acuerdo con el comportamiento y las limitaciones de las series analizadas.

\subsubsection{Objetivos específicos}

\begin{itemize}
    \item Diseñar e implementar una herramienta de software que integre funcionalidades de consulta, filtrado, visualización y procesamiento de información histórica oficial del ICCP e ICOCIV para el análisis técnico de variaciones de precios en obras civiles.
    \item Incorporar funcionalidades de análisis y proyección basadas en series históricas que permitan generar resultados técnicamente interpretables de acuerdo con el comportamiento y las limitaciones de las series analizadas, para apoyar procesos de evaluación económica y planeación de proyectos de obra civil.
    \item Evaluar distintos enfoques de análisis y procesamiento estadístico sobre la información histórica utilizada por la herramienta, con el fin de identificar comportamientos consistentes y limitaciones asociadas a las series analizadas dentro del contexto de aplicación definido.
    \item Realizar pruebas de funcionamiento y evaluación de resultados utilizando información histórica representativa, con el propósito de analizar el comportamiento de la herramienta y la consistencia de las salidas generadas de acuerdo con las características y limitaciones de las series procesadas.
    \item Elaborar un manual de usuario que documente el funcionamiento general del software, los procedimientos de uso y las recomendaciones básicas para su operación.
\end{itemize}

\subsection{Justificación}

El análisis de variaciones de precios en obras civiles depende de datos oficiales, criterios técnicos y documentación suficiente para soportar decisiones contractuales o presupuestales. Una aplicación que integre consulta de índices, análisis de series, proyección, empalme y exportación de resultados aporta valor porque reduce operaciones manuales repetitivas y facilita la revisión de cada cálculo.

La revisión de una herramienta previa en Microsoft Excel permitió reconocer un flujo operativo útil para la entidad: organización de índices, uso de fórmulas, proyección de valores y estructuración de salidas. No obstante, al depender de actualizaciones manuales, gráficas generadas en hoja de cálculo, lectura de parámetros, copia de ecuaciones y arrastre de fórmulas, el proceso presenta riesgos asociados a trazabilidad, reproducibilidad y control de errores. La aplicación desarrollada responde a esa necesidad mediante validaciones, selección controlada de rutas, registro de valores \(I\) e \(I_0\), diferenciación entre cálculo general y cálculo especial para acero, y generación de salidas documentales.

Además, la transición entre \ICCP{} e \ICOCIV{} requiere que las equivalencias seleccionadas sean explícitas y justificables. Por ello, el sistema no automatiza el criterio profesional de equivalencia; lo registra y lo presenta para revisión. Este enfoque conserva el papel del profesional responsable y concentra el aporte del software en la organización, cálculo y trazabilidad.

\subsection{Alcance y limitaciones}

El alcance del proyecto comprende el desarrollo de una aplicación para cargar anexos \ICOCIV, consultar series históricas, ejecutar análisis y proyecciones, aplicar cálculos de empalme \ICCP--\ICOCIV, generar reportes y exportar resultados. El sistema trabaja con información oficial o archivos cargados por el usuario y conserva trazabilidad de las rutas, periodos, índices y fórmulas usados.

El proyecto no pretende crear una metodología estadística oficial nueva, reemplazar los documentos técnicos del \DANE, definir equivalencias contractuales automáticas ni sustituir la revisión jurídica o presupuestal de una entidad. La confiabilidad de las proyecciones depende de la calidad, continuidad y longitud de las series disponibles; por ello, la aplicación incorpora advertencias, métricas y criterios de horizonte viable.

\section[MARCO TEÓRICO, CONTEXTO Y ANTECEDENTES]{MARCO TEÓRICO, CONTEXTO Y\\ ANTECEDENTES}

\subsection{Marco conceptual}

El análisis de variación de precios en obras civiles se apoya en la comparación de valores observados en distintos periodos. Para que esa comparación sea técnicamente revisable, el cálculo debe identificar la fuente del índice, el periodo inicial, el periodo final, la ruta de clasificación usada y la fórmula aplicada. En este trabajo se entiende el índice de costos como un número relativo que resume la variación de precios de una canasta de bienes o servicios respecto a un periodo de referencia. Por tanto, el índice no representa por sí mismo un precio contractual, sino una medida de cambio que puede servir como insumo para fórmulas de actualización.

El periodo base corresponde al momento frente al cual se expresan los valores de una serie. En el caso del \ICOCIV, la metodología del \DANE{} define una estructura de publicación que permite consultar resultados por niveles jerárquicos y periodos mensuales \parencite{dane_ficha_icociv,dane_metodologia_icociv}. La variación mensual describe el cambio entre un mes y el inmediatamente anterior; la variación acumulada resume cambios dentro de un intervalo mayor. Para el sistema desarrollado, el valor central no es solamente la variación publicada, sino la trazabilidad entre el índice inicial \(I_0\), el índice final \(I\), el tramo temporal y la ruta seleccionada.

El análisis de precios unitarios (APU) organiza cantidades, unidades, rendimientos e insumos que conforman un valor contractual. Cuando se revisan variaciones de precios, el APU ayuda a identificar el insumo contractual que se desea actualizar y la unidad asociada. La aplicación desarrollada no modifica el APU; toma el valor base informado por el usuario y lo relaciona con índices oficiales para apoyar el cálculo de actualización.

La trazabilidad del cálculo se entiende como la capacidad de reconstruir el resultado a partir de entradas, fuentes, fechas, índices, fórmulas y observaciones. En este proyecto, dicha trazabilidad es un requisito metodológico: permite revisar si la serie \ICCP{} o \ICOCIV{} seleccionada corresponde al insumo analizado, si las fechas pertenecen al tramo correcto y si el resultado fue obtenido por fórmula general o por el caso especial del acero.

\subsection{Contexto institucional}

La Secretaría de Infraestructura de la Gobernación del Cauca gestiona información asociada a proyectos de infraestructura, contratos, presupuestos y análisis técnicos de obra civil. En ese contexto, una herramienta de apoyo debe facilitar la consulta de fuentes oficiales y la documentación de resultados, sin reemplazar las competencias del profesional responsable. El software desarrollado se ubica en ese punto: organiza datos, aplica reglas de cálculo programadas y presenta resultados verificables para revisión técnica.

El desarrollo de una herramienta de software para el análisis de variaciones y la proyección de precios en obras civiles permite vincular tres necesidades institucionales: consulta ordenada de índices, reducción de operaciones manuales y generación de salidas que puedan revisarse posteriormente. La aplicación responde a estas necesidades mediante módulos separados de carga \ICOCIV, selección jerárquica, análisis de serie, proyección, empalme \ICCP--\ICOCIV{} y exportación.

El sistema debe interpretarse como una herramienta de apoyo técnico. La decisión sobre procedencia contractual, equivalencia de insumos, reconocimiento de ajustes o aplicación en un caso específico requiere revisión profesional, jurídica y administrativa. Esta delimitación evita presentar el software como una autoridad automática y mantiene su alcance dentro del procesamiento, organización y trazabilidad de información.

\subsection{Contratos de obra civil y variación de precios}

En los contratos de obra civil, los cambios de precios pueden afectar el equilibrio de presupuestos, la ejecución financiera y la interpretación de valores pactados. La revisión de precios exige soportes verificables, porque el cálculo no depende únicamente de una operación aritmética: también requiere identificar el insumo, el periodo, la fuente del índice y la metodología aplicable. La literatura institucional sobre equilibrio económico, revisión y ajuste de precios resalta la necesidad de analizar cada caso con soporte documental y criterios contractuales \parencite{cce2024_equilibrio}.

Dentro de un análisis técnico, los índices oficiales aportan una referencia externa para medir cambios de costos. Sin embargo, el uso de un índice no elimina la necesidad de justificar su pertinencia frente al insumo contractual. Por ejemplo, un material, una actividad o un grupo de obra puede tener más de una posible equivalencia entre estructuras \ICCP{} e \ICOCIV. Por ello, la aplicación registra la selección realizada por el usuario y la presenta en tablas de equivalencia, dejando claro que la equivalencia requiere criterio técnico.

La variación de precios también puede requerir separar tramos temporales. Cuando una fecha inicial se ubica antes de diciembre de 2021 y una fecha final se ubica después de ese periodo, el cálculo debe distinguir el tramo asociado al \ICCP{} y el tramo posterior asociado al \ICOCIV. Este problema explica la importancia del módulo de empalme y de la trazabilidad de índices \(I\) e \(I_0\).

\subsection{Índices de costos de construcción}

\subsubsection{Índice de Costos de Construcción Pesada (ICCP)}

El \ICCP{} fue una referencia histórica para medir variaciones de costos en construcción pesada. De acuerdo con la documentación metodológica del \DANE, antes del \ICOCIV{} existían mediciones orientadas al seguimiento de costos en obras de infraestructura, con estructuras de grupos de obra y grupos de costo \parencite{dane_ficha_icociv,dane_metodologia_icociv}. En este proyecto, el \ICCP{} se usa como fuente histórica para periodos anteriores o iguales a diciembre de 2021.

La aplicación conserva el Anexo 10 \ICCP{} histórico como fuente interna. Esta decisión reduce la carga operativa del usuario y evita que cada cálculo dependa de una importación manual del archivo histórico. Además, permite separar las series \ICCP{} por tipo: Total \ICCP, Canasta general y Grupo de obra. Esa separación es necesaria porque cada categoría cumple un papel diferente y no debe mezclarse en el selector de cálculo.

\subsubsection{Índice de Costos de Construcción de Obras Civiles (ICOCIV)}

El \ICOCIV{} es una operación estadística del \DANE{} orientada a medir la variación promedio de precios de una canasta de bienes y servicios usados en la construcción de obras civiles \parencite{dane_ficha_icociv,dane_metodologia_icociv}. Su diseño metodológico amplía la capacidad de desagregación para diferentes niveles de análisis, lo que permite consultar grupos de obra, subclases, tipologías, capítulos constructivos, grupos de costo e insumos, según la tabla disponible en los anexos oficiales.

Para la aplicación, esta estructura jerárquica es central. El usuario no selecciona solamente una columna de datos; selecciona una ruta dentro de la estructura \ICOCIV. Esa ruta se registra porque condiciona la serie histórica, la proyección y la equivalencia con el insumo contractual. Así, la trazabilidad no se limita al valor numérico del índice, sino que incluye el camino jerárquico usado para obtenerlo.

El \ICOCIV{} se usa en dos procesos del sistema. Primero, en el módulo de análisis y proyección, donde la serie histórica permite evaluar modelos y estimar valores futuros si el horizonte es viable. Segundo, en el módulo de empalme, donde el índice \ICOCIV{} inicial y final permiten calcular el segundo tramo de actualización después de la transición desde \ICCP.

\subsubsection{Transición ICCP--ICOCIV}

La transición entre \ICCP{} e \ICOCIV{} surge por el cambio metodológico y operativo en la medición de costos de obras civiles. Los lineamientos institucionales revisados reconocen que el \ICOCIV{} pasa a ser el referente posterior y que los usuarios deben realizar ajustes para aplicar sus resultados según sus necesidades \parencite{idu_circular_6_2022,invias2022_icociv}. En el sistema, esta transición se representa mediante un corte en diciembre de 2021.

El tratamiento implementado distingue tres escenarios. Si las fechas se ubican hasta diciembre de 2021, el cálculo usa solamente \ICCP. Si ambas fechas son posteriores a ese periodo, el cálculo usa solamente \ICOCIV. Si el rango cruza el periodo de transición, el cálculo se divide en dos tramos: primero \ICCP{} hasta diciembre de 2021 y luego \ICOCIV{} desde diciembre de 2021 hasta la fecha final. Esta regla permite aplicar las fórmulas de transición sin mezclar periodos ni fuentes.

\subsection{Metodología IDU para transición de índices}

\subsubsection{Transición general ICCP--ICOCIV}

Para la transición general, el documento de lineamientos establece un primer ajuste con el índice \ICCP{} y un segundo ajuste con el \ICOCIV{} \parencite{idu_circular_6_2022}. En el sistema, la base ajustable se define como \(P-A\), donde \(P\) corresponde al valor base a ajustar y \(A\) al anticipo amortizado.

\begin{equation}
\label{eq:r1_general}
R_1 = (P - A) \left[ \left( \frac{I}{I_0} \right) - 1 \right]
\end{equation}

Donde \(R_1\) es el primer valor parcial ajustado con \ICCP, \(I\) es el índice final del tramo \ICCP{} e \(I_0\) es el índice inicial del mismo tramo.

\begin{equation}
\label{eq:r2_general}
R_2 = [(P - A) + R_1] \left[ \left( \frac{I}{I_0} \right) - 1 \right]
\end{equation}

Donde \(R_2\) es el segundo valor parcial ajustado con \ICOCIV, calculado sobre la base ya ajustada por el tramo anterior.

\begin{equation}
\label{eq:r_total}
R = R_1 + R_2
\end{equation}

\begin{equation}
\label{eq:valor_actualizado_general}
\text{Valor actualizado} = (P - A) + R
\end{equation}

Estas fórmulas se implementan de forma trazable en el módulo de empalme. El resultado \(R\) corresponde a la suma de los valores parciales de ajuste. El valor actualizado incorpora la base ajustable más el ajuste total. La equivalencia entre la serie \ICCP{} seleccionada y la ruta \ICOCIV{} equivalente debe ser validada por el profesional responsable.

\subsubsection{Caso especial del acero}

El caso especial del acero tiene tratamiento diferenciado en los lineamientos revisados \parencite{idu_circular_6_2022}. En este caso se usa una base \(P_0\) asociada al valor del insumo acero. El primer ajuste se calcula con \ICCP:

\begin{equation}
\label{eq:r1_acero}
R_1 = P_0 \left[ \left( \frac{I}{I_0} \right) - 1 \right]
\end{equation}

El segundo ajuste se calcula con \ICOCIV{} sobre la base \(P_0 + R_1\):

\begin{equation}
\label{eq:r2_acero}
R_2 = (P_0 + R_1) \left[ \left( \frac{I}{I_0} \right) - 1 \right]
\end{equation}

\begin{equation}
\label{eq:r_acero}
R = R_1 + R_2
\end{equation}

Finalmente, para la comparación del valor facturado frente al valor ajustado se usa:

\begin{equation}
\label{eq:z_acero}
Z = (I_x \times q) - (R + P_0)
\end{equation}

Donde \(I_x\) es el valor facturado por kilogramo, \(q\) es la cantidad ejecutada y \(Z\) representa el valor adicional por fluctuación del acero. La aplicación presenta esta ruta como cálculo especial y conserva las variables separadas para que el usuario pueda revisar la diferencia frente al cálculo general.

\subsubsection{Selección de la fecha base para el índice inicial}
\label{sec:seleccion_i0}

Las fórmulas de ajuste dependen de la razón \(I/I_0\), por lo que la fecha asociada al índice inicial \(I_0\) determina directamente la magnitud del ajuste. La elección de esa fecha base no es un dato calculado por el sistema, sino una decisión contractual que la aplicación recibe como parámetro; por ello conviene precisar qué establecen las fuentes de referencia y qué ocurre en la práctica de los contratos de obra.

Las fuentes institucionales sugieren momentos determinados para fijar la fecha base. El Manual de Interventoría de Obra Pública del \INVIAS{} asocia esa referencia al cierre del proceso de selección, entendido como el instante en que quedan definidos los precios de la oferta \parencite{invias2010_interventoria}. En la misma línea, el procedimiento de implementación del \ICOCIV{} en los contratos de obra del \INVIAS{} describe cómo aplicar el índice para el ajuste, tomando como punto de partida las condiciones económicas iniciales del contrato \parencite{invias2022_icociv}. Estas referencias ofrecen un criterio uniforme y verificable para definir \(I_0\).

La fecha base aplicable depende, en cambio, de las condiciones y disposiciones del contrato particular; SAVIP no la determina automáticamente y exige que sea ingresada y validada por el profesional responsable. Esa dependencia es relevante porque la fórmula de ajuste es sensible a la fecha elegida: si \(I_0\) se fija en un mes muy próximo a los momentos sugeridos por las fuentes de referencia, la razón \(I/I_0\) tiende a ser menor y el ajuste reconocido resulta más bajo, lo que puede afectar el equilibrio económico del contrato, principio reconocido por la normativa y la doctrina de contratación pública colombiana \parencite{cce2024_equilibrio}.

Debe entenderse, entonces, que la fecha sugerida por las fuentes de referencia constituye un criterio orientador y no una regla universal de obligatorio cumplimiento: el momento exacto para definir \(I_0\) depende del acuerdo entre las partes dentro de lo que permita cada contrato. La aplicación no impone esa fecha ni la deduce automáticamente; conserva de forma explícita la fecha y el índice \(I_0\) empleados, de modo que el cálculo pueda revisarse y justificarse frente a lo pactado. La equivalencia entre la serie \ICCP{} o \ICOCIV{} seleccionada y el insumo contractual, así como la fecha base adoptada, quedan bajo el criterio del profesional responsable.

\subsection{Series de tiempo aplicadas a índices de costos}

Una serie de tiempo organiza observaciones de una variable en orden cronológico. En este proyecto, cada serie mensual \ICOCIV{} contiene valores de índice asociados a una ruta jerárquica. El análisis de la serie permite reconocer continuidad temporal, longitud disponible, tendencia, cambios abruptos y posibles valores atípicos. Estas condiciones son relevantes porque una proyección no debe interpretarse igual cuando la serie es estable y continua que cuando tiene pocos datos o saltos temporales.

La proyección de índices se usa como apoyo cuando el usuario requiere un valor futuro que todavía no se encuentra en el archivo oficial cargado. La aplicación no inventa índices ni reemplaza la publicación oficial del \DANE. Cuando se proyecta, el resultado queda identificado como estimación y debe interpretarse junto con el modelo seleccionado, las métricas de validación, el horizonte y las advertencias. Esta delimitación es coherente con la literatura de pronóstico, que recomienda evaluar los modelos fuera de muestra y no confiar únicamente en el ajuste histórico \parencite{hyndman2021fpp3}.

\subsection{Preparación y validación de series mensuales}

La posición temporal forma parte del dato. Por ello, antes de modelar se normaliza cada periodo al formato año--mes, se ordena cronológicamente y se comprueba que dos observaciones consecutivas estén separadas por un mes. Si una serie pasa de enero a marzo sin registrar febrero, el cambio observado corresponde a dos meses y no debe tratarse como variación mensual. Del mismo modo, dos valores para un mismo periodo generan ambigüedad sobre cuál representa la observación oficial.

La preparación revisa continuidad, duplicados, valores faltantes, valores no numéricos, valores no positivos y longitud disponible. Una serie corta puede describirse, pero no necesariamente reúne ventanas de validación suficientes, y ambas cosas se informan por separado. Las cifras de 18 y 24 observaciones se conservan como \textbf{referencias operativas documentadas}: por debajo de 18 se advierte y 24 se cita como longitud cómoda para modelación mensual básica. Son parámetros de la herramienta, no reglas estadísticas universales, y \textbf{ninguno impide proyectar}. El mínimo de 8 observaciones que hasta agosto de 2026 negaba la proyección fue retirado: era un literal sin fuente que cancelaba un pronóstico calculable. Lo único que impide modelar es un error crítico de los datos o una serie sin valores numéricos.

\begin{table}[H]
\centering
\caption{Controles previos al modelado de una serie ICOCIV}
\label{tab:controles_previos_modelado}
\scriptsize
\begin{adjustbox}{max width=\textwidth}
\begin{tabularx}{\textwidth}{@{}p{3.2cm}YYY@{}}
\toprule
\textbf{Control} & \textbf{Procedimiento} & \textbf{Riesgo evitado} & \textbf{Respuesta de la aplicación} \\
\midrule
Continuidad mensual & Comparar cada periodo con el mes inmediatamente anterior. & Interpretar un cambio acumulado como cambio mensual. & Advertir los meses faltantes y limitar el análisis. \\
Duplicados & Contar ocurrencias por periodo normalizado. & Asociar dos índices al mismo mes. & Marcar error crítico hasta resolver la ambigüedad. \\
Valores faltantes o no numéricos & Convertir el índice a formato numérico y registrar conversiones fallidas. & Ejecutar fórmulas con celdas vacías o texto. & Bloquear o advertir según el impacto. \\
Valores no positivos & Comprobar \(y_t>0\) cuando se evalúan transformaciones logarítmicas. & Aplicar \(\ln(y_t)\) fuera de su dominio. & Descartar esos candidatos, sin alterar la serie oficial. \\
Longitud & Comparar \(n\) con mínimos de modelado y ventanas walk-forward. & Seleccionar modelos con evidencia insuficiente. & Describir la serie o restringir la proyección. \\
\bottomrule
\end{tabularx}
\end{adjustbox}
\end{table}

\subsection{Estadística descriptiva, variaciones y valores atípicos}

La estadística descriptiva resume la serie antes de pronosticarla. La media, la desviación estándar muestral y el coeficiente de variación se calculan como:

\begin{align}
\bar{y} &= \frac{1}{n}\sum_{t=1}^{n}y_t, \\
s &= \sqrt{\frac{1}{n-1}\sum_{t=1}^{n}(y_t-\bar{y})^2}, \\
CV &= \frac{s}{\bar{y}}\times100.
\end{align}

La media describe el nivel promedio histórico, mientras \(s\) y \(CV\) ayudan a reconocer dispersión absoluta y relativa. En una serie con tendencia, estas medidas se leen junto con la gráfica, el valor inicial, el valor final y las variaciones; no sustituyen la validación predictiva.

La variación mensual, el factor entre el último valor observado y una proyección, y la variación acumulada se expresan como:

\begin{align}
\Delta_t &= \left(\frac{y_t}{y_{t-1}}-1\right)\times100, \\
F_h &= \frac{\hat{y}_{T+h}}{y_T}, \\
VA_h &= (F_h-1)\times100.
\end{align}

La estadística descriptiva sigue la práctica estándar de análisis exploratorio de datos \parencite{nist_eda_summary}. Para detectar cambios inusuales se usa el puntaje modificado con mediana y desviación absoluta mediana (MAD), definido en la Ecuación~\ref{eq:zmod}:

\begin{equation}
M_i=0,6745\frac{x_i-\widetilde{x}}{MAD},
\qquad MAD=\operatorname{mediana}(|x_i-\widetilde{x}|).
\label{eq:zmod}
\end{equation}

El factor $0{,}6745$ hace consistente el estadístico con la desviación estándar de una distribución normal. El criterio bibliográfico adoptado marca como posible atípico un caso con \(|M_i|>3,5\), umbral recomendado por el manual NIST/SEMATECH \parencite{nist_outliers}. La detección se aplica sobre las variaciones ---no sobre el nivel del índice, que por su tendencia haría aparecer como atípico todo extremo de la muestra--- y la aplicación no elimina automáticamente esos datos: conserva el valor oficial, registra el periodo afectado y lo comunica como advertencia técnica, porque un salto real de precios es información económica y no un error de medición.

El puntaje modificado puede calcularse sobre varias escalas de la misma serie ---variación porcentual mensual, diferencia simple y diferencia logarítmica---, que son transformaciones distintas de un mismo movimiento de precios. Cuando así se procede, un único hecho económico puede quedar señalado por más de uno de esos procedimientos, de modo que el recuento de detecciones deja de coincidir con el número de observaciones anómalas. La práctica de análisis exploratorio recomienda por ello consolidar las detecciones por periodo ---conservando qué escalas lo señalaron y el mayor $|M_i|$ observado--- antes de contabilizarlas, y clasificar cada periodo en categorías excluyentes \parencite{nist_outliers,nist_eda_summary}. Una tipología habitual distingue cuatro casos:

\begin{itemize}
  \item \textbf{Patrón calendario de cambio de año}: enero cuyo signo coincide con el del salto mediano y cuya serie cumple los criterios de recurrencia descritos en la sección siguiente. No corresponde contabilizarlo como valor atípico, porque expresa un comportamiento estructural del índice y no una anomalía.
  \item \textbf{Posible error de datos}: pico extremo, con $|M_i|$ superior al doble del umbral, que revierte al mes siguiente. Procede señalarlo para verificación contra la fuente oficial.
  \item \textbf{Posible cambio de nivel}: salto que persiste en los meses posteriores, compatible con un desplazamiento del nivel de la serie.
  \item \textbf{Posible atípico aislado}: el resto de los casos.
\end{itemize}

En ninguna de las cuatro categorías corresponde eliminar, interpolar ni suavizar la observación: la clasificación es descriptiva y su efecto se limita a cómo se comunica el hallazgo y a qué se contabiliza como anomalía en la evaluación de calidad de datos. La verificación concreta de este comportamiento sobre las series del \ICOCIV{} se documenta en el capítulo de pruebas.

\subsection{Modelos de pronóstico considerados}

\subsubsection{Modelo Naive}

El modelo Naive usa el último valor observado como pronóstico (Ecuación~\ref{eq:naive}) \parencite{hyndman2021fpp3}. Aunque es simple, cumple una función importante: sirve como referencia mínima. Si un modelo más complejo no mejora de forma razonable frente a esta base, su uso puede no estar justificado. En la aplicación, Naive se conserva como benchmark para comparar desempeño.

\begin{equation}
\hat{y}_{t+h}=y_t
\label{eq:naive}
\end{equation}

\subsubsection{Modelo Drift}

El modelo Drift proyecta el cambio promedio observado entre el primer y el último dato de la serie (Ecuación~\ref{eq:drift}) \parencite{hyndman2021fpp3}. A diferencia de Naive, incorpora una tendencia lineal simple. Su utilidad radica en ofrecer una referencia interpretable cuando existe una dirección clara en la serie, sin introducir demasiados parámetros. Por ejemplo, con $y_1=100$, $y_T=146$, $T=24$ y $h=3$, Naive proyecta $146$ y Drift proyecta $146+3\cdot 2=152$.

\begin{equation}
\hat{y}_{t+h}=y_t+h\left(\frac{y_t-y_1}{t-1}\right)
\label{eq:drift}
\end{equation}

\subsubsection{Regresiones sobre el nivel del índice}

La regresión lineal ordinaria representa una tendencia mediante:

\begin{equation}
\hat{y}_t=\beta_0+\beta_1t.
\end{equation}

En esta especificación \(\beta_0\) es el intercepto, es decir el valor esperado del índice cuando \(t=0\), y \(\beta_1\) es la pendiente, que expresa el cambio esperado del índice por cada unidad de tiempo transcurrida. Los coeficientes se estiman por mínimos cuadrados ordinarios, procedimiento que elige los valores que minimizan la suma de los cuadrados de los residuos \(\sum_t (y_t-\hat{y}_t)^2\) \parencite{montgomery2012introduccion}. La validez de la inferencia asociada a este estimador descansa en cuatro supuestos: linealidad en los parámetros, errores de media cero, varianza constante y ausencia de correlación entre ellos. El incumplimiento de alguno no invalida el ajuste como descripción, pero sí las medidas de precisión derivadas de él, razón por la cual el diagnóstico residual acompaña siempre a la estimación.

La regresión logarítmica temporal usa \(\ln(t+1)\) para representar trayectorias que crecen con rapidez al inicio y luego se desaceleran:

\begin{equation}
\hat{y}_t=\beta_0+\beta_1\ln(t+1).
\end{equation}

La alternativa exponencial o log-lineal ajusta \(\ln(y_t)=\beta_0+\beta_1t+\varepsilon_t\). Al regresar a la escala original, la aplicación puede aplicar el factor de corrección de retransformación de Duan:

\begin{equation}
\hat{y}_t=\exp(\hat{\beta}_0+\hat{\beta}_1t)
\left(\frac{1}{n}\sum_{i=1}^{n}\exp(\hat{\varepsilon}_i)\right).
\end{equation}

Esta corrección evita asumir que aplicar la función exponencial al valor ajustado reproduce sin sesgo la esperanza en escala original \parencite{duan1983}; la decisión de modelar en logaritmos o en niveles sigue las consideraciones de \parencite{manning_mullahy2001}. Las regresiones sobre nivel se tratan como referencias predictivas e interpretables, no como relaciones causales. En series con tendencia pueden producir un ajuste visual alto y aun así pronosticar mal; por eso se comparan fuera de muestra y no se seleccionan por \(R^2\) aislado \parencite{granger_newbold1974,springer_spurious_regressions}.

La regresión robusta de Huber reduce la influencia de observaciones extremas y actúa como candidato condicionado, no como mecanismo para borrar atípicos. En lugar de minimizar únicamente errores cuadrados, estima los parámetros mediante una pérdida que es cuadrática cerca de cero y lineal para residuos grandes:

\begin{equation}
\rho_{\delta}(e)=
\begin{cases}
\frac{1}{2}e^2, & |e|\leq\delta,\\
\delta\left(|e|-\frac{1}{2}\delta\right), & |e|>\delta.
\end{cases}
\end{equation}

El umbral \(\delta\) ---llamado \(\epsilon\) en buena parte de la literatura y del software--- controla a partir de qué magnitud un residuo deja de influir de forma cuadrática. La formulación procede del trabajo original de Huber \parencite{huber1964}, que introdujo esta pérdida como estimador de posición resistente a contaminación de la muestra; el tratamiento moderno del estimador y de su regularización se recoge en la literatura de regresión robusta \parencite{penn_state_robust_regression}. Un valor pequeño de \(\delta\) acerca el estimador a una regresión sobre la mediana y uno grande lo acerca a mínimos cuadrados ordinarios. La regresión robusta no depura la serie: reduce el peso de las observaciones extremas en la estimación de los coeficientes, mientras el dato observado permanece intacto.

\subsubsection{Modelos sobre variación y log-variación}

Una serie de índices admite modelarse sobre tres magnitudes distintas: el nivel \(y_t\), la variación relativa \(\Delta_t=(y_t-y_{t-1})/y_{t-1}\) y la log-variación o retorno logarítmico, definido como la diferencia de logaritmos de dos observaciones consecutivas:

\begin{equation}
r_t=\ln(y_t)-\ln(y_{t-1})=\ln\!\left(\frac{y_t}{y_{t-1}}\right).
\label{eq:logvariacion}
\end{equation}

La log-variación es una transformación de uso extendido en el análisis de series económicas y financieras \parencite{tsay2010analysis}. Su interpretación se apoya en la aproximación \(\ln(1+x)\approx x\) para valores pequeños de \(x\): dado que \(r_t=\ln(1+\Delta_t)\), cuando la variación relativa es pequeña el retorno logarítmico se aproxima numéricamente a la tasa de crecimiento del periodo. Para una variación mensual del 1\,\%, por ejemplo, la log-variación vale 0,00995, con una diferencia inferior a cinco milésimas de punto porcentual; la discrepancia crece a medida que la variación se aleja de cero.

Dos propiedades explican su conveniencia. La primera es la aditividad temporal: la suma de log-variaciones consecutivas equivale al cambio logarítmico acumulado del intervalo completo, mientras que las variaciones relativas deben componerse de forma multiplicativa. La segunda es que convierte un efecto proporcional al nivel de la serie en uno aditivo, lo que hace comparables movimientos ocurridos en niveles distintos del índice. Resulta apropiada cuando la serie es estrictamente positiva y su dinámica se comporta de forma proporcional al nivel, condición habitual en índices de precios. No es adecuada cuando la serie admite valores nulos o negativos, porque el logaritmo deja de estar definido.

Modelar la variación o la log-variación implica pronosticar un cambio y reconstruir después el nivel, lo que traslada la acumulación del error al proceso de reconstrucción. Conviene precisar que una regresión logarítmica del tiempo sobre el nivel no equivale a un modelo de log-variación mensual: las transformaciones y sus supuestos son diferentes.

\subsubsection{Holt lineal}

El modelo Holt lineal pertenece a la familia de suavizamiento exponencial e incorpora componentes de nivel y tendencia (Ecuación~\ref{eq:holt}) \parencite{holt1957}. En una serie de índices puede capturar cambios graduales sin exigir una estructura compleja. La aplicación lo considera cuando la longitud y continuidad de la serie permiten estimar sus parámetros.

\begin{equation}
\hat{y}_{t+h}=l_t+hb_t
\label{eq:holt}
\end{equation}

\subsubsection{Holt amortiguado}

El modelo Holt amortiguado introduce un parámetro $\phi\in(0,1)$ que reduce la tendencia progresivamente (Ecuación~\ref{eq:holt_amortiguado}) \parencite{gardner_mckenzie1985}. Esta característica es útil cuando extender una tendencia lineal completa podría producir valores poco defendibles en horizontes largos. En el sistema, esta alternativa se evalúa como opción de tendencia moderada.

\begin{equation}
\hat{y}_{t+h}=l_t+(\phi+\phi^2+\cdots+\phi^h)b_t
\label{eq:holt_amortiguado}
\end{equation}

\subsubsection{Suavizamiento exponencial ETS}
\label{sec:teoria_ets}

La familia ETS reformula el suavizamiento exponencial como un modelo de espacio de estados con tres componentes explícitos: el error \(E\), la tendencia \(T\) y la estacionalidad \(S\) \parencite{hyndman2021fpp3}. Cada componente admite variantes, lo que da lugar a la notación ETS(\(\cdot,\cdot,\cdot\)): el error puede ser aditivo (A) o multiplicativo (M); la tendencia, inexistente (N), aditiva (A) o aditiva amortiguada (A\textsubscript{d}); y la estacionalidad, inexistente (N), aditiva (A) o multiplicativa (M). La forma aditiva supone que las fluctuaciones conservan una magnitud comparable a lo largo de la serie, mientras que la multiplicativa supone que crecen proporcionalmente al nivel, lo que la hace apropiada para series cuya amplitud aumenta con la escala.

La relación con el suavizamiento exponencial clásico es directa: cada método de la familia de suavizamiento tiene su contraparte en espacio de estados. El suavizamiento exponencial simple corresponde a ETS(A,N,N); el modelo de Holt con tendencia lineal, a ETS(A,A,N); y el modelo de Holt con tendencia amortiguada, a ETS(A,A\textsubscript{d},N). Las ecuaciones de actualización del nivel y de la tendencia son las mismas descritas para Holt en los apartados anteriores; lo que ETS añade es una formulación probabilística del término de error que permite estimar los parámetros por máxima verosimilitud, seleccionar la variante mediante criterios de información y derivar intervalos de predicción analíticos coherentes con el modelo.

Esa formulación explica sus ventajas y también sus condiciones de uso. Entre las ventajas figuran la selección automática de la estructura de componentes, la obtención de intervalos de predicción con respaldo distribucional y la capacidad de representar estacionalidad, ausente en los métodos de tendencia pura. Entre las limitaciones, la estimación conjunta de los parámetros de suavizamiento y de los estados iniciales exige más información que un método con coeficientes fijos: identificar una componente estacional mensual requiere varios ciclos completos, y las variantes multiplicativas no admiten valores no positivos. La superficie de verosimilitud puede además presentar óptimos locales, de modo que el resultado depende de la implementación y de la inicialización empleadas.

En cuanto a la interpretabilidad, ETS conserva la lectura por componentes ---nivel, tendencia y estacionalidad son magnitudes con significado--- pero introduce parámetros estimados cuya variación entre series dificulta explicar por qué dos series semejantes reciben especificaciones distintas. La literatura recomienda por ello contrastar su desempeño fuera de muestra frente a alternativas más simples antes de adoptarlo, en lugar de suponer que una familia más general pronostica mejor \parencite{hyndman2021fpp3,hyndman_koehler2006}.

\subsubsection{Combinación de pronósticos}
\label{sec:teoria_combinacion}

La combinación de pronósticos consiste en obtener una predicción única a partir de varios modelos en lugar de seleccionar uno solo. Es una práctica extendida en el análisis de series temporales \parencite{hyndman2021fpp3}. Su forma más simple es el promedio aritmético de los \(k\) pronósticos disponibles, que no requiere estimar ningún parámetro adicional. Una alternativa habitual pondera cada modelo de forma inversa a su error fuera de muestra, de manera que los modelos más precisos reciban mayor peso:

\begin{equation}
w_i=\frac{1/RMSE_i}{\sum_{j=1}^{k}1/RMSE_j}
\end{equation}

El argumento teórico a favor de combinar es la reducción de varianza. Si los errores de los modelos no están perfectamente correlacionados, la varianza del error del promedio es menor que la varianza media de los errores individuales, y la combinación resulta además más robusta frente a la incertidumbre de modelo: no obliga a acertar cuál es el modelo correcto, sino a distribuir el riesgo entre varias especificaciones razonables.

Ese beneficio, sin embargo, es condicional. Cuando los modelos comparten estructura y datos, sus errores están fuertemente correlacionados y la reducción de varianza se desvanece. Los pesos estimados sobre el mismo conjunto de errores con el que después se evalúa la combinación producen una mejora aparente que no se sostiene al medirla sin fuga de información, por lo que la evaluación exige separar la muestra de calibración de la muestra de medición. La combinación reduce asimismo la interpretabilidad del resultado: deja de existir una ecuación única y unos parámetros que expliquen la trayectoria proyectada, lo que es relevante cuando el pronóstico debe justificarse ante un tercero. Por último, si los pesos se recalculan para cada horizonte, los primeros valores de la trayectoria pueden variar según el horizonte solicitado, lo que compromete la consistencia de la proyección. La literatura recomienda, en consecuencia, sostener la decisión de combinar en validación fuera de muestra y no en su plausibilidad teórica \parencite{hyndman2021fpp3,tashman2000}.

\subsection{Patrón de cambio de año en las series ICOCIV}

Los modelos de tendencia descritos reparten el crecimiento de forma aproximadamente uniforme entre los meses. Varias series \ICOCIV{}, en cambio, concentran una parte importante de su variación anual en la transición de diciembre a enero, asociada a la actualización anual de precios que recoge el índice. Cuando ese patrón existe y no se representa, el modelo subestima enero y sobrestima el resto del año, aun cuando el crecimiento anual total esté bien capturado. El tratamiento de componentes de calendario y estacionalidad en series mensuales es un problema clásico del análisis de pronóstico \parencite{hyndman2021fpp3}.

La detección se realiza sobre variaciones logarítmicas \(r_t=\ln(y_t/y_{t-1})\). El uso de la transformación logarítmica se justifica porque el salto de cambio de año es aproximadamente proporcional al nivel del índice, y la diferencia de logaritmos convierte un efecto multiplicativo en uno aditivo, comparable entre meses \parencite{hyndman2021fpp3}. Se separan las transiciones de diciembre a enero del resto de meses y se resumen con estadística robusta: \(\gamma\) es la mediana de las transiciones de cambio de año y \(m\) la mediana del valor absoluto de las variaciones restantes. La mediana y la desviación absoluta respecto de la mediana son estimadores resistentes a valores extremos, recomendados cuando pocas observaciones podrían distorsionar la media \parencite{nist_outliers}. Se considera que existe un patrón de calendario cuando se cumplen simultáneamente las tres condiciones de la Ecuación~\ref{eq:calendario}: al menos dos transiciones de cambio de año, razón \(|\gamma|/m\) superior a 1,5 y consistencia de signo de al menos 0,6. Estas tres condiciones constituyen reglas de decisión definidas para la implementación y calibradas con el anexo vigente; no provienen de una fuente externa y no se presentan como un umbral universal. \textbf{Desde el 12 de agosto de 2026 estas condiciones no activan ningún ajuste}: se conservan únicamente como criterio de \emph{descripción} del patrón, según se explica más adelante en este mismo apartado.

\begin{equation}
K\geq2,\qquad \frac{|\gamma|}{m}>1{,}5,\qquad \text{consistencia de signo}\geq0{,}6
\label{eq:calendario}
\end{equation}

donde \(K\) es el número de transiciones observadas. El uso de medianas evita que un único año extremo determine la magnitud estimada, y la exigencia de recurrencia y consistencia de signo distingue un patrón de calendario de un valor atípico aislado, que no se repite en la misma posición del calendario. Los valores de corte son parámetros de decisión internos calibrados con el anexo vigente: las series con patrón presentan razones entre 4 y 9, mientras la serie de control (acero) queda en 0,36.

La \emph{consistencia de signo} se define como la proporción de las \(K\) transiciones de diciembre a enero cuyo signo coincide con el signo de la mediana \(\gamma\). Formalmente, si \(s_k=\operatorname{sgn}(r_k)\) es el signo de la \(k\)-ésima transición, la consistencia es \(\tfrac{1}{K}\sum_{k=1}^{K}\mathbf{1}\{s_k=\operatorname{sgn}(\gamma)\}\), un valor entre 0 y 1. Su interpretación estadística se apoya en la lógica de la prueba de signos: si no existiera un efecto sistemático de calendario, un salto tendría aproximadamente la misma probabilidad de ser positivo o negativo, de modo que una proporción cercana a 0,5 indica ausencia de dirección y una proporción cercana a 1 indica un efecto recurrente en un mismo sentido \parencite{nist_eda_summary}. Por ejemplo, si en cinco años cuatro transiciones de cambio de año son positivas y una es negativa, la consistencia es \(4/5=0{,}8\); si fueran dos positivas y tres negativas, sería \(2/5=0{,}4\). Dentro del criterio de \emph{detección}, esta medida cumple la función de exigir que el salto se repita en la misma dirección: un valor de referencia de al menos 0{,}6 evita confundir un movimiento aislado con un patrón anual estable. Conviene subrayar que ese valor, como los demás de la Ecuación~\ref{eq:calendario}, fue calibrado sobre el propio anexo de aplicación y \textbf{no habilita ningún tratamiento}: la versión final describe el patrón cuando se cumple el criterio, pero no corrige por él los valores proyectados.

\paragraph{Tratamiento retirado el 12 de agosto de 2026.} La aplicación
incorporó hasta esa fecha un ajuste que reconcentraba el salto en enero mediante
el factor de la Ecuación~\ref{eq:factor_calendario}. \textbf{Ese tratamiento se
retiró del cálculo productivo}, y la ecuación se conserva aquí únicamente para
documentar qué se hacía y por qué dejó de hacerse. Los motivos son cuatro, y
ninguno se refiere al fenómeno, que está medido y es contundente:

\begin{enumerate}
  \item usar la mediana de los saltos pasados como estimador del salto
  \emph{futuro} no cuenta con fuente: la robustez de la mediana está publicada,
  su uso predictivo en esta forma no;
  \item el término \(h/12\) supone que el modelo base repartió \(\gamma\) de
  manera uniforme entre los doce meses. Esa hipótesis \textbf{es falsa para el
  modelo ingenuo}, cuyo pronóstico es plano y no reparte crecimiento alguno, y
  es distinta para cada uno de los demás candidatos; el factor, sin embargo, se
  aplicaba a todos por igual sin considerar cuál había sido seleccionado;
  \item los tres cortes de detección ---al menos dos transiciones, razón
  superior a 1,5 y consistencia de signo de al menos 0,6--- fueron calibrados
  sobre el propio anexo de aplicación, lo que impide presentarlos como criterio
  independiente;
  \item la documentación oficial del \ICOCIV{} consultada no describe ninguna
  discontinuidad de diciembre a enero, de modo que el ajuste tampoco podía
  presentarse como una corrección oficial del índice.
\end{enumerate}

\textbf{El patrón se sigue midiendo y publicando como diagnóstico descriptivo}
---número de transiciones, salto mediano, razón frente al movimiento típico y
consistencia de signo---, pero \textbf{no modifica ningún valor proyectado}. No
se introdujo ningún método sustituto: hacerlo habría añadido un candidato al
catálogo y reabierto la definición del conjunto de modelos comparados.

El factor retirado era:

\begin{equation}
f_h=\exp\!\left(\gamma\left(n_h-\frac{h}{12}\right)\right),\qquad \hat{y}^{\,\text{aj}}_{T+h}=\hat{y}_{T+h}\, f_h,
\label{eq:factor_calendario}
\end{equation}

donde \(n_h\) era el número de eneros contenidos en los \(h\) pasos proyectados.
El factor se evaluaba paso a paso y en \(h=12\) valía exactamente 1, de modo que
redistribuía el salto dentro del año sin alterar el crecimiento anual. Esa
propiedad era correcta; lo que no lo era es el resto de la construcción, según
los cuatro motivos enumerados arriba.

La alternativa metodológicamente limpia ---incorporar una variable indicadora de
enero \emph{dentro} de la regresión, de modo que su coeficiente se estime
conjuntamente con la tendencia--- queda documentada como línea de trabajo
futuro: elimina de raíz los tres cortes de detección y el término \(h/12\),
porque no habría doble contabilidad que descontar. No se adoptó en este trabajo
porque incorpora un modelo nuevo al catálogo comparado y exige cerrar antes la
definición del primer origen de la validación temporal, que sigue abierta.

\subsection{Validación predictiva y métricas de error}

\subsubsection{Backtesting y walk-forward}

El backtesting temporal consiste en evaluar modelos usando cortes históricos de entrenamiento y prueba. En lugar de ajustar un modelo con toda la serie y reportar únicamente su ajuste interno, se simula cómo habría pronosticado el modelo en periodos anteriores.

Conviene fijar antes tres nociones que estructuran este procedimiento. El \emph{origen de pronóstico} es el instante desde el cual se proyecta: separa lo que el modelo puede conocer de lo que aún no ha ocurrido. La \emph{ventana de entrenamiento} es el conjunto de observaciones empleado para ajustar el modelo en ese origen, y contiene únicamente datos anteriores o iguales a él. El \emph{horizonte de pronóstico} es el número de periodos que se proyectan hacia adelante desde el origen, de modo que el paso $h$ corresponde a la observación situada $h$ periodos después.

La validación walk-forward repite el ajuste en sucesivos orígenes que avanzan en el tiempo. En una \emph{ventana expansiva} el inicio permanece fijo en la primera observación y el final se desplaza con cada origen, de manera que cada iteración incorpora la historia nueva sin descartar la anterior; en una ventana deslizante, en cambio, la longitud se mantiene constante y el inicio avanza junto con el final. En ambos casos el pronóstico generado en un origen se compara con observaciones posteriores que no intervinieron en el ajuste, lo que preserva el orden temporal de los datos y evita la fuga de información desde el futuro hacia la estimación. Es esa separación, y no el tamaño de la muestra, la que distingue una medida de capacidad predictiva de una medida de ajuste histórico \parencite{tashman2000,hyndman_tscv_2021}.

Repetir la evaluación en varios orígenes permite además valorar la estabilidad del desempeño: un modelo puede acertar en promedio y comportarse de forma errática entre orígenes, lo que un único corte no revelaría.

La ausencia de fuga de información se garantiza con dos reglas: el modelo se reestima en cada origen usando solo el pasado y la escala del MASE se calcula dentro de cada ventana de entrenamiento. Hasta el 12 de agosto de 2026 existía una tercera, referida al parámetro \(\gamma\) del ajuste de cambio de año, que se reestimaba por origen; al retirarse el tratamiento del cálculo productivo esa regla dejó de tener objeto. En la aplicación, esta validación ayuda a decidir si la proyección es defendible o si debe acompañarse de una advertencia sobre la evidencia disponible. El criterio no depende solamente de producir un número; también considera advertencias, continuidad, errores extremos e incertidumbre.

\subsubsection{MAE, RMSE, MAPE, sMAPE y MASE}

Las métricas de error se calculan sobre los $n$ errores de predicción fuera de muestra $e_i=y_i-\hat{y}_i$, siguiendo las definiciones estándar del análisis de exactitud de pronósticos \parencite{hyndman_koehler2006,hyndman2021fpp3}. El error absoluto medio y la raíz del error cuadrático medio se definen en las Ecuaciones~\ref{eq:mae} y~\ref{eq:rmse}:

\begin{equation}
MAE=\frac{1}{n}\sum_{i=1}^{n}|e_i|
\label{eq:mae}
\end{equation}

\begin{equation}
RMSE=\sqrt{\frac{1}{n}\sum_{i=1}^{n}e_i^{2}}
\label{eq:rmse}
\end{equation}

El MAE resume el error absoluto promedio, mientras el RMSE penaliza con mayor fuerza los errores grandes. Los errores relativos porcentuales se expresan mediante el MAPE (Ecuación~\ref{eq:mape}) y el sMAPE (Ecuación~\ref{eq:smape}); el primero excluye las observaciones nulas para evitar divisiones por cero y el segundo usa la forma acotada en $[0,200]$:

\begin{equation}
MAPE=\frac{100}{n}\sum_{i=1}^{n}\left|\frac{e_i}{y_i}\right|,\qquad y_i\neq0
\label{eq:mape}
\end{equation}

\begin{equation}
sMAPE=\frac{100}{n}\sum_{i=1}^{n}\frac{2\,|e_i|}{|y_i|+|\hat{y}_i|}
\label{eq:smape}
\end{equation}

El sesgo medio (Ecuación~\ref{eq:sesgo}) mide si el modelo subestima o sobrestima de forma sistemática; un valor alejado de cero indica error direccional y no aleatorio:

\begin{equation}
Sesgo=\frac{1}{n}\sum_{i=1}^{n}e_i
\label{eq:sesgo}
\end{equation}

El MASE escala el error frente al método naive dentro de la muestra de entrenamiento, lo que lo hace comparable entre series de escalas distintas \parencite{hyndman_koehler2006}. Su denominador se calcula \emph{dentro de cada ventana de entrenamiento} del walk-forward, nunca con datos de prueba, para no incorporar información futura (Ecuación~\ref{eq:mase}):

\begin{equation}
MASE=\frac{\frac{1}{n}\sum_{i=1}^{n}|e_i|}{\frac{1}{T-1}\sum_{t=2}^{T}|y_t-y_{t-1}|}
\label{eq:mase}
\end{equation}

Un MASE mayor que 1 indica que, en promedio, el error del modelo sobre esa ventana de entrenamiento es superior al de la referencia naive dentro de muestra; un valor menor que 1 indica mejora frente a ella. Es una lectura comparativa, no un veto: el valor se publica como información y no cambia por sí solo el estado del horizonte. A modo de ejemplo, con errores absolutos de prueba $\{1;1;1;2\}$ (MAE $=1{,}25$) y un tramo de entrenamiento $\{90;92;96\}$ cuya escala naive es $\tfrac{1}{2}(|2|+|4|)=3$, resulta $MASE=1{,}25/3=0{,}4167$. Estas seis métricas son complementarias entre sí y se reportan todas; el modelo entregado lo decide únicamente el RMSE fuera de muestra sobre la muestra común a todos los candidatos (§\ref{sec:criterios_metricas}), y ninguna de ellas acota por sí sola el horizonte recomendado. Las fórmulas están verificadas contra valores calculados a mano mediante pruebas automatizadas, cuyos resultados se presentan en el capítulo de pruebas.

\subsubsection{Diagnóstico residual e intervalos de predicción}

Los residuos \(e_t=y_t-\hat{y}_t\) admiten cuatro contrastes formales: la nulidad de su media, la ausencia de autocorrelación conjunta, la constancia de su varianza y su normalidad. El estadístico Durbin--Watson se calcula como se indica en la Ecuación~\ref{eq:dw}:

\begin{equation}
DW=\frac{\sum_{t=2}^{T}(e_t-e_{t-1})^2}{\sum_{t=1}^{T}e_t^2}.
\label{eq:dw}
\end{equation}

Un valor cercano a 2 sugiere ausencia de autocorrelación de primer orden fuerte, se acerca a 0 con autocorrelación positiva y a 4 con negativa. Su contraste formal requiere las tablas de significancia $d_L/d_U$, que dependen del tamaño muestral y del número de regresores, por lo que en modelos no lineales o de suavizamiento se interpreta únicamente como lectura descriptiva.

El estadístico de Durbin--Watson examina un solo rezago. Para contrastar de forma conjunta la presencia de autocorrelación en varios rezagos simultáneamente se emplea la prueba de Ljung--Box \parencite{ljung_box1978}. Su hipótesis nula sostiene que las primeras $h$ autocorrelaciones de la serie de residuos son conjuntamente nulas, es decir que los residuos se comportan como ruido blanco; la hipótesis alternativa afirma que al menos una de ellas es distinta de cero. El estadístico se calcula según la Ecuación~\ref{eq:ljungbox}:

\begin{equation}
Q=n(n+2)\sum_{k=1}^{h}\frac{\hat{\rho}_k^{2}}{n-k},
\label{eq:ljungbox}
\end{equation}

donde $n$ es el número de residuos disponibles, $h$ el número de rezagos considerados y $\hat{\rho}_k$ la autocorrelación muestral de los residuos en el rezago $k$. El factor $n(n+2)/(n-k)$ corrige el sesgo del estadístico en muestras finitas, mejora que distingue esta prueba de la formulación previa de Box y Pierce.

Bajo la hipótesis nula, $Q$ sigue asintóticamente una distribución $\chi^2$. Cuando la prueba se aplica a residuos de un modelo autorregresivo y de media móvil, los grados de libertad son $h-p-q$, donde $p$ y $q$ son los órdenes estimados de ese modelo; cuando los residuos no provienen de una estructura de ese tipo, no hay parámetros de autocorrelación que descontar y los grados de libertad son $h$. Un valor-p pequeño conduce a rechazar la nulidad y a concluir que los residuos conservan estructura temporal no capturada por el modelo; un valor-p alto indica que no hay evidencia suficiente para afirmar esa estructura, lo que no equivale a demostrar su ausencia. La elección de $h$ condiciona el resultado: un valor demasiado grande reparte la evidencia entre muchos rezagos con pocas observaciones cada uno y reduce la potencia de la prueba \parencite{hyndman2021fpp3}.

La nulidad de la media residual admite un contraste directo. Bajo la hipótesis nula \(H_0:\mu_e=0\), el estadístico de la prueba \(t\) de una muestra se construye según la Ecuación~\ref{eq:tmedia}:

\begin{equation}
t=\frac{\bar{e}-0}{s_e/\sqrt{n}},
\label{eq:tmedia}
\end{equation}

donde \(\bar{e}\) es la media muestral de los residuos, \(s_e\) su desviación típica muestral y \(n\) el número de residuos. Bajo la nulidad el estadístico sigue una distribución \(t\) de Student con \(n-1\) grados de libertad. Su validez descansa en la aproximación normal de la media muestral, que con muestras pequeñas es más débil; la prueba no es calculable con menos de dos residuos ni con dispersión nula. Un valor-p pequeño indica que el modelo yerra de forma sistemática en una dirección; uno alto no demuestra la ausencia de sesgo, solo que no hay evidencia suficiente para afirmarlo al nivel elegido.

La constancia de la varianza residual se contrasta con la prueba de Breusch y Pagan \parencite{breusch_pagan1979}. Su hipótesis nula sostiene que la varianza de los residuos no depende de los regresores considerados. El procedimiento regresa los residuos al cuadrado sobre esa matriz de regresores y contrasta la nulidad conjunta de sus pendientes; el estadígrafo del multiplicador de Lagrange sigue asintóticamente una distribución \(\chi^2\) con tantos grados de libertad como regresores distintos de la constante. Existe una alternativa más general, la prueba de White \parencite{white1980}, que añade los cuadrados y los productos cruzados de los regresores y no exige especificar la forma de la heterocedasticidad; a cambio consume más grados de libertad, de modo que con un único regresor y series de pocas decenas de observaciones pierde potencia sin aportar información adicional. Como resultado asintótico, la aproximación \(\chi^2\) se degrada con muestras cortas, y esa limitación debe declararse junto al resultado.

Los tres contrastes anteriores y el de normalidad que sigue comparten una propiedad que conviene enunciar: no rechazar la hipótesis nula no equivale a demostrarla. Un valor-p alto solo indica que la evidencia disponible no basta para descartarla al nivel de significancia elegido, elección convencional y no un valor universal.

El contraste de normalidad de Jarque--Bera evalúa si la distribución de los residuos es compatible con la normalidad, a partir de la asimetría $S$ y la curtosis $K$ muestrales, según la Ecuación~\ref{eq:jb}:

\begin{equation}
JB=\frac{n}{6}\left(S^2+\frac{(K-3)^2}{4}\right).
\label{eq:jb}
\end{equation}

Bajo la hipótesis nula de normalidad, $JB$ sigue asintóticamente una distribución $\chi^2$ con dos grados de libertad \parencite{jarque_bera1987}. El valor-p corresponde a la función de supervivencia de esa distribución, que para dos grados de libertad admite la forma cerrada $p=\exp(-JB/2)$. Al tratarse de un resultado asintótico, su validez depende del tamaño de la muestra de residuos: con muy pocas observaciones la aproximación pierde exactitud, y con dispersión nula el estadístico no está definido. El nivel de referencia habitual es $\alpha=0{,}05$.

\paragraph{Advertencia sobre el alcance de este apartado.} \textbf{La versión productiva de la aplicación no publica intervalos de predicción.} Ninguno de los métodos evaluados alcanzó una construcción completa sustentada bajo los requisitos metodológicos adoptados ---formulación, supuestos, nivel nominal, grados de libertad, dispersión, tratamiento por horizonte, tamaño mínimo, muestras pequeñas y evaluación de cobertura, todo ello con fuente verificable---, de modo que entregar sus límites equivaldría a afirmar una precisión que la aplicación no puede defender. El intervalo se retiró de la interfaz, del CSV, del informe HTML y de los documentos DOCX y PDF, junto con su cobertura observada, su clasificación y su ancho relativo. Retirarlo \emph{no} significa que la incertidumbre no exista ni que el pronóstico sea exacto: significa que no se puede acotar con un método sustentado; el pronóstico puntual sí se publica cuando es calculable. Lo que sigue documenta el método construido y evaluado, y las razones por las que no se adoptó. Se conserva porque una omisión sin explicación no es auditable, y porque de esa evaluación depende la decisión de retirarlo.

Un pronóstico puntual no expresa la incertidumbre asociada a la predicción. Un intervalo de predicción acota el rango en que puede situarse una observación futura, y su construcción a partir de errores de validación exige que cada paso emplee los errores observados a su propio horizonte, sin reescalar los de otro: la dispersión del error crece con la distancia predictiva y mezclar horizontes distorsiona la banda. Sobre esa muestra de errores pueden combinarse dos construcciones, ambas centradas en la media $\bar{e}_h$ de los errores del horizonte, de modo que un sesgo sistemático del modelo desplace la banda hacia donde efectivamente caen las observaciones:

\begin{equation}
IP_{1-\alpha,h}=\hat{y}_{T+h}+\bar{e}_h\pm\max\left(Q_{1-\alpha}^{\text{conf}},\;t_{1-\alpha/2,\,n_h-1}\,s_h\sqrt{1+\tfrac{1}{n_h}}\right).
\label{eq:ip_conformal}
\end{equation}

El primer término es un cuantil de orden con corrección de muestra finita: ordenados los $n_h$ valores $|e_i-\bar{e}_h|$, se toma el $k$-ésimo con $k=\lceil (n_h+1)(1-\alpha)\rceil$. Ese índice es el que emplea la predicción conformal dividida para su cuantil de calibración \parencite{lei2018conformal,angelopoulos2023conformal}, y la corrección $(n_h+1)$ evita el sesgo hacia adentro que presenta el percentil empírico simple en muestras pequeñas.

Debe precisarse el alcance de esa referencia. La predicción conformal dividida garantiza cobertura de al menos $1-\alpha$ en muestra finita bajo \emph{intercambiabilidad} entre los errores de calibración y el error futuro. Los errores de validación temporal de una serie \ICOCIV{} no satisfacen ese supuesto: los orígenes sucesivos comparten casi toda su historia de entrenamiento, el modelo se reestima en cada corte y la dispersión del error varía a lo largo del tiempo. A ello se añade que aquí el estadístico se centra con la media de la propia muestra de calibración, lo que rompe la intercambiabilidad estricta que exige el resultado teórico. Por consiguiente, la garantía de cobertura en muestra finita del método conformal no es invocable en este contexto: el cuantil de orden conserva su valor como criterio conservador, inspirado en esa literatura, pero su validez debe sostenerse en la verificación empírica de la cobertura alcanzada, no en el teorema.

El umbral $n_h\ge 1/\alpha-1$ ---19 errores para el nivel del 95\,\%--- conserva sentido con independencia de la garantía: por debajo de él ninguna observación de la muestra alcanza la cola correspondiente y el cuantil de orden degenera en el máximo observado. En ese caso opera el segundo término, el intervalo de predicción $t$ de Student clásico \parencite{hyndman2021fpp3}, y la banda pasa a descansar en un supuesto paramétrico de normalidad que debe declararse de forma explícita. La amplitud se fuerza no decreciente con el paso, porque la incertidumbre declarada a $p+1$ pasos no puede ser menor que la declarada a $p$ pasos, y el límite inferior se recorta en cero por tratarse de un índice positivo.

Debe subrayarse que se trata de intervalos de \emph{predicción}, que cubren una observación futura del índice, y no de intervalos de \emph{confianza}, que cubrirían un parámetro poblacional. La distinción es relevante en el uso presupuestal: la banda comunica el rango donde puede situarse el índice publicado, no la precisión con la que se estimó un coeficiente del modelo.

\subsubsection{Selección de modelo, clasificación por horizonte y reproducibilidad}

El desempeño relativo de los modelos puede variar con \(h\), pero la aplicación no elige un modelo distinto para cada horizonte: eso haría que los primeros meses proyectados cambiaran según el horizonte que el usuario solicitara, para la misma serie. La selección agrega la evidencia de todos los horizontes evaluados en un único puntaje ---el error cuadrático medio de validación temporal calculado sobre la muestra común a todos los candidatos--- y de él resulta un modelo único por serie. Las demás métricas ---MAE, MAPE, sMAPE, MASE, sesgo, errores extremos, estabilidad y comparación contra Naive y Drift--- no puntúan la selección y tampoco vetan: se publican como información descriptiva junto al resultado. El ancho relativo del intervalo se calculó como diagnóstico interno durante la evaluación del método (§\ref{sec:criterios_metricas}), pero no forma parte de las salidas de la versión final. Los criterios de información dentro de muestra, como el AICc, se reportan pero no sustituyen ni complementan el backtesting en la decisión.

Se distinguen horizonte solicitado, máximo recomendado, máximo permitido como escenario y horizonte no recomendable. La interfaz ofrece accesos rápidos de 1, 3, 6, 12 y 18 meses, pero 18 no constituye un máximo estadístico fijo. Para reproducir una salida deben conservarse archivo y ruta fuente, periodos, observaciones, modelo, ecuación, parámetros, definición de \(t\), horizonte, proyecciones, métricas fuera de muestra con su número de ventanas, factor y variación acumulada. Los intervalos no forman parte de la salida y por tanto tampoco de lo reproducible.

\subsection{Antecedentes y trabajos relacionados}

Los antecedentes del proyecto se organizan en tres planos que convergen en la necesidad atendida: el antecedente institucional y operativo, el antecedente normativo de la actualización de precios, y el antecedente técnico del análisis de series de costos.

\textbf{Antecedente institucional y operativo:} la Secretaría de Infraestructura de la Gobernación del Cauca adelanta de manera recurrente procesos de actualización, ajuste y revisión de precios en contratos de obra civil, apoyándose en información oficial del \DANE. Para ello ha empleado un archivo institucional en Microsoft Excel que filtra los índices y ejecuta los cálculos de ajuste. Este archivo constituye el antecedente directo del presente trabajo: demostró la utilidad del enfoque, pero su dependencia de manipulación manual, su complejidad creciente y la necesidad de conocimiento específico para operarlo limitaban su apropiación por el equipo técnico y aumentaban el riesgo de errores. Casos concretos de revisión de precios en contratos del departamento ---como las solicitudes de ajuste asociadas a fluctuaciones de insumos como el acero--- evidenciaron la necesidad de un procedimiento más trazable, estandarizado y verificable. El presente sistema recoge ese flujo operativo y lo traslada a una aplicación que conserva la lógica técnica pero mejora la trazabilidad, la reproducibilidad y el control de errores.

\textbf{Antecedente normativo de la actualización de precios:} la actualización de precios en contratos de obra pública se enmarca en la figura del restablecimiento del equilibrio económico del contrato, reconocida por la doctrina de contratación pública colombiana \parencite{cce2024_equilibrio}. En el plano técnico, los índices oficiales del \DANE{} proporcionan la base metodológica para medir la variación de costos \parencite{dane_ficha_icociv,dane_metodologia_icociv}, mientras que los lineamientos institucionales del \IDU{} y el procedimiento de implementación del \ICOCIV{} del \INVIAS{} establecen cómo aplicar el ajuste y cómo tratar la transición entre el \ICCP{} y el \ICOCIV{} \parencite{idu_circular_6_2022,invias2022_icociv}. El Manual de Interventoría de Obra Pública del \INVIAS{} aporta el criterio de referencia sobre el momento contractual frente al cual se define la fecha base del ajuste \parencite{invias2010_interventoria}. Estos documentos constituyen el antecedente normativo que el módulo de empalme del sistema implementa de forma trazable.

\textbf{Antecedente técnico del análisis de series de costos:} la componente de proyección se apoya en el cuerpo de conocimiento del pronóstico de series de tiempo. La comparación de modelos mediante validación fuera de muestra y evaluación de origen rodante, en lugar del solo ajuste histórico, es una práctica consolidada en la literatura \parencite{hyndman2021fpp3,tashman2000}, así como el uso de métricas comparables entre series y de referencias simples como base de contraste \parencite{hyndman_koehler2006}. El tratamiento robusto de valores atípicos mediante estadística basada en la mediana \parencite{nist_outliers}, la corrección del sesgo de retransformación en modelos logarítmicos \parencite{duan1983} y la regresión robusta frente a observaciones extremas \parencite{huber1964} conforman los antecedentes metodológicos que el sistema adapta al caso de las series mensuales del \ICOCIV. La advertencia sobre regresiones espurias en series con tendencia compartida \parencite{granger_newbold1974} justifica que la selección de modelos no se base en el ajuste dentro de muestra. Con estos elementos, el sistema se plantea como una integración aplicada de consulta de datos oficiales, análisis estadístico de series, empalme normativo de índices y generación de evidencias verificables, orientada a apoyar la planeación y la estructuración de presupuestos de referencia en la entidad.

\section{METODOLOGÍA}

La metodología se estructuró en fases orientadas a transformar una necesidad técnica de análisis de variaciones de precios en una herramienta computacional verificable. Para ello, se revisaron documentos institucionales y técnicos, se identificó el flujo de trabajo utilizado por la entidad, se diseñó una arquitectura modular, se implementaron procedimientos de análisis estadístico y empalme de índices, y se verificó el funcionamiento de la aplicación mediante pruebas funcionales y numéricas. El proceso mantuvo una relación directa entre las fuentes revisadas, los requerimientos del sistema, los módulos implementados y las salidas documentales generadas.

\subsection{Flujo metodológico del análisis estadístico}
\label{sec:flujo_metodologico_estadistico}

El análisis estadístico se definió como una cadena reproducible y no como una selección aislada de modelos. Primero se selecciona una serie oficial del \ICOCIV{} mediante su tabla y ruta jerárquica; luego se valida su estructura temporal, se calculan descriptivos y variaciones, se construye el conjunto de candidatos y se evalúa cada horizonte sin usar información futura. La selección del modelo se decide únicamente por el desempeño fuera de muestra (§\ref{sec:criterios_metricas}); la comparación contra referencias simples es información diagnóstica y la incertidumbre medida informa la evaluación del horizonte, pero ninguna de las dos interviene en la elección del modelo.

Como se observa en la Figura~\ref{fig:flujo_preparacion_modelado}, el proceso empieza en la ruta seleccionada y aplica controles de calidad antes de conformar el conjunto de modelos candidatos. La Figura~\ref{fig:flujo_validacion_exportacion} presenta la etapa posterior: validación walk-forward, métricas por horizonte, decisión de admisibilidad y exportación reproducible.

\afterpage{%
\begin{landscape}
\begin{figure}[H]
\centering
\begin{adjustbox}{max width=\linewidth,max height=0.88\textheight}
\begin{tikzpicture}[
    node distance=0.62cm and 1.35cm,
    paso/.style={rectangle,rounded corners=2pt,draw=black!70,fill=blue!5,align=center,
                 text width=3.6cm,minimum height=0.9cm,inner sep=3pt,font=\scriptsize},
    grupo/.style={rectangle,rounded corners=2pt,draw=black!75,fill=green!8,align=center,
                  text width=4.6cm,minimum height=1.0cm,inner sep=3pt,font=\scriptsize},
    cond/.style={rectangle,rounded corners=2pt,draw=black!75,fill=orange!10,align=center,
                 text width=4.6cm,minimum height=1.0cm,inner sep=3pt,font=\scriptsize},
    flecha/.style={-{Latex[length=2mm]},thick}
]
% --- Columna 1: descendente ---
\node[paso] (inicio) {Inicio};
\node[paso,below=of inicio] (carga) {Carga de datos ICOCIV};
\node[paso,below=of carga] (ruta) {Selección de ruta jerárquica};
\node[paso,below=of ruta] (identifica) {Identificación de la serie};
\node[paso,below=of identifica] (normaliza) {Normalización temporal};
\node[paso,below=of normaliza] (valida) {Validación de datos};

% --- Columna 2: ascendente (arranca abajo, alineada con 'valida') ---
\node[paso,right=of valida] (faltantes) {Revisión de valores faltantes};
\node[paso,above=of faltantes] (duplicados) {Revisión de duplicados};
\node[paso,above=of duplicados] (nopositivos) {Revisión de valores no positivos};
\node[paso,above=of nopositivos] (longitud) {Revisión de longitud mínima};
\node[paso,above=of longitud] (descriptivo) {Análisis descriptivo};
\node[paso,above=of descriptivo] (variaciones) {Variaciones mensuales y acumuladas};
\node[paso,above=of variaciones] (outliers) {Detección robusta de atípicos};
\node[paso,above=of outliers] (grilla) {Construcción de grilla mensual continua};

% --- Columna 3: descendente (arranca arriba, alineada con 'grilla') ---
\node[grupo,right=of grilla] (base)
    {Modelos mínimos obligatorios\\Naive, Drift, Holt lineal y Holt amortiguado};
\node[cond,below=of base] (ampliados)
    {Modelos ampliados condicionales\\OLS, logarítmico, exponencial/log-lineal,\\variación, log-variación y robusto};
\node[paso,below=of ampliados,text width=4.6cm] (salida)
    {Conjunto de modelos candidatos\\para validación temporal};

% --- Conexiones: dentro de cada columna y por los bordes, sin cruces ---
\draw[flecha] (inicio) -- (carga);
\draw[flecha] (carga) -- (ruta);
\draw[flecha] (ruta) -- (identifica);
\draw[flecha] (identifica) -- (normaliza);
\draw[flecha] (normaliza) -- (valida);
\draw[flecha] (valida) -- (faltantes);      % borde inferior, horizontal
\draw[flecha] (faltantes) -- (duplicados);
\draw[flecha] (duplicados) -- (nopositivos);
\draw[flecha] (nopositivos) -- (longitud);
\draw[flecha] (longitud) -- (descriptivo);
\draw[flecha] (descriptivo) -- (variaciones);
\draw[flecha] (variaciones) -- (outliers);
\draw[flecha] (outliers) -- (grilla);
\draw[flecha] (grilla) -- (base);            % borde superior, horizontal
\draw[flecha] (base) -- (ampliados);
\draw[flecha] (ampliados) -- (salida);
\end{tikzpicture}
\end{adjustbox}
\caption{Flujo de preparación, validación y modelado de la serie ICOCIV.}
\label{fig:flujo_preparacion_modelado}
\end{figure}
\end{landscape}
%
\begin{landscape}
\begin{figure}[H]
\centering
\begin{adjustbox}{max width=\linewidth,max height=0.88\textheight}
\begin{tikzpicture}[
    node distance=0.7cm and 1.7cm,
    paso/.style={rectangle,rounded corners=2pt,draw=black!70,fill=blue!5,align=center,
                 text width=3.7cm,minimum height=0.9cm,inner sep=3pt,font=\scriptsize},
    bloque/.style={rectangle,rounded corners=2pt,draw=black!75,fill=green!8,align=center,
                   text width=3.7cm,minimum height=0.9cm,inner sep=3pt,font=\scriptsize},
    decision/.style={diamond,aspect=2.4,draw=black!75,fill=orange!12,align=center,
                     inner sep=1pt,font=\scriptsize},
    negar/.style={rectangle,rounded corners=2pt,draw=red!60!black,fill=red!7,align=center,
                  text width=3.5cm,minimum height=0.9cm,inner sep=3pt,font=\scriptsize},
    flecha/.style={-{Latex[length=2mm]},thick}
]
% --- Columna 1: descendente ---
\node[paso] (bt) {Backtesting walk-forward};
\node[paso,below=of bt] (reestima) {Reestimación del modelo en cada origen \(r\)};
\node[paso,below=of reestima] (proyh) {Proyección por horizonte \(h\)};
\node[paso,below=of proyh] (errores) {Cálculo de errores \(e_{r,h}\)};
\node[bloque,below=of errores] (metricas)
    {Métricas por horizonte: MAE, RMSE, MAPE, sMAPE, MASE, sesgo y errores extremos};

% --- Columna 2: descendente (arranca arriba, a la derecha de 'bt') ---
% B.1, 18-08-2026 (pasada mecanica). Se retira el paso "Evaluacion del
% intervalo del 95% y su ancho relativo" del flujo decisorio: el intervalo es
% diagnostico interno, no publicado, y no interviene entre la comparacion
% contra benchmarks y la seleccion del modelo. `bench` conecta ahora
% directamente con `seleccion`.
\node[paso,right=of bt] (bench) {Comparación contra referencias Naive y Drift};
\node[paso,below=of bench] (seleccion) {Selección del modelo: RMSE OOS mínimo sobre muestra común};
\node[paso,below=of seleccion] (clasifica) {Evidencia fuera de muestra por horizonte (\(W\))};
\node[bloque,below=of clasifica] (estados)
    {Cada horizonte conserva su propia evidencia; un hueco no invalida los posteriores};
\node[decision,below=of estados] (admisible) {¿Horizonte solicitado\\disponible?};

% --- Bifurcación de la decisión y exportación ---
\node[paso,below left=0.9cm and 0.2cm of admisible] (generar) {Generar proyección solicitada};
\node[negar,below right=0.9cm and 0.2cm of admisible] (negar) {Negar proyección y mostrar el motivo};
\node[paso,below=of generar] (pdf) {Exportación PDF, DOCX y HTML};
\node[paso,below=of pdf] (csv) {Exportación CSV reproducible};
\node[paso,below=of csv] (fin) {Fin};

% --- Conexiones ---
\draw[flecha] (bt) -- (reestima);
\draw[flecha] (reestima) -- (proyh);
\draw[flecha] (proyh) -- (errores);
\draw[flecha] (errores) -- (metricas);
% metricas (abajo, col1) -> bench (arriba, col2): canal vertical en el hueco
\draw[flecha] (metricas.east) -- ++(0.85,0) |- (bench.west);
\draw[flecha] (bench) -- (seleccion);
\draw[flecha] (seleccion) -- (clasifica);
\draw[flecha] (clasifica) -- (estados);
\draw[flecha] (estados) -- (admisible);
\draw[flecha] (admisible.south west) -- node[above left,font=\scriptsize,inner sep=1pt]{Sí} (generar.north);
\draw[flecha] (admisible.south east) -- node[above right,font=\scriptsize,inner sep=1pt]{No} (negar.north);
\draw[flecha] (generar) -- (pdf);
\draw[flecha] (pdf) -- (csv);
\draw[flecha] (csv) -- (fin);
% negar -> fin: baja por la derecha y entra a fin por su lado derecho
\draw[flecha] (negar.south) |- (fin.east);
\end{tikzpicture}
\end{adjustbox}
\caption{Flujo de validación walk-forward, selección del modelo y exportación de resultados.}
\label{fig:flujo_validacion_exportacion}
\end{figure}
\end{landscape}
%
}

\subsection{Paso a paso del análisis estadístico realizado por la aplicación}
\label{sec:paso_a_paso_estadistico}

El recorrido completo desde una ruta \ICOCIV{} hasta una salida aceptada o negada se documenta a continuación. Los pasos corresponden al orden lógico del servicio; algunas comprobaciones internas pueden ejecutarse juntas, pero ninguna se omite en la decisión.

\begin{enumerate}[leftmargin=*,label=\textbf{\arabic*.}]
    \item \textbf{Carga de datos ICOCIV.} El usuario selecciona el anexo; el cargador identifica tablas, encabezados mensuales y columnas descriptivas sin alterar el archivo fuente.
    \item \textbf{Selección de ruta jerárquica.} Los desplegables dependientes restringen sucesivamente grupo, subclase, tipología y demás niveles disponibles.
    \item \textbf{Identificación de la serie.} Al completar la ruta se resuelve una única fila fuente, cuya tabla, códigos y descripciones quedan asociados al análisis.
    \item \textbf{Normalización temporal.} Los encabezados de periodo se convierten a año--mes, se ordenan cronológicamente y se vinculan con sus valores numéricos.
    \item \textbf{Validación de datos.} El servicio comprueba que periodos y valores formen una serie mensual procesable antes de habilitar modelos.
    \item \textbf{Revisión de valores faltantes.} Se detectan celdas vacías, conversiones inválidas y meses sin observación; el diagnóstico indica si procede continuar, advertir o bloquear.
    \item \textbf{Revisión de duplicados.} Se exige una observación por ruta y periodo; un mes repetido se informa como error porque vuelve ambigua la serie.
    \item \textbf{Revisión de valores no positivos.} Se marcan valores menores o iguales que cero, pues impiden transformaciones logarítmicas y pueden revelar problemas de lectura.
    \item \textbf{Revisión de longitud mínima.} El número de observaciones se contrasta con el mínimo computacional y con la cantidad necesaria para generar ventanas de evaluación.
    \item \textbf{Análisis descriptivo.} Se calculan tamaño, extremos, media, mediana, desviación estándar, coeficiente de variación y periodos inicial y final.
    \item \textbf{Variaciones mensuales y acumuladas.} Se obtiene el cambio frente al mes anterior y el cambio total entre los extremos de la serie, conservando las fechas reales.
    \item \textbf{Detección robusta de atípicos.} El puntaje modificado basado en mediana y MAD marca variaciones inusuales; los datos oficiales no se eliminan automáticamente.
    \item \textbf{Construcción de grilla mensual continua.} Se genera la secuencia completa de meses entre el primer y último periodo para evidenciar huecos y alinear observaciones.
    \item \textbf{Definición de modelos candidatos.} Compiten los diez candidatos del catálogo ---Naive y Drift como referencias, más las regresiones, las transformaciones, el modelo robusto y los suavizamientos---. No hay activación condicionada: la única condición de elegibilidad es que el modelo pueda estimarse con las observaciones disponibles (§\ref{sec:criterios_catalogo}).
    \item \textbf{Validación walk-forward.} Para cada modelo y horizonte se simulan decisiones históricas mediante orígenes sucesivos que respetan el orden pasado--futuro \parencite{tashman2000}.
    \item \textbf{Reestimación por origen.} En cada origen \(r\), los parámetros se ajustan únicamente con observaciones disponibles hasta \(r\); no se reutiliza el ajuste hecho con toda la serie.
    \item \textbf{Pronóstico por horizonte.} El modelo reestimado produce \(\hat y_{r+h}\) para cada horizonte \(h\) evaluable.
    \item \textbf{Cálculo de errores.} Cada pronóstico se contrasta con el dato observado mediante \(e_{r,h}=y_{r+h}-\hat y_{r+h}\).
    \item \textbf{Métricas por horizonte.} Los errores se agregan por \(h\) para obtener MAE, RMSE, MAPE, sMAPE, MASE y sesgo medio. Se informa además cuántas ventanas presentan un error inusual según el puntaje z modificado, dato descriptivo que no interviene en ninguna decisión.
    \item \textbf{Comparación contra benchmarks.} Naive y Drift compiten en igualdad de condiciones con los demás candidatos y sirven además de referencia de lectura: el error relativo frente a ellos se publica como información y \textbf{no sustituye} al modelo elegido ni condiciona su aceptación.
    \item \textbf{Evaluación de intervalos (diagnóstico interno, no publicado).} Se construye internamente el intervalo de predicción del 95\,\% con los errores del horizonte exacto de cada paso y se calcula su ancho relativo como señal de incertidumbre.
    \item \textbf{Selección del modelo.} Se elige un único modelo por serie: el que minimiza el error cuadrático medio fuera de muestra sobre la muestra común a todos los candidatos, de modo que la trayectoria no dependa del horizonte solicitado. La regla no contiene umbrales ni ponderaciones elegidas.
    \item \textbf{Clasificación de horizontes.} Cada \(h\) se clasifica con la evidencia fuera de muestra de ese mismo horizonte: cuántas ventanas reúne y qué error midió. La estabilidad y los diagnósticos residuales se informan y no degradan el estado; un límite de interfaz no es un máximo estadístico.
    \item \textbf{Recomendación o restricción.} El resultado distingue horizonte recomendado, escenario con advertencia, primer horizonte no viable o rechazo, y explica el criterio aplicado.
    \item \textbf{Generación o negación.} Si el pronóstico puntual es calculable, se genera y se publica, acompañado del estado metodológico que declara qué no puede sostenerse. Solo las imposibilidades aritméticas del propio punto ---valor no finito--- niegan la salida: una deficiencia del intervalo o una carencia de evidencia no la bloquean por sí solas, conforme al principio de que una deficiencia del intervalo no invalida automáticamente el pronóstico puntual. El motivo se conserva siempre para revisión.
    \item \textbf{Exportación de resultados.} La interfaz y los informes presentan ruta, serie, modelo, métricas, horizonte y advertencias de forma coherente. \textbf{No presentan intervalos}: esta versión los retiró de todas las superficies públicas.
    \item \textbf{Reproducibilidad.} CSV y reportes conservan archivo, fila fuente, corte, observaciones, parámetros, periodos y valores proyectados para repetir la salida con la misma versión.
\end{enumerate}

Este recorrido enlaza explícitamente la selección jerárquica con el resultado. Si la serie no supera una validación crítica o el horizonte no es admisible, el flujo termina con una negación explicada; la aplicación no completa el análisis con valores inventados.

\subsection{Criterios para la conformación del catálogo de modelos}
\label{sec:criterios_catalogo}

La aplicación no evalúa el universo de métodos de pronóstico disponibles en la literatura, sino un \textbf{catálogo acotado} de diez candidatos. Esa delimitación es una decisión metodológica previa al análisis y, por tanto, debe justificarse antes de describir cómo compiten los modelos entre sí. Conviene enunciar primero lo que la delimitación \emph{no} afirma: no se sostiene que estos sean los mejores métodos existentes, ni que los métodos excluidos sean inferiores para series \ICOCIV{}. Afirmar cualquiera de las dos cosas exigiría una evaluación comparativa que este trabajo no realizó.

Lo que sí se sostiene es que el conjunto adoptado es \emph{compatible} con el problema, con los datos disponibles y con el alcance del proyecto. Diez consideraciones lo determinan.

\begin{enumerate}[leftmargin=*,label=\textbf{\arabic*.}]
    \item \textbf{Naturaleza univariada de las series.} Cada ruta \ICOCIV{} se analiza como una única serie mensual de índices. El problema que la aplicación resuelve es extrapolar esa serie a partir de su propia historia, de modo que los candidatos son métodos univariados de pronóstico.
    \item \textbf{Longitud limitada de la historia disponible.} Las series del anexo vigente tienen sesenta y cinco observaciones mensuales. Con esa cantidad, los métodos con muchos parámetros libres disponen de poca información por parámetro, y su ventaja aparente dentro de muestra no se traduce necesariamente en mejor desempeño fuera de ella \parencite{hyndman2021fpp3}.
    \item \textbf{Reestimación repetida dentro del esquema de validación.} Cada modelo no se ajusta una sola vez: se reestima en cada origen de la validación temporal, para cada horizonte evaluado y para cada candidato del catálogo. Un método cuyo ajuste dependa de búsquedas costosas o de intervención manual no puede someterse a ese esquema de forma automática y reproducible.
    \item \textbf{Interpretabilidad técnica de la salida.} El resultado se emplea en el análisis de variaciones de precios de contratos de obra civil, donde el profesional responsable debe poder explicar de dónde procede una cifra. Todos los candidatos admiten una ecuación explícita y parámetros con lectura directa.
    \item \textbf{Reproducibilidad.} La misma entrada debe producir el mismo modelo, los mismos parámetros y el mismo pronóstico. Los candidatos adoptados tienen estimación determinista, sin inicializaciones aleatorias ni dependencia del orden de ejecución.
    \item \textbf{Parsimonia respecto de la información disponible.} El catálogo cubre desde un parámetro efectivo, en Naive, hasta cinco, en Holt amortiguado. Ese rango es proporcionado a la longitud de las series y evita que la comparación deba decidirse entre métodos que la muestra no alcanza a distinguir.
    \item \textbf{Inclusión de referencias simples.} Naive y Drift se incorporan como \emph{benchmarks}: sin una referencia elemental no es posible afirmar que un método más elaborado aporte algo. Comparar contra un método ingenuo es práctica estándar en la evaluación de pronósticos \parencite{hyndman_koehler2006}.
    \item \textbf{Cobertura de hipótesis distintas sobre la dinámica.} Los candidatos no son variantes de una misma idea: representan supuestos diferentes sobre cómo evoluciona la serie, según se detalla en la Tabla~\ref{tab:catalogo_hipotesis}.
    \item \textbf{Aplicabilidad uniforme del mismo esquema de validación.} Todas las rutas \ICOCIV{} ---agregados nacionales, grupos de obra, subclases, tipologías e insumos--- se evalúan con el mismo procedimiento de origen móvil, sin ajustes por ruta. Un catálogo homogéneo es condición para que esa comparación sea la misma en todos los casos.
    \item \textbf{Alcance de un trabajo de grado.} Evaluar exhaustivamente la literatura de pronóstico excede el propósito de este trabajo, que es construir una herramienta defendible y reproducible para un uso institucional concreto.
\end{enumerate}

En síntesis, el catálogo no corresponde a una acumulación arbitraria de algoritmos, sino a un conjunto parsimonioso de candidatos que representa diferentes hipótesis plausibles sobre el comportamiento de las series y que puede estimarse, interpretarse y compararse de forma reproducible mediante validación temporal fuera de muestra.

\subsubsection*{Hipótesis que representa cada candidato}

\begin{table}[H]
\centering
\caption{Catálogo de modelos: hipótesis representada y papel en la aplicación}
\label{tab:catalogo_hipotesis}
\scriptsize
\begin{adjustbox}{max width=\textwidth}
\begin{tabular}{>{\raggedright\arraybackslash}p{0.16\textwidth}>{\raggedright\arraybackslash}p{0.27\textwidth}>{\raggedright\arraybackslash}p{0.30\textwidth}>{\raggedright\arraybackslash}p{0.17\textwidth}}
\toprule
\textbf{Modelo} & \textbf{Hipótesis representada} & \textbf{Razón de inclusión} & \textbf{Papel en SAVIP} \\
\midrule
Naive & El mejor pronóstico del futuro es el último valor observado. & Referencia elemental: cualquier método debe justificar su complejidad frente a ella. & Benchmark \\
Drift & La serie continúa con el cambio promedio observado entre su primer y su último dato. & Referencia tendencial mínima, en forma cerrada y sin optimización. & Benchmark \\
Lineal (OLS) & Existe una tendencia determinista aproximadamente lineal en el tiempo. & Interpretabilidad paramétrica directa: la pendiente tiene unidades de índice por mes. & Candidato \\
Logarítmica temporal & La tendencia crece, pero a ritmo decreciente. & Cubre la hipótesis de desaceleración, que una recta no puede representar. & Candidato \\
Exponencial / log-lineal & El crecimiento es aproximadamente proporcional al nivel de la serie. & Hipótesis natural en índices de precios; exige serie estrictamente positiva. & Candidato sobre transformaciones \\
Huber & Hay tendencia lineal, pero algunas observaciones no deben influir en la misma medida. & Acomoda observaciones extremas \textbf{sin eliminar ningún dato oficial} \parencite{huber1964}. & Candidato robusto \\
Holt lineal & El nivel y la tendencia se actualizan con cada nueva observación. & Tendencia dinámica, no fijada de una vez por toda la historia como en una regresión. & Candidato de suavizamiento \\
Holt amortiguado & Como el anterior, pero la extrapolación de la tendencia pierde intensidad con el horizonte. & Evita proyectar indefinidamente una pendiente reciente; respaldado en horizontes largos \parencite{gardner_mckenzie1985,gardner2006}. & Candidato de suavizamiento \\
Variación lineal & Lo que sigue una trayectoria regular es el cambio absoluto mensual, no el nivel. & Traslada la hipótesis de tendencia a la primera diferencia de la serie. & Candidato sobre transformaciones \\
Log-variación & Lo regular es el cambio relativo o porcentual mensual. & Complementa al anterior cuando la variación es proporcional al nivel. & Candidato sobre transformaciones \\
\bottomrule
\end{tabular}
\end{adjustbox}
\end{table}

\textbf{Elegibilidad dentro del catálogo.} Los diez candidatos compiten siempre que puedan estimarse. La única condición de elegibilidad es la \emph{estimabilidad matemática} con los datos disponibles: cada modelo declara un mínimo operativo \(N_{\min}(m)\) y queda excluido, con su motivo registrado, cuando la serie no lo alcanza. No existe activación condicionada por horizonte, por longitud de la serie ni por alertas de diagnóstico: una regla de ese tipo decidiría qué modelos pueden ganar, y establecerla sin sustento equivaldría a decidir el resultado.

Conviene precisar qué es y qué no es ese mínimo. \(N_{\min}(m)\) es el número de observaciones por debajo del cual el modelo \(m\) \emph{no puede intentarse}; no significa suficiencia estadística, ni validación, ni estabilidad garantizada, ni un mínimo universal. Que una serie alcance \(N_{\min}(m)\) habilita el intento de ajuste y nada más: la evidencia fuera de muestra se mide después y por separado. Tampoco procede de una regla única aplicada a todas las familias, sino de tres justificaciones distintas, según cómo obtiene sus parámetros cada modelo (tabla \ref{tab:nmin_por_modelo}).

\begin{table}[H]
\centering
\caption{Mínimo operativo de estimabilidad por modelo y justificación de la que procede}
\label{tab:nmin_por_modelo}
\scriptsize
\begin{adjustbox}{max width=\textwidth}
\begin{tabular}{lcp{0.62\linewidth}}
\toprule
\textbf{Modelo} & \textbf{\(N_{\min}\)} & \textbf{Justificación} \\
\midrule
Naive & 1 & Forma cerrada \(\hat{y}=y_T\): no minimiza nada y no estima ningún parámetro; basta que exista la última observación. \\
Drift & 2 & Forma cerrada \((y_T-y_1)/(T-1)\): exige las dos observaciones extremas y que sus tiempos difieran. \\
Lineal, Logarítmico & 3 & Ajuste por mínimos cuadrados con \(k=2\). Con \(n=2\) la recta interpola exactamente y el residuo es nulo, de modo que se requiere \(n \geq k+1\). \\
Exponencial log-lineal & 3 & Mínimos cuadrados sobre \(\ln y\) con \(k=2\); el factor de corrección de retransformación es función de los residuos, no un parámetro libre adicional. \\
Huber & 4 & El estimador ajusta además la \emph{escala} junto con el intercepto y la pendiente \parencite{sklearn_huberregressor}, de modo que \(k=3\) y \(n \geq 4\). Es un caso en que contar solo los coeficientes de regresión subestimaría el mínimo. \\
Variación mensual, Log-variación & 4 & Requisito propio del modelo: la serie se transforma en variaciones, y se exigen al menos tres variaciones finitas, lo que implica cuatro observaciones. No se deriva del recuento de parámetros. \\
Holt lineal & 5 & Se estiman \(\alpha\), \(\beta^{*}\) y los dos estados iniciales \(l_0\) y \(b_0\) minimizando la suma de cuadrados de un paso \parencite{hyndman2021fpp3}, luego \(k=4\) y \(n \geq 5\). \\
Holt amortiguado & 6 & Los cuatro anteriores más \(\varphi\), luego \(k=5\) y \(n \geq 6\). Es el máximo del catálogo y, por tanto, el que fija el primer origen del backtesting. \\
\bottomrule
\end{tabular}
\end{adjustbox}
\end{table}

Las tres justificaciones no son intercambiables. Aplicar \(n \geq k+1\) a Naive o a Drift carecería de sentido, porque no resuelven ningún problema de minimización; aplicarla a los modelos sobre variaciones daría un valor menor que el que su propia transformación exige; y contar en Huber solo los dos coeficientes de regresión, ignorando la escala, produciría un mínimo insuficiente. El criterio general es contar todo parámetro que el procedimiento estime y exigir una observación más, allí donde el modelo efectivamente estima.

\textbf{Métodos considerados y no adoptados.} Dos candidatos que llegaron a implementarse fueron retirados del catálogo ---el promedio móvil y la variación reciente---, ambos por depender de una ventana de seis periodos sin fuente identificada; un parámetro sin sustento no puede decidir un resultado publicado. La familia ETS se estudió e implementó, y se retiró tras comprobar que duplicaba exactamente la trayectoria de Holt amortiguado bajo la configuración adoptada; queda registrada como trabajo futuro, condicionada a una estimación completa de los parámetros de espacio de estados. La combinación de pronósticos se evaluó y no se adoptó, tras medir que no aportaba bajo el esquema de un modelo único por serie.

Otros métodos ampliamente utilizados ---entre ellos la familia ARIMA y SARIMA--- \textbf{no fueron evaluados sistemáticamente} en este trabajo y quedaron fuera del catálogo delimitado. No se afirma que sean inadecuados para series \ICOCIV{}: su incorporación exigiría ampliar el alcance metodológico e implicaría, como con cualquier otro candidato, someterlos al mismo esquema de validación fuera de muestra antes de admitirlos.

El catálogo adoptado no pretende abarcar exhaustivamente los métodos de pronóstico existentes ni demostrar que los métodos excluidos sean inferiores. Constituye una delimitación metodológica compatible con las características de los datos, los requisitos de interpretabilidad y reproducibilidad, y el alcance del proyecto.

\subsection{Criterios para la selección de métricas y del modelo final}
\label{sec:criterios_metricas}

Conviene distinguir con precisión tres conceptos que el lenguaje corriente confunde, porque solo el tercero determina qué modelo entrega la aplicación.

\begin{description}[leftmargin=*]
    \item[Parámetros del modelo.] Son las cantidades que definen cada método ---\(\beta_0\) y \(\beta_1\) en una regresión; \(\alpha\), \(\beta^{*}\), \(\phi\) y los estados iniciales de nivel y tendencia en los suavizamientos--- y se estiman mediante el procedimiento propio de cada formulación. \textbf{No son métricas de error.}
    \item[Métricas de desempeño.] Resumen, con dimensiones distintas, la magnitud de los errores cometidos fuera de muestra. Se reportan todas porque ninguna es suficiente por sí sola.
    \item[Criterio de selección.] Es la función objetivo, única, con la que se decide qué modelo produce la trayectoria de una serie.
\end{description}

\subsubsection*{Las métricas y su función}

Las métricas adoptadas son complementarias y cubren dimensiones que no se sustituyen entre sí. MAE expresa la magnitud promedio del error en unidades del índice y es la lectura más directa para el usuario. RMSE expresa esa magnitud penalizando en mayor medida las desviaciones grandes, por tratarse de una media cuadrática. MAPE ofrece una lectura porcentual relativa, con la limitación conocida de no estar definido cuando algún valor observado es cero o próximo a cero. sMAPE atenúa parte de la asimetría de MAPE y se interpreta junto a él. MASE escala el error frente a una referencia \emph{naive} calculada dentro de cada ventana de entrenamiento, lo que lo hace comparable entre series de escalas distintas \parencite{hyndman_koehler2006}; su lectura frente a la unidad es informativa y \textbf{no constituye un veto}: un valor mayor que uno se comunica como advertencia y no cancela la proyección. El sesgo medio, por último, no mide magnitud sino dirección, y revela la tendencia sistemática a sobrestimar o subestimar que las métricas de valor absoluto ocultan por construcción.

\begin{table}[H]
\centering
\caption{Métricas reportadas y su papel en la decisión}
\label{tab:metricas_papel}
\scriptsize
\begin{adjustbox}{max width=\textwidth}
\begin{tabular}{>{\raggedright\arraybackslash}p{0.11\textwidth}>{\raggedright\arraybackslash}p{0.28\textwidth}>{\raggedright\arraybackslash}p{0.33\textwidth}>{\raggedright\arraybackslash}p{0.16\textwidth}}
\toprule
\textbf{Métrica} & \textbf{Dimensión que representa} & \textbf{Uso en SAVIP} & \textbf{¿Define el modelo seleccionado?} \\
\midrule
MAE & Magnitud absoluta promedio, en unidades del índice. & Lectura directa del error para el usuario. & No \\
RMSE & Magnitud con penalización mayor de los errores grandes. & Evaluación del desempeño \textbf{y función objetivo de la selección}. & \textbf{Sí} \\
MAPE & Error relativo porcentual. & Lectura comparable entre escalas; no definido con observados nulos. & No \\
sMAPE & Error relativo porcentual simétrico. & Complemento de MAPE ante asimetría. & No \\
MASE & Error escalado frente a una referencia \emph{naive}. & Comparación entre series de distinta escala; advertencia, no veto. & No \\
Sesgo medio & Dirección sistemática del error. & Diagnóstico de sobrestimación o subestimación. & No \\
\bottomrule
\end{tabular}
\end{adjustbox}
\end{table}

\textbf{Las métricas no se combinan.} No existe en la aplicación una puntuación compuesta del tipo \(w_1\,\mathrm{RMSE}+w_2\,\mathrm{MAPE}+w_3\,\mathrm{MASE}\). Una combinación de ese tipo obligaría a justificar los pesos, y unos pesos elegidos sin fuente decidirían el modelo entregado. Las métricas distintas del RMSE se publican como información de desempeño y no intervienen en la elección. Los criterios de información dentro de muestra, como el AICc, y las medidas de bondad de ajuste, como el \(R^2\), se reportan como diagnósticos y \textbf{no participan en la selección predictiva}: describen el ajuste a los datos ya observados, no la capacidad de pronosticar los que aún no lo están.

\subsubsection*{Por qué el RMSE fuera de muestra define el modelo seleccionado}

Adoptar el RMSE como función objetivo no supone afirmar que sea la mejor métrica en abstracto, sino que es la adecuada para la decisión concreta que debe tomarse. Seis razones lo sustentan.

Primero, es una medida estándar de exactitud predictiva, de uso extendido y lectura conocida. Segundo, los candidatos se comparan siempre \emph{dentro de la misma serie} y, por tanto, sobre la misma escala; en esa situación las medidas dependientes de escala son las apropiadas, mientras que las relativas o escaladas están pensadas para comparar entre series distintas \parencite{hyndman_koehler2006,hyndman2021fpp3}. Tercero, la pérdida cuadrática penaliza en mayor medida los errores grandes, lo que es coherente con un uso presupuestal en el que una desviación considerable tiene consecuencias mayores que varias pequeñas. Cuarto, se calcula sobre errores genuinamente fuera de muestra, obtenidos por validación de origen móvil con reestimación en cada corte \parencite{tashman2000,hyndman2021fpp3}, y no sobre el ajuste a los datos con los que el modelo fue estimado. Quinto, todos los candidatos se evalúan sobre una muestra común, según se justifica más adelante. Sexto, es un criterio único, lo que evita introducir una función multicriterio cuyos pesos no podrían sustentarse.

\textbf{Equivalencia con la suma de cuadrados.} La implementación compara los candidatos mediante la suma de errores al cuadrado, calculada con aritmética exacta para que dos modelos no queden empatados por efecto del redondeo. La equivalencia es inmediata: si \(S\) es la muestra común y \(\mathrm{SSE}_m=\sum_{(r,h)\in S}e_{m,r,h}^{2}\), entonces

\begin{equation}
\mathrm{RMSE}_m=\sqrt{\frac{\mathrm{SSE}_m}{|S|}},
\label{eq:rmse_sse}
\end{equation}

y como \(|S|\) es idéntico para todos los candidatos por construcción de la muestra común, la raíz y el divisor son transformaciones monótonas comunes, de modo que \(\arg\min_m \mathrm{RMSE}_m=\arg\min_m \mathrm{SSE}_m\). \textbf{Esta equivalencia vale únicamente cuando el número de pares es el mismo para todos los candidatos}: con muestras de distinto tamaño, minimizar la suma de cuadrados no equivale a minimizar el error cuadrático medio.

\subsubsection*{Por qué la comparación se hace sobre una muestra común}

La muestra de evaluación es el conjunto de pares \((\text{periodo objetivo},\ \text{horizonte})\) en los que \emph{todos} los candidatos producen un error finito. Sin esa restricción, un modelo que solo pudiera pronosticar en los orígenes más tardíos ---o en un número distinto de casos--- se compararía sobre un subconjunto diferente del de los demás, y su ventaja podría deberse a que enfrentó casos más fáciles, o menos casos, y no a que pronostique mejor. La muestra común elimina esa posibilidad: la comparación se realiza sobre exactamente los mismos objetivos y con exactamente el mismo número de contrastes.

\subsubsection*{Por qué se selecciona un único modelo por serie}

La selección no busca un ganador distinto para cada horizonte. Si el modelo se eligiera por horizonte, los primeros meses de la trayectoria cambiarían según el horizonte final que el usuario solicitara, aunque la información histórica disponible fuese idéntica: pedir tres meses y pedir doce devolvería valores distintos para el mismo mes. Como el modelo entregado debe producir una trayectoria mensual coherente desde el mes siguiente al último dato observado hasta el objetivo solicitado, la muestra de evaluación pertinente es exactamente el conjunto de pronósticos que ese modelo tendrá que producir: todos los orígenes de la validación por todos los horizontes entregables. De ahí que el RMSE se agregue sobre ese conjunto y que el resultado sea un modelo único, que constituye una propiedad de la serie y no del horizonte consultado.


\subsection{Regla de trazabilidad y clasificación de criterios}
\label{sec:trazabilidad_criterios}

Cada valor presentado por la aplicación debe responder a cinco preguntas: de qué archivo y ruta provino, qué transformación recibió, qué función lo calculó, con qué criterio se interpretó y qué limitación acompaña su uso. Para no presentar decisiones internas como leyes estadísticas, los criterios se clasificaron en cuatro tipos. La matriz completa se presenta más adelante en esta misma sección (Sección~\ref{sec:matriz_criterios}).

\begin{table}[H]
\centering
\caption{Clasificación usada para la trazabilidad de criterios}
\label{tab:clasificacion_criterios}
\scriptsize
\begin{adjustbox}{max width=\textwidth}
\begin{tabularx}{\textwidth}{@{}p{3.4cm}YYY@{}}
\toprule
\textbf{Tipo} & \textbf{Origen} & \textbf{Uso permitido} & \textbf{Forma de comunicación} \\
\midrule
Bibliográfico & Fuente académica o técnica identificable. & Sustentar conceptos, fórmulas y umbrales publicados. & Cita y referencia bibliográfica. \\
Estadístico estándar & Convención matemática o distribución conocida. & Calcular métricas, cuantiles o pruebas bajo sus supuestos. & Fórmula, supuestos y cautelas. \\
Operativo interno & Parámetro configurable del software, sin fuente externa. & Solo aspectos de interfaz, cómputo, formato o control técnico que \textbf{no} cambian el modelo, el pronóstico, el horizonte, la admisibilidad, el intervalo, la cobertura ni los valores publicados. Un parámetro interno que sí cambia alguno de esos elementos requiere sustento, estimación o derivación; la etiqueta ``interno'' no habilita por sí sola una decisión. & Identificarlo expresamente como criterio interno, no universal, y declarar si tiene o no efecto decisorio. \\
Regla institucional & Documento, lineamiento o requisito del proceso de obra. & Aplicar empalme, periodos de corte, variables y exportables. & Citar el documento y conservar variables de entrada. \\
\bottomrule
\end{tabularx}
\end{adjustbox}
\end{table}

\subsection{Matriz de trazabilidad conceptual}
\label{sec:matriz_trazabilidad_conceptual}

La Tabla~\ref{tab:matriz_conceptual} enlaza definición, finalidad, implementación e interpretación. Su función es impedir que una cifra se presente de manera aislada: cada componente debe responder qué calcula, dónde interviene y qué límite conserva.

\begin{landscape}
\begin{table}[H]
\centering
\caption{Matriz de trazabilidad conceptual del análisis estadístico}
\label{tab:matriz_conceptual}
\setstretch{1}
\tiny
\begin{adjustbox}{max width=\linewidth,max totalheight=0.84\textheight}
\begin{tabularx}{\linewidth}{@{}p{3.1cm}YYY@{}}
\toprule
\textbf{Componente} & \textbf{Definición y propósito} & \textbf{Implementación en la aplicación} & \textbf{Interpretación y referencia} \\
\midrule
Validación temporal & Comprueba que la serie represente meses reales, consecutivos y ordenados. & Normaliza periodos a año--mes y revisa duplicados, saltos, faltantes, dominio y archivo fuente. & Sin continuidad, un salto de calendario puede confundirse con una variación económica. \\
Estadística descriptiva & Resume nivel, dispersión y recorrido histórico. & Calcula observaciones, extremos, media, mediana, desviación, coeficiente de variación y periodos. & Describe la serie antes de pronosticar; no sustituye validación predictiva. \\
Variación mensual & Mide el cambio porcentual frente al mes anterior. & Calcula \((y_t/y_{t-1}-1)\times100\) sobre periodos consecutivos. & Permite localizar aceleraciones, saltos y episodios inusuales. \\
Atípicos robustos & Identifica cambios inusuales sin depender de media y desviación estándar. & Usa puntaje modificado con mediana y MAD sobre variaciones. & Marca y comunica; no elimina automáticamente el valor oficial \parencite{nist_outliers}. \\
Naive & Benchmark que conserva el último valor observado. & Para \(h\) meses usa \(\hat y_{T+h}=y_T\). & Un modelo complejo que no lo supera ofrece poca ganancia predictiva \parencite{hyndman2021fpp3}. \\
Drift & Benchmark que prolonga el cambio promedio histórico. & Extiende \((y_T-y_1)/(T-1)\) al horizonte \(h\). & Proporciona una tendencia simple, reproducible y comparable. \\
Holt lineal y amortiguado & Suavizamiento con nivel y tendencia. & Estima nivel, pendiente y, cuando aplica, amortiguación; reentrena en cada origen. & La amortiguación limita extrapolaciones lineales prolongadas \parencite{holt1957,gardner_mckenzie1985}. \\
Regresiones tradicionales & Modelos interpretables sobre el nivel del índice. & Evalúa OLS lineal, logarítmica temporal, log-lineal y ajuste robusto cuando los datos lo permiten. & Son referencias predictivas, no evidencia causal; se validan fuera de muestra \parencite{granger_newbold1974,huber1964}. \\
Walk-forward & Validación fuera de muestra respetando pasado--futuro. & En cada origen entrena con datos anteriores y predice un horizonte específico. & Simula el uso real y evita fuga de información \parencite{tashman2000}. \\
Métricas de error & Cuantifican magnitud, penalización y escala del error. & Calcula MAE, RMSE, MAPE, sMAPE, MASE, sesgo y comparaciones por horizonte. & El RMSE fuera de muestra sobre la muestra común es la función objetivo que decide el modelo seleccionado (§\ref{sec:criterios_metricas}); las demás son complementarias y se reportan como información \parencite{hyndman_koehler2006}. \\
Diagnóstico residual & Busca sesgo y estructura no capturada por el modelo. & Contrasta media residual nula (prueba \(t\)), autocorrelación conjunta (Ljung--Box), varianza constante (Breusch--Pagan) y normalidad (Jarque--Bera); Durbin--Watson se reporta como descriptor. & Informativa: no modifica el modelo, el pronóstico ni el horizonte. \\
Intervalo de predicción 95\,\% (retirado de las salidas) & Comunicaba la incertidumbre alrededor del pronóstico puntual. & Empleaba errores fuera de muestra del horizonte exacto de cada paso, con cuantil empírico corregido o respaldo $t$. & No se publica en esta versión: la construcción completa carece de sustento verificable y permanece solo como diagnóstico interno. \\
Evidencia por horizonte (\(W\)) & Determina si existe error fuera de muestra para ese \(h\). & \(W=n-N_0-h+1\); \(W=0\) sin evidencia, \(W\in\{1,2\}\) evidencia muy limitada, \(W\geq3\) evidencia disponible. & \(W\) es existencia, no calidad; ningún valor de \(W\) certifica ni valida un horizonte. Los accesos de interfaz no constituyen un límite estadístico fijo. \\
Reproducibilidad & Permite reconstruir una proyección. & Conserva archivo, ruta, fila, corte, datos, modelo, parámetros, métricas y salida. & Una cifra sin estos metadatos no constituye evidencia completa. \\
\bottomrule
\end{tabularx}
\end{adjustbox}
\end{table}
\end{landscape}

\subsection{Matriz de criterios estadísticos}
\label{sec:matriz_criterios}

Los umbrales y reglas de decisión de la aplicación están centralizados en un módulo interno de criterios, donde cada criterio se registra con su tipo, su valor, su sustento, su justificación y la acción de gobierno metodológico que le corresponde. A partir de ese registro se genera automáticamente la Tabla~\ref{tab:matriz_criterios}, de modo que el documento y la aplicación no pueden divergir sin que la diferencia sea visible: si un umbral cambia en la implementación, la tabla del documento cambia al regenerarse. La columna de tipo es la salvaguarda central del enfoque, porque solo los criterios marcados como bibliográficos o estándares invocan una fuente externa, mientras que los operativos internos se declaran expresamente como parámetros propios y configurables, nunca como reglas estadísticas universales.

Cada criterio se conecta con una etapa concreta del flujo y con un efecto sobre la decisión. Los criterios de preparación de datos determinan si una serie es siquiera modelable y con cuántas ventanas de validación. Los criterios de atípicos y de cambio de año solo marcan o ajustan bajo condiciones acumulativas verificables, sin alterar el dato oficial de forma silenciosa. El criterio de error que decide la selección del modelo es único: el RMSE fuera de muestra sobre la muestra común a todos los candidatos (§\ref{sec:criterios_metricas}). Los criterios de sesgo, estabilidad y comparación contra referencias simples no intervienen en esa selección; se publican como advertencias descriptivas, y ninguna métrica actúa como veto. Los criterios de horizonte traducen la evidencia fuera de muestra medida en la clasificación que aprueba, advierte, restringe a escenario o niega una proyección; los del intervalo dejaron de intervenir cuando el intervalo se retiró de la salida. Cada fórmula asociada a estos criterios aparece numerada en el marco teórico y está verificada mediante pruebas automatizadas; en particular, las fórmulas de error y el ajuste de cambio de año cuentan con pruebas específicas cuyos resultados se presentan en el capítulo de pruebas.

\begin{landscape}
% Tabla generada automaticamente desde la matriz de criterios del codigo.
% No editar a mano; regenerar con genera_tabla_criterios.py (misma carpeta).
{\scriptsize
\begin{longtable}{@{}>{\raggedright\arraybackslash}p{0.238\linewidth}>{\raggedright\arraybackslash}p{0.152\linewidth}>{\raggedright\arraybackslash}p{0.266\linewidth}>{\raggedright\arraybackslash}p{0.295\linewidth}@{}}
\caption{Matriz de criterios estad\'isticos de la aplicaci\'on, generada autom\'aticamente desde el m\'odulo interno de criterios.}
\label{tab:matriz_criterios}\\
\toprule
\textbf{Criterio} & \textbf{Tipo} & \textbf{Valor o umbral} & \textbf{Sustento} \\
\midrule
\endfirsthead
\multicolumn{4}{l}{\emph{Continuaci\'on de la Tabla~\ref{tab:matriz_criterios}}}\\
\toprule
\textbf{Criterio} & \textbf{Tipo} & \textbf{Valor o umbral} & \textbf{Sustento} \\
\midrule
\endhead
\bottomrule
\endlastfoot
Mínimo computacional para modelacion & operativo interno, sin fuente externa & 8 & Parámetro interno documentado \\
Advertencia por serie corta & operativo interno, sin fuente externa & 18 & Parámetro interno documentado \\
Longitud recomendada para modelacion mensual & operativo interno, sin fuente externa & 24 & Criterio metodológico interno \\
Mínimo de iteraciones walk-forward (descriptivo) & operativo interno, sin fuente externa & 6 & Parámetro interno documentado; ninguna fuente fija el valor 6 \\
Primer origen de la validación temporal & pendiente de decisión & N0 = max sobre los candidatos de su mínimo operativo codificado para intentar el ajuste (hoy 6, por Holt amortiguado) & Tashman (2000) y FPP3 §5.10 sustentan el esquema de origen móvil; ninguna fuente determina el primer origen \\
Valor atípico por Z robusto modificado & bibliográfico & |M\_\allowbreak{}i| $>$ 3.5 & NIST/SEMATECH e-Handbook of Statistical Methods, sección 1.3.5.17 \\
Clasificación de alertas por periodo (deduplicadas) & operativo interno, sin fuente externa & patron\_\allowbreak{}calendario / posible\_\allowbreak{}error\_\allowbreak{}datos / posible\_\allowbreak{}cambio\_\allowbreak{}nivel / posible\_\allowbreak{}atipico\_\allowbreak{}aislado & Politica interna; el patron calendario se confirma con C-CAL-001 \\
Subrangos leve/moderado/fuerte de atípicos & operativo interno, sin fuente externa & 3.5-4.5; 4.5-5.0; >5.0 & Sin calibración empírica documentada \\
Durbin-Watson como estadistico descriptivo & bibliográfico & Se reporta el valor en [0, 4]; sin cortes de clasificacion & Durbin y Watson (1951), Biometrika 38(1-2), pp. 159-177 \\
Nivel de significancia Ljung-Box/Jarque-Bera & bibliográfico & alpha = 0.05 & Convencion estadística usual al 5\%; Ljung y Box; Jarque y Bera \\
Heterocedasticidad de los residuos & bibliográfico & Breusch-Pagan sobre [1, t]; se rechaza si p $<$ 0.05 & Breusch y Pagan (1979); statsmodels.stats.diagnostic.het\_\allowbreak{}breuschpagan \\
Media residual igual a cero & bibliográfico & Prueba t de una muestra bilateral; se rechaza si p $<$ 0.05 & Prueba t de una muestra; scipy.stats.ttest\_\allowbreak{}1samp \\
MASE como señal de cautela & bibliográfico & > 1 frente a escala naive in-sample & Hyndman y Koehler (2006) \\
Errores de backtesting inusuales por puntaje z modificado & bibliográfico & |M\_\allowbreak{}i| $>$ 3.5 con M\_\allowbreak{}i = 0.6745 (|e\_\allowbreak{}i| - mediana) / MAD & Iglewicz y Hoaglin (1993); NIST/SEMATECH e-Handbook 1.3.5.17 \\
Construccion del intervalo de prediccion (FPP3 5.5) & bibliográfico & pronostico +/- c * sigma\_\allowbreak{}h, con sigma\_\allowbreak{}h = sqrt(SUM e\textasciicircum{}2 / n) sobre los errores OOS del paso exacto y c el cuantil t con n grados de libertad del nivel nominal & Hyndman y Athanasopoulos (2021), FPP3 5.5: yhat +/- c*sigma\_\allowbreak{}h, con sigma\_\allowbreak{}h estimada de los errores. El multiplicador t es la derivacion exacta cuando sigma\_\allowbreak{}h se estima \\
Condicion de existencia del cuantil de orden al 95 \% & derivación matemática & k = ceil((n+1)(1-alfa)) $\leq$ n, es decir n $\geq$ 1/alfa - 1 = 19 errores al 95 \% & Lei et al. (2018); Angelopoulos y Bates (2023) \\
Cobertura empírica por evaluacion de origen movil & operativo técnico & cada error se contrasta contra el rango construido con los errores ANTERIORES del mismo paso; con n errores se obtienen n-2 contrastes & Christoffersen (1998), International Economic Review 39(4), 841-862, para la evaluacion de una SECUENCIA de intervalos por su sucesion de aciertos y su cobertura incondicional; Tashman (2000) y Hyndman y Athanasopoulos (2021) 5.10 para la evaluacion de origen rodante \\
Ancho relativo de IC95 por horizonte & operativo interno, sin fuente externa & NUEVE cortes: 0.10 operativo; 0.20 corto; 0.30 medio; 0.40 largo; 0.50 extendido cercano; 0.55 extendido; 0.75 exploratorio; mas 0.45 y 0.65 escritos en linea sin constante propia & Politica interna de comunicacion de incertidumbre. Ninguna fuente respalda estos valores \\
Clasificación del intervalo del 95 \% por cobertura empírica & operativo interno, sin fuente externa & ningun corte de cobertura degrada el horizonte; solo degrada que la cobertura no sea calculable & Cobertura observada por origen movil (D-12b-C); Hyndman y Koehler (2006) \\
Seleccion del modelo por RMSE fuera de muestra global & bibliográfico & m* = argmin RMSE\_\allowbreak{}global(m) sobre la muestra comun de pares (objetivo, horizonte) & RMSE: Hyndman y Koehler (2006); FPP3 5.8 (medidas dependientes de escala son las adecuadas dentro de una misma serie). Validacion de origen movil: Tashman (2000); FPP3 5.10. La muestra de evaluacion es el rango entregable h=1..h\_\allowbreak{}max, deducido del contrato de un modelo por trayectoria \\
Continuidad de la evaluacion pese a un horizonte no recomendable (parada retirada) & derivación matemática & el bucle evalua TODOS los horizontes de la rejilla; un horizonte no recomendable no detiene la evaluacion de los siguientes (continue, no break) & Consecuencia tecnica: cada horizonte tiene su propia evidencia fuera de muestra, independiente de la de los demas \\
Horizonte maximo publicado: RETIRADO el prefijo consecutivo & operativo técnico & el maximo publicado es el mayor horizonte que cumple con su propia evidencia; los horizontes validos NO tienen que formar un prefijo & Ninguna fuente exige continuidad; cada horizonte se evalua con su propia muestra de errores \\
Seleccion del horizonte y del modelo entregados & operativo técnico & mayor horizonte permitido menor o igual al solicitado & Regla de entrega; no es un criterio de calidad \\
Salvaguarda por benchmark, retirada como decisora & operativo interno, sin fuente externa & diagnostico: evalua drift y naive y publica el resultado; no sustituye & Decision operativa D-3; el criterio de aceptacion no tenia fuente \\
Horizontes operativos de interfaz & sin efecto decisorio & 1, 3, 6, 12, 18 & Decisión de UX institucional \\
Límite de búsqueda del horizonte & operativo técnico & HORIZONTE\_\allowbreak{}MAXIMO\_\allowbreak{}AUDITORIA = 30 meses, constante & Politica computacional interna \\
Niveles de evidencia fuera de muestra por horizonte & operativo interno, sin fuente externa & 0 ventanas (h $>$ n - N0): no evaluable, por inexistencia del error fuera de muestra; 1 o mas: se entrega, con el numero de ventanas declarado & Derivacion aritmetica: evaluar h con ventana expansiva exige n - N0 - h + 1 $\geq$ 1, luego h $\leq$ n - N0. Por encima no existe ningun error fuera de muestra que medir \\
Parametros de Holt estimados por minimizacion del SSE & bibliográfico & alpha, beta*, l0 y b0 estimados; phi estimado en [0.80, 0.98] en la variante amortiguada; restriccion 0 $<$ beta* $\leq$ alpha & Hyndman y Athanasopoulos (2021), FPP3, secciones 8.1-8.2; Gardner (2006), Exponential smoothing: the state of the art - Part II \\
Elegibilidad de candidatos por estimabilidad (activación por niveles retirada) & derivación matemática & Compiten los 10 candidatos del catalogo; un modelo se excluye solo si n $<$ N\_\allowbreak{}min(m), derivado de su propia formulacion & N\_\allowbreak{}min se deriva por familia, no por una regla unica: forma cerrada (naive, drift) exige solo que existan las observaciones que la formula usa; ajuste por minimizacion (OLS, Huber, Holt) exige n $\geq$ k + 1 contando TODO parametro estimado -en Huber tambien la escala; en Holt tambien alpha, beta*, phi y los dos estados iniciales-; y los modelos sobre variaciones exigen ademas un minimo de variaciones finitas \\
Comparación contra Naive/Drift: constantes retiradas de la logica decisoria & sin efecto decisorio & ningun corte vigente; la lectura viva es C-BEN-002 frente a 1 & Sin fuente para 1,10 ni para 1,25; el inventario RC3 los declaro margen de gracia sin fuente \\
Error relativo frente al benchmark, lectura descriptiva & bibliográfico & no es\_\allowbreak{}benchmark y rRMSE frente a naive $>$ 1; se informa, no decide & Hyndman y Koehler (2006): el error relativo se define frente a UN metodo de referencia y 1 es su punto de equivalencia \\
Detección del salto de cambio de anio & operativo interno, sin fuente externa & $\geq$ 2 transiciones; ratio $>$ 1.5; signo $\geq$ 0.6 & Estadística robusta (medianas); calibrado con el anexo ICOCIV de mayo de 2026 \\
Tratamiento del salto de cambio de año: RETIRADO, no se aplica & operativo interno, sin fuente externa & aplicado = False de forma incondicional; y\_\allowbreak{}ajustado = y\_\allowbreak{}futuro & P0-F, 12-08-2026: ningun componente del tratamiento tiene sustento completo (gamma como estimador del salto futuro, la forma del factor, las puertas de activacion calibradas sobre el propio anexo y el conjunto de validacion) \\
Separacion entre patron detectado y efecto en el horizonte & operativo interno, sin fuente externa & patron\_\allowbreak{}detectado\_\allowbreak{}en\_\allowbreak{}serie, horizonte\_\allowbreak{}cruza\_\allowbreak{}cambio\_\allowbreak{}anio y efecto\_\allowbreak{}en\_\allowbreak{}horizonte\_\allowbreak{}solicitado se informan por separado & Politica interna autorizada el 29-jul-2026 (hallazgo H-10); redaccion completada el 14-08-2026 \\
Estabilidad del error entre ventanas (coeficiente de variación) & operativo interno, sin fuente externa & inestable $>$ 1.0 & Coeficiente de variación de |errores| OOS; parámetro interno \\
\end{longtable}
}

\end{landscape}

El proyecto conserva, además, un anexo técnico independiente y compilable de criterios estadísticos que reúne en un solo documento cada criterio con su fórmula numerada, su fuente verificable y su prueba de respaldo, pensado para estudiarse por separado. Ese anexo incorpora la misma Tabla~\ref{tab:matriz_criterios}, generada desde el mismo módulo interno, por lo que ambos documentos permanecen sincronizados con la aplicación.

\subsection{Preparación, validación y reproducibilidad como fases metodológicas}

La tabla \ref{tab:fases_estadisticas_metodologia} conecta los datos de entrada con la decisión permitida. Un fallo crítico de periodos o valores no produce una proyección silenciosa; una advertencia de incertidumbre tampoco altera el dato puntual, pero sí cambia su clasificación.

\begin{table}[H]
\centering
\caption{Fases metodológicas del procesamiento estadístico}
\label{tab:fases_estadisticas_metodologia}
\scriptsize
\begin{adjustbox}{max width=\textwidth}
\begin{tabularx}{\textwidth}{@{}p{2.7cm}YYYY@{}}
\toprule
\textbf{Fase} & \textbf{Entrada} & \textbf{Acción} & \textbf{Producto verificable} & \textbf{Decisión} \\
\midrule
Preparación & Fila y periodos del anexo. & Normalizar, ordenar y convertir valores. & Serie con periodo real y metadatos. & Continuar o informar formato inválido. \\
Validación & Serie normalizada. & Revisar continuidad, duplicados, faltantes, dominio y longitud. & Diagnóstico y lista de advertencias/errores. & Modelar, limitar o bloquear. \\
Modelado & Serie válida y horizonte. & Ajustar candidatos progresivos. & Parámetros y pronósticos por modelo. & Mantener candidatos ajustables. \\
Validación temporal & Candidatos y ventanas históricas. & Reentrenar, pronosticar y comparar por origen. & Errores y métricas por \(h\). & Seleccionar o descartar. \\
Incertidumbre & Pronóstico y errores fuera de muestra. & Medir el error fuera de muestra por horizonte. & Métricas agregadas (RMSE, MAE, MASE) con su número de ventanas. & Informar el alcance de la evidencia; el intervalo no se publica. \\
Reproducibilidad & Datos, parámetros y resultado. & Serializar ruta, modelo, periodos y métricas. & Interfaz, DOCX y CSV/Excel. & Permitir revisión independiente. \\
\bottomrule
\end{tabularx}
\end{adjustbox}
\end{table}

\subsection{Estudio de las necesidades del sistema}

En esta etapa se revisó el problema operativo asociado al análisis de variaciones y proyecciones de precios en obras civiles. La revisión incluyó documentos oficiales del \DANE, lineamientos de transición \ICCP--\ICOCIV, anexos técnicos disponibles en la documentación disponible del proceso, el Anexo 10 \ICCP{} histórico, anexos \ICOCIV{}, documentos de soporte de revisión de precios y la herramienta previa construida en Microsoft Excel. Con esa revisión se identificaron necesidades de consulta jerárquica, control de periodos, reducción de operaciones manuales, cálculo trazable de índices y generación de salidas revisables.

\subsubsection{Documentos y archivos revisados}

Para el desarrollo del trabajo se revisaron documentos técnicos, normativos y de datos. La tabla \ref{tab:documentos_necesidades} resume cada fuente consultada, su naturaleza y el aporte técnico que tuvo dentro de la definición del sistema. No se incluyen fuentes no disponibles.

\begin{table}[H]
\centering
\caption{Documentos y archivos revisados para el estudio de necesidades}
\label{tab:documentos_necesidades}
\scriptsize
\begin{adjustbox}{max width=\textwidth}
\begin{tabularx}{\textwidth}{@{}p{3.6cm}YY@{}}
\toprule
\textbf{Fuente revisada} & \textbf{Información aportada} & \textbf{Necesidad identificada} \\
\midrule
Documentos oficiales del \DANE{} sobre \ICOCIV{} & Estructura metodológica del índice, publicación mensual, niveles de desagregación y uso de datos oficiales. & Cargar anexos, conservar periodos y respetar la clasificación jerárquica de las series. \\
Lineamientos de empalme \ICCP--\ICOCIV{} & Fórmulas de transición, variables \(I\), \(I_0\), \(R_1\), \(R_2\), \(R\), periodo de corte y tratamiento del acero. & Implementar cálculo general y cálculo especial con trazabilidad explícita. \\
Anexo 10 \ICCP{} histórico & Series del \ICCP{} total, canasta general y grupos de obra. & Separar las listas \ICCP{} y evitar mezclar categorías en el selector. \\
Anexos \ICOCIV{} cargables & Tablas con grupos de obra, subclases, tipologías, capítulos constructivos, grupos de costos, insumos y periodos mensuales. & Diseñar el selector jerárquico y la construcción de series históricas. \\
Anexo técnico y documentos contractuales disponibles & Estructura de insumos, actividades, unidades, APU, cronogramas y soportes de solicitud de ajuste. & Relacionar insumo contractual, unidad, precio base, fechas y equivalencias técnicas. \\
Herramienta Microsoft Excel \textit{5. AJUSTE DE PRECIOS.xlsx} & Hojas de listas, índices, fórmula de ajuste, proyección, cronograma, meses de acta y resumen. & Automatizar parcialmente un flujo ya conocido, reduciendo errores de rangos, fórmulas y actualización manual. \\
\bottomrule
\end{tabularx}
\end{adjustbox}
\end{table}

\subsubsection{Flujo de trabajo identificado}

El flujo operativo identificado parte de una solicitud o necesidad de análisis, continúa con la revisión contractual y termina con una salida revisable. En términos funcionales, la entidad requiere identificar el insumo o actividad, revisar su unidad y valor base, seleccionar una serie \ICCP{} o ruta \ICOCIV{} pertinente, ubicar fechas iniciales y finales, consultar índices, aplicar la fórmula correspondiente y dejar evidencia del resultado.

\begin{figure}[H]
\centering
\fbox{\begin{minipage}{0.92\textwidth}
\centering
Solicitud de análisis \(\rightarrow\) revisión contractual \(\rightarrow\) identificación de insumo/APU \(\rightarrow\) selección de índice o equivalencia \(\rightarrow\) consulta de índices históricos \(\rightarrow\) definición de fechas \(\rightarrow\) cálculo de variación o ajuste \(\rightarrow\) revisión técnica \(\rightarrow\) reporte o exportable.
\end{minipage}}
\caption{Flujo de trabajo identificado para análisis de variaciones de precios}
\label{fig:flujo_trabajo_entidad}
\end{figure}

\subsubsection{Herramienta existente en Microsoft Excel}

La herramienta existente en Microsoft Excel permitió reconocer la lógica operativa utilizada para organizar índices, aplicar fórmulas de ajuste y apoyar la estimación de valores futuros. El archivo revisado contiene hojas como \textit{LISTAS}, \textit{ÍNDICES - ICCP}, \textit{ÍNDICES - ICOCIV}, \textit{FORMULA DE AJUSTE - CVP}, \textit{PROYECCIÓN DE ÍNDICES}, \textit{CRONOGRAMA}, \textit{CRONOGRAMA INVERSIÓN}, doce hojas mensuales de acta y una hoja \textit{RESUMEN}. Esta estructura muestra un flujo de trabajo amplio: listas para equivalencias, bases históricas de índices, hoja de empalme, proyección de índices, distribución temporal del contrato, cálculo mensual y consolidación de resultados.

En la hoja \textit{LISTAS} se observaron categorías \ICCP{} como Total \ICCP, Equipos, Materiales, Transporte, Mano de obra, Costos indirectos, Obras de explanación, Sub bases y bases, Transporte de materiales, Aceros y elementos metálicos, Acero estructural y cables de acero, Concretos, morteros y obras varias, Concreto para estructura de puentes y Pavimentaciones con asfalto. En las hojas de índices, el \ICCP{} se organiza con descripciones por fila y periodos mensuales desde enero de 2000, mientras el \ICOCIV{} se organiza por ruta jerárquica y periodos desde enero de 2021. La hoja \textit{PROYECCIÓN DE ÍNDICES} usa concatenaciones de niveles \ICOCIV{} y búsquedas con \texttt{INDEX}, \texttt{MATCH}, \texttt{OFFSET} e \texttt{IFERROR} para traer valores por periodo; además contiene gráficos para apoyar la lectura de tendencias.

Las hojas mensuales reflejan el cálculo operativo por acta. En ellas se identifican columnas de cantidad, valor unitario, valor total, fechas de inicio y fin, inversión diaria, índices \ICCP{} \(I_0\) e \(I\), índices \ICOCIV{} \(I_0\) e \(I\), \(R_1\), \(R_2\) y valor total de ajuste. Las fórmulas observadas aplican el corte de diciembre de 2021, consultan índices desde las hojas históricas y calculan \(R_1\), \(R_2\) y la suma del ajuste mensual. La hoja \textit{RESUMEN} consolida valores por mes mediante referencias a cada hoja mensual, acumulados y porcentajes.

La revisión de esta herramienta no implica una descalificación del archivo; por el contrario, sirvió para comprender el procedimiento técnico real. También evidenció riesgos propios de un flujo manual: actualización parcial de índices, selección de rangos, dependencia de fórmulas arrastradas, lectura de parámetros desde gráficos, referencias indirectas entre hojas, cambios inadvertidos en celdas y dificultad para reconstruir una decisión de equivalencia después de varios cálculos.

\begin{table}[H]
\centering
\caption{Comparación entre la herramienta previa en Excel y la aplicación desarrollada}
\label{tab:excel_vs_app}
\scriptsize
\begin{adjustbox}{max width=\textwidth}
\begin{tabularx}{\textwidth}{@{}p{3.1cm}YYY@{}}
\toprule
\textbf{Aspecto} & \textbf{Herramienta previa en Excel} & \textbf{Aplicación desarrollada} & \textbf{Mejora esperada} \\
\midrule
Actualización de índices \ICOCIV{} & Requería actualización manual de índices publicados por el \DANE. & Organiza la información cargada desde anexos oficiales y controla periodos. & Reduce errores de digitación y mejora trazabilidad. \\
Organización de índices & Separaba información en pestañas y rangos. & Centraliza carga, lectura y rutas de datos. & Facilita reproducibilidad. \\
Selección de rutas jerárquicas & Dependía de selección manual de tablas o rangos. & Usa selector jerárquico \ICOCIV. & Disminuye ambigüedad en la ruta usada. \\
Obtención de ecuaciones & Se apoyaba en gráficas de Excel y parámetros copiados. & Evalúa modelos y registra métricas. & Mejora control del procedimiento. \\
Cálculo de proyecciones & Podía requerir arrastre de fórmulas. & Ejecuta proyección desde servicios de análisis. & Reduce manipulación manual. \\
Trazabilidad de \(I\) e \(I_0\) & Dependía de celdas y rangos. & Registra valores, fechas y ruta seleccionada. & Facilita revisión técnica. \\
Empalme \ICCP--\ICOCIV{} & Dependía de fórmulas manuales. & Implementa fórmulas con variables explícitas. & Evita pérdida de evidencia del cálculo. \\
Cálculo especial para acero & Requería control manual de variables. & Diferencia cálculo general y acero. & Mejora claridad metodológica. \\
Presentación de resultados & Quedaba distribuida entre hojas. & Presenta tablas acumulativas y detalle técnico. & Mejora lectura del resultado. \\
Exportación de resultados & Dependía del archivo de trabajo. & Exporta resultados y trazabilidad. & Facilita archivo y revisión. \\
Riesgo de error manual & Asociado a digitación, rangos y fórmulas. & Usa validaciones de entrada. & Reduce errores operativos. \\
Reproducibilidad & Dependía del estado de la hoja. & Conserva entradas, rutas, índices y fórmulas. & Permite repetir y auditar el cálculo. \\
\bottomrule
\end{tabularx}
\end{adjustbox}
\end{table}

\subsubsection{Resultado de la fase: requerimientos y reglas}

Como resultado de esta etapa se consolidaron requerimientos funcionales asociados a carga de datos, consulta, filtrado, visualización, proyección, empalme, reportes y exportación. La tabla \ref{tab:req_funcionales} sintetiza los requerimientos principales. Cada requerimiento conserva un código (RF para los funcionales y RNF para los no funcionales) porque esos identificadores se utilizan más adelante para trazar la relación entre cada requerimiento y las pruebas que lo verifican; por ejemplo, el requerimiento RF-10 se contrasta en el capítulo de pruebas con el caso de exportación correspondiente.

\begin{table}[H]
\centering
\caption{Requerimientos funcionales del sistema}
\label{tab:req_funcionales}
\scriptsize
\begin{adjustbox}{max width=\textwidth}
\begin{tabularx}{\textwidth}{@{}p{1.4cm}p{4.1cm}YY@{}}
\toprule
\textbf{Código} & \textbf{Requerimiento funcional} & \textbf{Descripción} & \textbf{Módulo asociado} \\
\midrule
RF-01 & Cargar archivos \ICOCIV{} & Leer anexos oficiales, detectar periodos y conservar filas fuente. & Carga de datos \\
RF-02 & Seleccionar rutas jerárquicas & Permitir escoger niveles \ICOCIV{} y guardar la ruta completa. & Selector \ICOCIV{} \\
RF-03 & Consultar serie histórica & Mostrar valores mensuales asociados a la ruta seleccionada. & Serie histórica \\
RF-04 & Ejecutar análisis y proyección & Evaluar modelos, métricas y advertencias. & Proyección \ICOCIV{} \\
RF-05 & Generar informe DOCX & Exportar ruta, serie, modelo, métricas, gráfica y metodología. & Reportes \\
RF-06 & Seleccionar serie \ICCP{} & Diferenciar Total \ICCP, Canasta general y Grupo de obra. & Empalme \ICCP--\ICOCIV{} \\
RF-07 & Calcular empalme general & Aplicar \(R_1\), \(R_2\), \(R=R_1+R_2\) y valor actualizado. & Empalme \\
RF-08 & Calcular caso especial de acero & Aplicar las variables y comparación propias del acero. & Empalme acero \\
RF-09 & Generar tablas acumulativas & Registrar equivalencias y tabla de valor ajustado por cálculo válido. & Resultados \\
RF-10 & Exportar Excel & Entregar una hoja con resultados, fórmulas y trazabilidad. & Exportación \\
\bottomrule
\end{tabularx}
\end{adjustbox}
\end{table}

\begin{table}[H]
\centering
\caption{Requerimientos no funcionales del sistema}
\label{tab:req_no_funcionales}
\scriptsize
\begin{adjustbox}{max width=\textwidth}
\begin{tabularx}{\textwidth}{@{}p{1.4cm}p{3.7cm}YY@{}}
\toprule
\textbf{Código} & \textbf{Requerimiento no funcional} & \textbf{Descripción} & \textbf{Criterio de cumplimiento} \\
\midrule
RNF-01 & Usabilidad & Interfaz clara para usuarios técnicos no programadores. & Formularios, botones y mensajes comprensibles. \\
RNF-02 & Trazabilidad & Conservación de fechas, rutas, índices, fórmulas y resultados. & Detalle técnico y exportables revisables. \\
RNF-03 & Reproducibilidad & Posibilidad de reconstruir el cálculo. & Entradas y salidas persistidas en tablas. \\
RNF-04 & Validación de errores & Bloqueo de cálculos incompletos o inconsistentes. & Mensajes explícitos y ausencia de fallos silenciosos. \\
RNF-05 & Modularidad & Separación entre interfaz, datos, estadística, empalme y reportes. & Cambios locales sin romper módulos no relacionados. \\
RNF-06 & Reducción de manipulación manual & Disminución de selección de rangos y arrastre de fórmulas. & Cálculos ejecutados por servicios internos. \\
\bottomrule
\end{tabularx}
\end{adjustbox}
\end{table}

Las reglas básicas del proceso quedaron definidas así: toda actualización requiere fecha inicial y fecha final; los índices \(I\) e \(I_0\) deben mostrarse explícitamente; la equivalencia \ICCP--\ICOCIV{} requiere criterio técnico; el \ICCP{} se usa como fuente histórica hasta el corte de transición cuando corresponda; el \ICOCIV{} se usa para el tramo posterior y para proyección si el horizonte es viable; el cálculo general y el cálculo del acero usan fórmulas diferentes; el resultado debe conservar trazabilidad; y no se calcula cuando faltan índices, unidad, serie \ICCP{} o ruta \ICOCIV{} obligatoria.

\subsection{Diseño de la plataforma de software}

El diseño de la plataforma se estructuró en módulos para evitar que la lógica de interfaz, datos, estadística, empalme y reportes quedara mezclada. Se definió una aplicación con una ventana principal, paneles de selección, pestañas de resultados, tablas acumulativas y servicios internos. La arquitectura separó la carga de anexos \ICOCIV, el análisis de series, la validación estadística, la proyección, el empalme \ICCP--\ICOCIV{} y la generación de exportables.

\begin{figure}[H]
\centering
\fbox{\begin{minipage}{0.92\textwidth}
\centering
Usuario \(\rightarrow\) capa de presentación \(\rightarrow\) servicios de procesamiento \(\rightarrow\) datos \ICOCIV/\ICCP{} \(\rightarrow\) validación, proyección y empalme \(\rightarrow\) reportes DOCX, tablas y Excel.
\end{minipage}}
\caption{Arquitectura conceptual de la aplicación}
\label{fig:arquitectura_conceptual_app}
\end{figure}

Durante esta etapa se definió que el selector jerárquico \ICOCIV{} sería el punto de entrada para construir la ruta técnica de una serie. También se estableció que el Anexo 10 \ICCP{} histórico debía cargarse internamente para evitar que el usuario repitiera esa operación en cada cálculo. La interfaz se organizó para que el usuario pudiera distinguir el módulo de análisis/proyección \ICOCIV{} del módulo de actualización de precios o empalme, conservando una ruta de información compartida cuando el empalme necesitara un índice \ICOCIV{} real o proyectado.

\begin{table}[H]
\centering
\caption{Módulos principales definidos en el diseño de la plataforma}
\label{tab:modulos_diseno}
\scriptsize
\begin{adjustbox}{max width=\textwidth}
\begin{tabularx}{\textwidth}{@{}p{3.2cm}YYYY@{}}
\toprule
\textbf{Módulo} & \textbf{Propósito} & \textbf{Entradas} & \textbf{Procesamiento} & \textbf{Salidas} \\
\midrule
Carga de datos \ICOCIV{} & Leer anexos oficiales. & Archivo seleccionado por el usuario. & Detección de tablas, periodos y filas. & Datos normalizados. \\
Selector jerárquico \ICOCIV{} & Escoger una ruta técnica. & Niveles disponibles. & Filtrado dependiente por nivel. & Ruta completa y fila fuente. \\
Serie histórica & Mostrar datos mensuales. & Ruta seleccionada. & Ordenamiento temporal. & Tabla y gráfica de serie. \\
Análisis y proyección & Estimar valores futuros viables. & Serie histórica y horizonte. & Validación, modelos, métricas y advertencias. & Resultado puntual, métricas fuera de muestra y diagnóstico. \\
Reportes DOCX & Documentar resultados. & Resultado estadístico. & Composición de tablas, gráfica y metodología. & Informe técnico. \\
Empalme \ICCP--\ICOCIV{} & Calcular actualización de precios. & Fechas, serie \ICCP, ruta \ICOCIV, valores contractuales. & Búsqueda de índices y fórmulas \(R_1\), \(R_2\), \(R\). & Resultado y tablas acumulativas. \\
Cálculo especial de acero & Aplicar tratamiento diferenciado. & \(P_0\), \(I_x\), \(q\), índices y fechas. & Fórmulas específicas del acero. & Ajuste y comparación técnica. \\
Exportación Excel & Conservar trazabilidad externa. & Historial de cálculos. & Organización de resultados y fórmulas. & Libro de resultados en una sola hoja, verificado por prueba automática. \\
Validaciones y mensajes & Evitar cálculos incompletos. & Entradas de formularios. & Reglas de obligatoriedad y consistencia. & Mensajes claros al usuario. \\
\bottomrule
\end{tabularx}
\end{adjustbox}
\end{table}

El flujo de navegación previsto inicia con la carga del archivo \ICOCIV{}, continúa con la selección de una ruta jerárquica, consulta la serie histórica, ejecuta análisis o proyección, revisa resultados y genera informe si corresponde. Para el empalme, el usuario ingresa al módulo de actualización de precios, selecciona fechas, tipo de serie \ICCP, serie \ICCP, unidad, insumo contractual y ruta \ICOCIV{} equivalente; luego calcula, revisa las tablas acumulativas y exporta los resultados.

\subsection{Desarrollo e implementación del software}

La implementación se realizó mediante módulos de Python organizados por responsabilidad. El cargador de datos quedó encargado de leer anexos oficiales, detectar encabezados y periodos mensuales, extraer tablas y conservar la relación entre fila fuente y ruta jerárquica. El módulo de proyección construyó la serie seleccionada, ejecutó validaciones, evaluó modelos candidatos, calculó métricas y entregó resultados con advertencias y con el estado metodológico de cada horizonte.

La tabla \ref{tab:bibliotecas_reales} resume las bibliotecas identificadas en \texttt{requirements.txt} y en imports del código fuente. Esta identificación evita presentar dependencias no verificadas como parte del sistema.

\begin{table}[H]
\centering
\caption{Bibliotecas reales utilizadas por la aplicación}
\label{tab:bibliotecas_reales}
\scriptsize
\begin{adjustbox}{max width=\textwidth}
\begin{tabularx}{\textwidth}{@{}p{3.2cm}YY@{}}
\toprule
\textbf{Biblioteca} & \textbf{Uso en el proyecto} & \textbf{Módulo asociado} \\
\midrule
\texttt{pandas} & Lectura, limpieza y presentación de tablas de índices. & Datos, interfaz, reportes y proyección. \\
\texttt{numpy} & Cálculos numéricos, métricas y manejo de series. & Estadística, validación y utilidades. \\
\texttt{scikit-learn} & Regresión lineal y robusta. & Modelos y estadística. \\
\texttt{scipy} & Cuantiles de la distribución t, empleados por el cálculo interno de la banda que no se publica. & Diagnóstico interno. \\
\texttt{pyxlsb} & Soporte para lectura de libros binarios cuando aplica. & Carga de datos. \\
\texttt{openpyxl} & Lectura/escritura de archivos Excel y exportación de empalme. & Empalme y exportables. \\
\texttt{python-docx} & Generación de informes técnicos en formato DOCX. & Reportes. \\
\texttt{matplotlib} & Gráficas de series y salidas visuales de reportes. & Interfaz y reportes. \\
\texttt{PySide6} & Construcción de la interfaz gráfica. & Ventana principal, widgets y controladores. \\
\texttt{statsmodels} & Dependencia obligatoria, versión fija. Uso exclusivo del diagnóstico Ljung-Box. & Diagnóstico residual. No interviene en modelos ni en intervalos. \\
\bottomrule
\end{tabularx}
\end{adjustbox}
\end{table}

El módulo de empalme \ICCP--\ICOCIV{} incorporó la separación entre Total \ICCP, Canasta general y Grupo de obra, el uso del periodo de transición de diciembre de 2021, la búsqueda de índices \ICCP{} e \ICOCIV, el cálculo general de transición y el tratamiento especial para acero. La interfaz agregó controles para tipo de serie \ICCP, serie \ICCP, unidad, insumo contractual, valores base, fechas, ruta \ICOCIV{} equivalente y botones de cálculo, limpieza, detalle, edición, eliminación y exportación.

Como ejemplo de funcionamiento, una ruta \ICOCIV{} seleccionada se transforma en una serie mensual, la serie pasa por validaciones de continuidad y longitud, los modelos candidatos generan pronósticos comparables y el servicio de proyección devuelve un resultado con modelo, horizonte, métricas fuera de muestra, advertencias y estado metodológico. En el empalme, el insumo contractual se vincula con una serie \ICCP{} y una ruta \ICOCIV{} equivalente; con las fechas seleccionadas se consultan \(I_0\) e \(I\), se calculan \(R_1\), \(R_2\), \(R\), el valor actualizado y se agrega una fila a las tablas acumulativas.

La implementación también contempló salidas documentales. El reporte de proyección se organizó para incluir ruta seleccionada, datos fuente, diagnóstico de la serie, modelo seleccionado, métricas, horizonte, resultados y metodología. El exportable de empalme se orientó a conservar equivalencias, índices, fechas, fórmulas, variables y trazabilidad del cálculo.

\subsection{Análisis estadístico y validación de resultados}

El análisis estadístico se aplicó sobre series históricas mensuales del \ICOCIV. Antes de proyectar, la aplicación revisó continuidad, longitud, valores disponibles y condiciones mínimas de modelado. Luego evaluó modelos interpretables y modelos condicionales de acuerdo con la información disponible. Entre las alternativas implementadas se incluyeron modelos de referencia como Naive y Drift, suavizamiento exponencial de Holt y Holt amortiguado, transformaciones logarítmicas y modelos robustos. Las decisiones de alcance sobre modelos considerados y no incorporados se documentan en la sección \ref{sec:decisiones_alcance_modelos}.

La validación de resultados se apoyó en backtesting temporal y métricas como MAE, RMSE, MAPE, sMAPE y MASE. Estas métricas permitieron comparar candidatos, justificar el modelo seleccionado y advertir al usuario cuando el horizonte solicitado no era recomendable. El sistema no presentó la proyección como un dato oficial del \DANE; la presentó como una estimación técnica condicionada por la calidad, longitud y comportamiento de la serie analizada.

La decisión metodológica fue evaluar cada horizonte como un problema propio. Por esta razón, los errores de \(h=1\) no se mezclan con los de \(h=6\) o \(h=12\). El flujo adoptado fue: preparar la serie mensual, validar datos, ejecutar modelos candidatos, repetir validación walk-forward, calcular métricas por horizonte, revisar residuos, clasificar el horizonte y declarar la evidencia fuera de muestra que lo sostiene, y solo entonces presentar la proyección. Esta secuencia conecta directamente la metodología estadística con la implementación de la aplicación.

En el empalme \ICCP--\ICOCIV, la validación se concentró en verificar fechas, disponibilidad de índices, correspondencia entre tipo de serie \ICCP{} y serie seleccionada, selección de unidad, ruta \ICOCIV{} equivalente y aplicación de las fórmulas de los lineamientos revisados. Para los cálculos de transición se conservaron las variables \(I\), \(I_0\), \(R_1\), \(R_2\), \(R\), base ajustable y valor actualizado, con el fin de que el resultado pudiera reconstruirse posteriormente.

\subsection{Integración y validación funcional del sistema}

La integración conectó la ventana principal, los controladores de interfaz, los servicios de datos, el servicio de proyección, el servicio de empalme y los generadores de reportes. En el flujo \ICOCIV, la selección jerárquica alimentó la construcción de la serie histórica, la vista de datos, la gráfica, el análisis estadístico y la proyección. En el flujo de empalme, la misma ruta \ICOCIV{} sirvió para obtener el índice final o para activar la proyección existente cuando la fecha solicitada superaba el último periodo real disponible.

La validación funcional revisó escenarios de carga, selección de rutas, cambio de tipo de serie \ICCP, limpieza de campos, cálculo con fechas antes y después del periodo de transición, acumulación de filas en tablas, visualización del detalle técnico y exportación. También se verificaron mensajes de error para evitar fallos silenciosos cuando faltaba una unidad, una serie, una ruta equivalente o un índice disponible. Esta etapa permitió corregir comportamientos de interfaz, listas \ICCP, controles numéricos de fechas y formatos de resultados.

\subsection{Documentación y manual de usuario}

La documentación se construyó a partir de los documentos oficiales revisados, los módulos reales de la aplicación, los anexos técnicos y las verificaciones realizadas durante el desarrollo. El presente documento técnico dejó organizados el marco conceptual, la metodología, el desarrollo del sistema, las pruebas, el manual de usuario, las fórmulas y las referencias. El manual incorporó procedimientos para cargar archivos, seleccionar rutas \ICOCIV, ejecutar análisis y proyecciones, interpretar resultados, usar el módulo de empalme, calcular el caso general, calcular el caso especial de acero y exportar resultados.

La documentación evita presentar el software como una herramienta de decisión automática. En su lugar, lo describe como un apoyo técnico que ordena información oficial, aplica procedimientos programados, presenta advertencias y deja evidencia para revisión profesional. Esta precisión fue necesaria porque las equivalencias entre insumos contractuales, series \ICCP{} y rutas \ICOCIV{} requieren criterio técnico del usuario responsable.

\subsection{Puesta en conocimiento a la Secretaría de Infraestructura de la Gobernación del Cauca}

Esta etapa quedó preparada mediante la consolidación del software, el documento técnico, el manual de usuario y los exportables que explican la operación de la herramienta. El material producido permite presentar a la Secretaría de Infraestructura de la Gobernación del Cauca el alcance del sistema, los datos que utiliza, las fórmulas aplicadas, las validaciones incorporadas, las limitaciones de las proyecciones y el papel del profesional responsable en la interpretación de resultados.

La socialización institucional deberá apoyarse en ejemplos de uso con información histórica representativa, mostrando la carga de anexos, la selección jerárquica, la proyección de una serie, el cálculo de empalme \ICCP--\ICOCIV{} y la exportación de evidencias. De esta manera, la entrega no se limitará a mostrar una interfaz, sino que permitirá explicar la trazabilidad completa entre datos de entrada, procesamiento, resultados y revisión técnica.

\section{DESARROLLO DEL SOFTWARE O DEL SISTEMA}

\subsection{Descripción general del sistema}

La aplicación desarrollada, SAVIP (Sistema de Análisis de Variaciones de Precios), integra consulta de información oficial, análisis de series históricas, proyección de índices, empalme \ICCP--\ICOCIV, generación de reportes y exportación de resultados. Su diseño responde a una necesidad práctica: transformar un flujo originalmente dependiente de hojas de cálculo, selección manual de rangos y fórmulas dispersas en una aplicación con entradas controladas, cálculos reproducibles y salidas trazables.

El sistema trabaja con dos familias de información. La primera corresponde al \ICOCIV, cargado desde anexos oficiales seleccionados por el usuario. La segunda corresponde al \ICCP{} histórico, integrado internamente a partir del Anexo 10 para soportar el tramo anterior o igual a diciembre de 2021. La figura \ref{fig:flujo_sistema} resume el flujo principal.

\begin{figure}[H]
\centering
\fbox{\begin{minipage}{0.9\textwidth}
\centering
Carga de datos oficiales \(\rightarrow\) selección jerárquica \(\rightarrow\) construcción de serie histórica \(\rightarrow\) análisis y proyección \(\rightarrow\) empalme \ICCP--\ICOCIV{} \(\rightarrow\) tablas, reportes y exportables.
\end{minipage}}
\caption{Flujo general de la aplicación desarrollada}
\label{fig:flujo_sistema}
\end{figure}

\subsection{Arquitectura de la aplicación}

La arquitectura se organizó por capas y responsabilidades. Esta decisión permitió corregir componentes puntuales sin romper otros módulos, por ejemplo ajustar el empalme \ICCP--\ICOCIV{} sin afectar el análisis/proyección \ICOCIV{} ya existente.

\subsubsection{Capa de presentación}

La capa de presentación contiene la ventana principal, los paneles de selección, las pestañas de resultados, las tablas, los formularios, los botones y los mensajes de validación. En esta capa se ubican componentes como la ventana principal, el selector jerárquico, los widgets de resultados, el panel de empalme y los diálogos de exportación. Su función es recibir entradas del usuario, mostrar resultados y comunicar errores de forma comprensible.

\subsubsection{Capa de procesamiento}

La capa de procesamiento contiene los servicios que ejecutan la lógica técnica: carga de datos, construcción de series, análisis estadístico, validación, proyección, empalme de índices y generación de reportes. En esta capa se concentran funciones como la resolución de la fila seleccionada, la construcción de la serie histórica, el backtesting, la selección de modelos, el cálculo de \(R_1\), \(R_2\), \(R\) y la preparación de salidas.

\subsubsection{Capa de datos}

La capa de datos contiene cargadores de anexos, datos \ICCP{} internos, rutas de archivos, estructuras de tablas, series normalizadas y valores seleccionados por el usuario. Esta capa no decide la interpretación del resultado; suministra datos ordenados y verificables para que la capa de procesamiento pueda aplicar las reglas definidas.

\subsubsection{Capa de salidas}

La capa de salidas genera documentos y archivos de apoyo. En el módulo de proyección se producen reportes con resultados, diagnóstico, metodología, gráfica, serie histórica y trazabilidad. En el módulo de empalme se exportan tablas de equivalencias, cálculo del valor ajustado, índices usados, fórmulas y observaciones. Estas salidas fueron diseñadas para facilitar revisión técnica posterior.

\subsection{Estructura modular de la aplicación}

La aplicación se organiza en módulos con responsabilidades separadas: datos, estadística, validación, proyección, servicios, interfaz, reportes y exportables. La tabla \ref{tab:modulos_app} resume el papel de cada uno dentro del sistema.

\begin{table}[H]
\centering
\caption{Módulos principales de la aplicación}
\label{tab:modulos_app}
\scriptsize
\begin{adjustbox}{max width=\textwidth}
\begin{tabularx}{\textwidth}{@{}p{3.2cm}YY@{}}
\toprule
\textbf{Módulo} & \textbf{Responsabilidad principal} & \textbf{Componentes principales} \\
\midrule
Datos & Lectura de anexos, detección de tablas y construcción de estructuras base. & Cargador de anexos, detección de encabezados, periodos y tablas. \\
Estadística & Preparación, diagnóstico y modelado de series. & Cálculo de métricas, análisis de continuidad, modelos interpretables y diagnóstico del patrón de cambio de año (no modifica la trayectoria). \\
Validación & Evaluación temporal de pronósticos y selección comparativa. & Backtesting, walk-forward, selección de mejor modelo. \\
Proyección & Ejecución del flujo de proyección \ICOCIV. & Servicio de proyección, construcción de serie y horizonte viable. \\
Servicios & Cálculos técnicos del empalme contractual. & Servicio de empalme \ICCP--\ICOCIV. \\
Interfaz & Ventanas, widgets, controladores, formularios, tablas y botones. & Ventana principal, controlador principal y formulario de empalme. \\
Reportes & Generación de reportes técnicos y evidencias exportables. & Generador de reportes DOCX, PDF, HTML y tablas de reproducibilidad. \\
\bottomrule
\end{tabularx}
\end{adjustbox}
\end{table}

\subsection{Módulo de carga y selección ICOCIV}

\subsubsection{Carga de anexos ICOCIV}

El sistema lee archivos oficiales del \DANE{} y extrae tablas mensuales para construir las series disponibles. El cargador identifica títulos de tabla, inicio de datos, fin de tabla y columnas de periodo. Esta lectura evita que el usuario seleccione manualmente rangos de celdas y reduce el riesgo de tomar una fila equivocada.

El proceso contempla encabezados con periodos mensuales, campos fijos de clasificación y valores numéricos asociados a cada ruta. Cuando una tabla se carga correctamente, la aplicación conserva la estructura necesaria para mostrarla en la interfaz, construir una serie histórica y reconstruir la fila fuente usada en el análisis.

\subsubsection{Selector jerárquico ICOCIV}

El selector jerárquico permite escoger niveles de la estructura \ICOCIV{} y conserva la ruta completa de la fila seleccionada. Esta ruta es importante porque el \ICOCIV{} no se consulta como un dato aislado, sino como una posición dentro de una clasificación. La aplicación registra niveles como grupo, subgrupo, clase, tipología, capítulo, grupo de costo o insumo, según la tabla disponible.

La selección jerárquica alimenta tres procesos: consulta de serie histórica, análisis/proyección y empalme \ICCP--\ICOCIV. En el empalme, la ruta seleccionada se usa como insumo equivalente \ICOCIV{} y se presenta en la tabla de equivalencias.

\subsubsection{Consulta de serie histórica}

La consulta de serie histórica presenta los valores mensuales asociados a la ruta seleccionada. La serie se transforma en una secuencia temporal ordenada por periodo, con valores numéricos listos para diagnóstico, graficación y modelado. Esta vista permite al usuario comprobar que la ruta seleccionada sí contiene datos antes de ejecutar una proyección.

\subsection{Módulo de análisis y proyección}

\subsubsection{Preparación de la serie}

La serie se normaliza por periodo mensual y se valida antes de modelar. La aplicación revisa longitud disponible, continuidad, valores faltantes, valores no numéricos y comportamiento general. Esta validación previa evita ejecutar modelos sobre series sin información suficiente o con estructuras que impidan interpretar el resultado.

\subsubsection{Modelos implementados}

La aplicación implementa los diez candidatos del catálogo: Naive y Drift como referencias; las regresiones lineal, logarítmica temporal y exponencial o log-lineal; el modelo robusto de Huber; los suavizamientos de Holt lineal y Holt amortiguado; y los modelos sobre variación lineal y log-variación. Todos compiten en cada serie siempre que puedan estimarse, sin activación por niveles ni por longitud, y la trayectoria final se genera siempre con un único modelo por serie. Los criterios que delimitaron este catálogo y la hipótesis que representa cada candidato se justifican en la metodología (§\ref{sec:criterios_catalogo}); aquí se describe su implementación.

El sistema no selecciona el modelo más complejo por defecto ni pondera varios criterios entre sí. La interpretabilidad y la estabilidad justificaron la \emph{conformación} del catálogo e informan la lectura del resultado, pero \textbf{no intervienen en la elección del ganador}: esta se decide con una única función objetivo, el error cuadrático medio fuera de muestra sobre la muestra común a todos los candidatos (§\ref{sec:criterios_metricas}).

\subsubsection{Backtesting y selección de modelo}

El sistema compara modelos mediante validación temporal. En lugar de depender solamente del ajuste histórico, simula pronósticos en cortes anteriores y calcula errores fuera de muestra por horizonte.

Con esa evidencia se selecciona \textbf{un único modelo por serie} para generar toda la trayectoria. Cada observación fuera de muestra queda identificada por el par \((t,h)\): el periodo objetivo y el número de pasos con que se pronosticó. Sea \(S\) el conjunto de pares en los que \emph{todos} los candidatos tienen error finito ---la \textbf{muestra común}---. El modelo seleccionado es el que minimiza el error cuadrático medio sobre esa muestra:
\begin{equation}
RMSE_{k}^{\,glob}=\sqrt{\frac{\sum_{(t,h)\in S} e_{k,t,h}^{2}}{|S|}},
\qquad k^{*}=\arg\min_k RMSE_{k}^{\,glob}.
\end{equation}
La restricción a la muestra común impide que un modelo obtenga ventaja por disponer de menos observaciones: todos se puntúan sobre exactamente las mismas. No se imputa ningún valor faltante ni se penaliza su ausencia. Sobre el anexo de mayo de 2026 la restricción no descarta ninguna observación ---los diez candidatos comparten los 348 pares de la grilla en las diez series---, pero la garantía es estructural y no depende de ese hecho.

\paragraph{Evaluación numérica de la expresión.} La implementación del \(RMSE^{\,glob}\) fue estabilizada numéricamente mediante una forma algebraicamente equivalente basada en la norma euclídea, preservando muestra, pesos, definición y criterio de selección. El motivo es de aritmética de punto flotante y no de método: con errores individualmente finitos pero de magnitud muy grande, elevar al cuadrado antes de sumar puede desbordar el rango del flotante de doble precisión aunque el resultado final sí quepa. Una primera corrección evaluó \(\sqrt{\sum_i e_i^{2}}\) como \(\operatorname{hypot}(e_1,\dots,e_n)\), que aplica un escalado interno y evita materializar los cuadrados. Esa forma \textbf{no resultó suficiente}, y conviene decir por qué: al escalar internamente \emph{pierde resolución}, y se midió que devuelve el \emph{mismo} flotante para dos candidatos cuyos errores difieren en una unidad en el último lugar. Como el desempate elige el primer candidato del banco, ese empate artificial hacía ganar al modelo equivocado: el resultado dependía del redondeo y del orden de inserción, no del error medido. La implementación vigente compara, por tanto, la \textbf{suma exacta de cuadrados}, acumulada sobre racionales a partir de la representación binaria exacta de cada flotante. Como \(\sqrt{\cdot}\) es estrictamente creciente y \(|S|\) es común a todos los candidatos, \(\arg\min RMSE = \arg\min SSE\): la métrica sigue siendo el RMSE y la regla no cambia; lo que cambia es la aritmética con que se compara. La comprobación sobre diez escenarios ordinarios reproduce los mismos modelos seleccionados y el mismo ordenamiento de candidatos, con diferencias relativas del orden de \(10^{-16}\), es decir del último dígito representable. No se trata de una metodología nueva ni de un cambio de regla.

La regla no contiene ningún parámetro libre. Conviene, aun así, declarar su ponderación implícita: cada observación pesa igual en el denominador, pero los horizontes largos aportan errores de mayor magnitud y por tanto dominan el numerador. En una de las series evaluadas, el horizonte de dieciocho meses aporta el 2,6\,\% de las observaciones y el 8,8\,\% de la suma de cuadrados. El peso lo determina la magnitud del error medido y no una constante elegida por el analista.

\paragraph{Por qué la muestra de evaluación es ésa.} La expresión anterior no es una agregación de métricas distintas: es el error cuadrático medio calculado sobre una \emph{muestra de evaluación}. Lo que requiere fundamento, por tanto, no es «agregar», sino \emph{cuál} es esa muestra. Y la respuesta se deduce del contrato del sistema, no de una preferencia.

El sistema selecciona \textbf{un único modelo por serie} y publica \textbf{todos los meses} comprendidos entre el siguiente al último dato observado y el mes objetivo que el usuario solicite, hasta el máximo que la evidencia permita. Pedir tres, seis, dieciocho o veinticuatro meses produce, respectivamente, tres, seis, dieciocho o veinticuatro filas mensuales, cada una con su pronóstico puntual. Si algún mes intermedio no está disponible, su fila se marca como tal y no se dibuja unida a sus vecinas: no se interpola ningún valor. En consecuencia, el modelo elegido deberá producir \emph{todos} esos pronósticos, y la muestra de evaluación pertinente es exactamente el conjunto de pronósticos que tendrá que producir: todos los orígenes de la validación temporal por todos los horizontes entregables. La rejilla evaluada coincide con ese rango.

\paragraph{Por qué no se escalan los errores.} Podría objetarse que las competencias de pronóstico agregan errores \emph{escalados} y no cuadrados en unidades originales. La objeción no aplica aquí, y conviene decir por qué. El escalado resuelve un problema concreto: comparar exactitud \textbf{entre series con unidades distintas}, que es la situación de una competencia con decenas de miles de series heterogéneas. \textcite{hyndman2021fpp3} lo enuncian directamente: las medidas dependientes de escala son las apropiadas cuando se comparan métodos aplicados a una sola serie temporal, o a series con las mismas unidades, mientras que los errores escalados se propusieron para comparar exactitud \emph{entre} series de unidades diferentes. SAVIP compara modelos \textbf{dentro de una misma serie} y todos los errores están expresados en puntos del mismo índice ---incluidos los de los modelos que operan sobre transformaciones logarítmicas, porque el error se mide siempre sobre el valor retransformado a nivel---.

\paragraph{Ponderación implícita y limitación.} Dentro del error cuadrático agregado, los horizontes largos pesan más porque sus errores son mayores. Eso no es un parámetro sino el significado de minimizar la pérdida cuadrática sobre la trayectoria: una alternativa de peso igual por horizonte sería una función de pérdida \emph{distinta} y exigiría su propio fundamento. Queda además declarada una limitación: \textcite{clements_hendry1993} demuestran que la comparación de errores cuadráticos medios en contextos multihorizonte no es invariante ante representaciones isomorfas del proceso, de modo que distintas representaciones pueden producir ordenaciones distintas. Evaluar todos los candidatos en la misma métrica de nivel acota ese problema, pero no lo elimina.

Como el conjunto de horizontes evaluados depende únicamente de la longitud de la serie, el modelo seleccionado ---y con él los primeros valores proyectados--- es una propiedad de la serie y no del horizonte que solicite el usuario: pedir 3 o 18 meses produce los mismos valores para los meses comunes.

Una versión anterior de esta regla promediaba el RMSE relativo al mejor modelo de cada horizonte ponderando por \(1/h\). Esa ponderación se retiró: no se identificó fuente que la sustente y sí modificaba el modelo entregado, de modo que constituía un criterio heurístico con efecto sobre el resultado.

Debe distinguirse entre el \emph{modelo seleccionado para la serie} y el \emph{horizonte máximo recomendado o admisible}: el modelo genera una trayectoria completa, pero la clasificación por horizonte determina qué parte de esa trayectoria se recomienda, se presenta con advertencia, se degrada a escenario o se bloquea. Ambas decisiones usan el mismo modelo, de modo que las métricas y las advertencias reportadas corresponden siempre al modelo que produce los valores. Desde agosto de 2026 la clasificación por horizonte tampoco exige contigüidad: si un horizonte intermedio no se sostiene, los posteriores conservan su propia evidencia y el hueco se publica como tal.

El criterio de la Ecuación anterior es el único que decide la selección. Métricas de ajuste dentro de muestra como el AICc, el \(R^2\) ajustado o el estadístico Durbin--Watson se calculan y se reportan como diagnósticos complementarios, pero no intervienen en la elección del modelo, precisamente porque el ajuste histórico no es evidencia de capacidad predictiva \parencite{granger_newbold1974}.

\paragraph{Salvaguarda por benchmark: diagnóstico, no sustitución.} Un resultado no viable en \(h=1\) impide entregar ese horizonte con el modelo principal; como se indicó arriba, eso no cancela por sí solo horizontes posteriores que sostengan su propia evidencia. Ante un \(h=1\) no viable, la aplicación reevalúa además los benchmarks Drift y Naive con los mismos criterios y publica hasta dónde llegaría cada uno, qué horizontes no cubre el modelo principal y si un benchmark habría ampliado el alcance.

Ese resultado es \textbf{exclusivamente diagnóstico}. Hasta agosto de 2026 la salvaguarda sustituía al modelo principal en toda la trayectoria cuando un benchmark ampliaba el horizonte admisible. El criterio de sustitución era \(h_{\text{bench}}>h_{\text{previo}}\) ---más horizonte, no menos error--- y carecía de fuente que lo sustentara: sobre el anexo de mayo entregaba un Drift con error cuadrático medio peor que el del modelo descartado en los horizontes de tres, seis, doce y dieciocho meses. La sustitución se retiró. \textbf{El modelo entregado es siempre el que selecciona el desempeño fuera de muestra}, y la salvaguarda solo informa del alcance que tendría una referencia simple.

Un caso de evaluación ilustra el procedimiento y, sobre todo, delimita cuándo la salvaguarda interviene. La serie corresponde al insumo Arena del anexo del \ICOCIV{} con corte en enero de 2026, con periodicidad mensual y la misma extensión histórica que el resto de las series analizadas ($n = 61$ observaciones). Ejecutado el procedimiento sobre esa serie, la selección por RMSE fuera de muestra sobre la muestra común elige Drift con el menor error del catálogo en $h = 1$ ($\mathrm{RMSE} = 0{,}7928$, frente a $0{,}8248$ del siguiente candidato y $1{,}0395$ del Naive). Ningún horizonte de la rejilla resulta no recomendable, de modo que la salvaguarda no llega a intentarse: queda registrada como no intentada y no activada, y el mayor horizonte permitido es el mismo antes y después de consultarla.

El caso es informativo precisamente por eso. Drift es el modelo entregado, pero lo es por haber ganado la comparación de error fuera de muestra, no por ninguna sustitución: un método de referencia puede ganar la selección cuando la serie se comporta como un paseo aleatorio con deriva, que es lo que ocurre en la mayoría de las rutas del anexo. La salvaguarda solo se activa ante un horizonte no recomendable atribuible al modelo, y aun entonces se limita a informar hasta dónde llegaría cada benchmark; no cambia el modelo entregado.

\subsubsection{Horizonte viable}

El horizonte viable se calcula según longitud, continuidad, errores, estabilidad y comportamiento de la serie. Cuando el usuario solicita una fecha final por fuera de la información real disponible, el módulo de empalme reutiliza el flujo de proyección existente para obtener el índice \ICOCIV{} final estimado. Si la proyección no resulta viable, el empalme no debe calcularse con un valor inventado.

\subsubsection{Intervalos y advertencias}

El módulo presenta advertencias y el estado metodológico cuando la incertidumbre o las restricciones afectan la interpretación; el intervalo de predicción no forma parte de la salida. Las salidas incluyen resultado de proyección, modelo seleccionado, horizonte, estado del horizonte, serie histórica, gráfica, fila fuente, explicación, diagnóstico y metodología. Estos elementos permiten que el usuario comprenda por qué se obtuvo un valor y bajo qué condiciones debe leerse.

\begin{figure}[H]
\centering
\fbox{\begin{minipage}{0.92\textwidth}
\centering
Selección de ruta \ICOCIV{} \(\rightarrow\) obtención de serie histórica \(\rightarrow\) validación de datos \(\rightarrow\) modelos candidatos \(\rightarrow\) backtesting \(\rightarrow\) cálculo de métricas \(\rightarrow\) evaluación de horizonte viable \(\rightarrow\) selección o rechazo de proyección \(\rightarrow\) presentación de resultados \(\rightarrow\) informe DOCX.
\end{minipage}}
\caption{Flujo interno del módulo estadístico ICOCIV}
\label{fig:flujo_estadistico_icociv}
\end{figure}

\subsection{Implementación estadística detallada del módulo ICOCIV}

La implementación estadística explica cómo la aplicación pasa de una ruta jerárquica seleccionada a una proyección defendible. La lógica real se concentra en los paquetes de estadística, validación, proyección y reportes, especialmente en los módulos de métricas, backtesting, servicio de proyección y generación documental. La idea central es sencilla: una proyección no se valida por verse razonable en una gráfica, sino por su desempeño fuera de muestra y por la estabilidad de sus errores \parencite{tashman2000,hyndman2021fpp3}.

\subsubsection{Preparación de la serie histórica}

Una vez el usuario selecciona una ruta \ICOCIV{}, la aplicación resuelve la fila fuente y extrae los valores mensuales disponibles. El vector de trabajo se representa como:

\begin{equation}
y_t = \text{valor del índice ICOCIV en el periodo mensual } t
\end{equation}

El índice interno \(t\) sirve para modelar, pero la interfaz y los reportes conservan el periodo real, por ejemplo 2025-01 o 2025-06. Antes de modelar se revisa que los periodos sean mensuales, que estén ordenados, que no existan duplicados y que los valores sean numéricos. La variación mensual usada para diagnóstico se calcula como:

\begin{equation}
\Delta_t = \left(\frac{y_t}{y_{t-1}} - 1\right) \times 100
\end{equation}

La integración estadística se contrastó con una serie real del Anexo de series del \ICOCIV{} publicado por el \DANE. La ruta localizada por el mismo cargador del backend fue: grupo de obra código 530201, subclase CPC 53211 y tipología 53211\_07, ``Vías urbanas'', dentro de la tabla \texttt{T\_16\_2}. La tabla \ref{tab:ruta_real_icociv} conserva esa trazabilidad.

\begin{table}[H]
\centering
\caption{Ruta jerárquica real utilizada en la integración estadística}
\label{tab:ruta_real_icociv}
\scriptsize
\begin{adjustbox}{max width=\textwidth}
\begin{tabular}{p{3.2cm}p{10.3cm}p{2.2cm}}
\toprule
\textbf{Nivel} & \textbf{Valor} & \textbf{Código} \\
\midrule
Grupo de obra & Carreteras, calles, vías férreas y pistas de aterrizaje, puentes, carreteras elevadas y túneles & 530201 \\
Subclase CPC & Carreteras (excepto carreteras elevadas); calles & 53211 \\
Tipología de obra & Vías urbanas & 53211\_07 \\
\bottomrule
\end{tabular}
\end{adjustbox}
\par\footnotesize Fuente: Anexo de series del \ICOCIV{} publicado por el \DANE, tabla \texttt{T\_16\_2}.
\end{table}

La serie contiene 61 observaciones consecutivas entre enero de 2021 y enero de 2026. La tabla \ref{tab:ejemplo_serie_icociv} presenta los primeros y últimos registros; el procesamiento y las tablas reproducibles fueron generados con las funciones reales de carga, normalización, validación y proyección de la aplicación.

\begin{table}[H]
\centering
\caption{Primeros y últimos registros reales de la serie Vías urbanas}
\label{tab:ejemplo_serie_icociv}
\scriptsize
\begin{adjustbox}{max width=\textwidth}
\begin{tabular}{lr}
\toprule
\textbf{Periodo} & \textbf{Índice \(y_t\)} \\
\midrule
2021-01 & 100,8010 \\
2021-02 & 100,9550 \\
2021-03 & 101,5740 \\
2021-04 & 101,9300 \\
2021-05 & 102,1510 \\
\multicolumn{2}{c}{\(\vdots\)} \\
2025-09 & 136,4500 \\
2025-10 & 136,2900 \\
2025-11 & 136,7100 \\
2025-12 & 136,7200 \\
2026-01 & 140,2000 \\
\bottomrule
\end{tabular}
\end{adjustbox}
\end{table}

Con esa serie, la variación de febrero de 2021 es:

\begin{equation}
\Delta_{2021-02} = \left(\frac{100,9550}{100,8010}-1\right)\times 100 = 0,1528\%
\end{equation}

La validación obtuvo 0 meses faltantes, 0 duplicados, 0 valores no numéricos y 0 valores no positivos. La media fue 120,2097; la mediana, 123,5700; la desviación estándar, 12,5135; y el coeficiente de variación, 10,41\%. El cambio acumulado entre 100,8010 y 140,2000 fue:

\begin{equation}
\left(\frac{140,2000}{100,8010}-1\right)\times100=39,0859\%.
\end{equation}

El diagnóstico robusto sobre variaciones mensuales identificó, entre otros, enero de 2023 con variación 4,9265\% y puntaje modificado 12,9966; febrero de 2023 con 2,5881\% y 6,4969; y enero de 2026 con 2,5453\% y 6,3781. Tras la consolidación por periodo, la serie agregada presenta cinco periodos señalados en lugar de las quince alertas que producía el conteo por escala. Cuatro de ellos ---los eneros de 2023, 2024, 2025 y 2026--- corresponden al patrón calendario de cambio de año detectado para esta serie y no se contabilizan como valores atípicos, mientras que febrero de 2023 se clasifica como posible cambio de nivel. La aplicación no elimina automáticamente ninguno de estos valores: los conserva y los comunica con su clasificación para revisión técnica \parencite{nist_outliers}.

\subsubsection{Modelos evaluados y ejemplos}

La aplicación evalúa los diez candidatos del catálogo sin escalonarlos por niveles: todos compiten en cada serie siempre que puedan estimarse. Naive y Drift cumplen además el papel de referencias simples y trazables, frente a las cuales se lee el aporte de los demás, pero no gozan de prioridad alguna en la elección. La trayectoria final se genera con el único modelo seleccionado por su error cuadrático medio fuera de muestra sobre la muestra común.

El candidato robusto merece una precisión de implementación, porque la pérdida de Huber descrita en el marco teórico admite varias realizaciones de software con resultados distintos. La aplicación emplea \texttt{sklearn.linear\_model.HuberRegressor} \parencite{sklearn_huberregressor,pedregosa2011}, ajustado sobre el índice temporal \(t\) como única variable explicativa, con los hiperparámetros fijos \(\epsilon=1{,}35\), \(\alpha=0{,}0001\) y \texttt{max\_iter}\(=2000\). El parámetro \(\epsilon\) es el umbral \(\delta\) de la Ecuación de la pérdida de Huber, cuyo estimador y su compromiso entre robustez y eficiencia bajo errores normales proceden de \textcite{huber1964}. El valor 1,35 es la constante de sintonización convencional asociada a una eficiencia asintótica cercana al 95\,\% frente a mínimos cuadrados en ese caso, y es la que la implementación adopta por omisión \parencite{sklearn_huberregressor}. Es, por tanto, el único de los tres con un fundamento estadístico propio.

Los otros dos son parámetros de la implementación y se declaran como tales, sin atribuirles una función inferencial. El parámetro \(\alpha\), que aplica una regularización \(L_2\), se conserva en el valor por omisión del estimador; una comprobación acotada sobre la ruta de Vías urbanas, recorriendo \(\alpha\) desde 0 hasta \(10^{-1}\), movió la pendiente estimada en menos de \(1{,}1\times 10^{-5}\) en términos relativos, de modo que la elección no gobierna el resultado en series de esta escala. El valor de \texttt{max\_iter} es una cota superior de iteraciones del optimizador, no una decisión estadística. Sobre las rutas de referencia el ajuste converge en 32 iteraciones en Vías urbanas y en 29 en Arena, de modo que la cota adoptada nunca llega a ser la restricción activa; elevarla por encima del valor por omisión de la biblioteca es un margen de seguridad frente a series que requieran más iteraciones, y en las series comprobadas los coeficientes resultan idénticos para cualquier cota a partir de la cual el optimizador converge. No se afirma que estos dos valores sean óptimos ni que estén sintonizados para \ICOCIV{}: se afirma que son los de la implementación, que están documentados, y que el resultado no depende de ellos en el rango comprobado.

No se emplea \texttt{statsmodels.RLM}: aunque esa biblioteca ofrece estimadores M con la misma función de pérdida \parencite{statsmodels_rlm}, no es el motor que la aplicación ejecuta, y describirla como tal impediría reconstruir el resultado desde el documento. Dentro del sistema, este candidato no depura la serie ni sustituye al análisis principal: compite como un modelo más del catálogo, sin condicionarse a la presencia de atípicos, y su permanencia en el ranking depende del mismo backtesting fuera de muestra que se aplica a los demás modelos.

\begin{table}[H]
\setstretch{1}
\centering
\caption{Modelos de proyección considerados por la aplicación}
\label{tab:modelos_proyeccion_detalle}
\scriptsize
\begin{tabularx}{\textwidth}{@{}>{\raggedright\arraybackslash}p{2.6cm}YYY@{}}
\toprule
\textbf{Modelo} & \textbf{Idea y fórmula básica} & \textbf{Cuándo puede funcionar} & \textbf{Cuándo puede fallar} \\
\midrule
Naive & Repite el último valor: \(\hat{y}_{T+h}=y_T\). & Series estables o como benchmark mínimo. & No captura tendencia; en horizontes largos puede quedarse rezagado. \\
Drift & Proyecta la pendiente promedio: \(\hat{y}_{T+h}=y_T+h\frac{y_T-y_1}{T-1}\). & Series con crecimiento gradual y tendencia simple. & Puede exagerar si el primer y último dato reflejan choques. \\
Lineal OLS & Ajusta \(\hat{y}_t=\beta_0+\beta_1t\). & Tendencia aproximadamente lineal e interpretable. & Puede ser espuria o sensible a escalones y cambios estructurales. \\
Logarítmico temporal & Ajusta \(\hat{y}_t=\beta_0+\beta_1\ln(t+1)\). & Trayectorias que desaceleran su crecimiento. & No representa log-variación y puede subajustar cambios recientes. \\
Exponencial/log-lineal & Ajusta \(\ln(y_t)\) y retransforma con corrección. & Crecimiento proporcional con valores positivos. & Puede sobreproyectar y exige controlar el sesgo de retransformación. \\
Robusto Huber & \texttt{sklearn.\allowbreak linear\_model.\allowbreak HuberRegressor} sobre \(t\), con \(\epsilon=1{,}35\), \(\alpha=0{,}0001\) y \texttt{max\_iter}\(=2000\). & Series con atípicos que afectan OLS. & No reemplaza la revisión del dato ni garantiza mejor pronóstico. \\
Holt lineal & Actualiza nivel \(l_t\) y tendencia \(b_t\): \(\hat{y}_{T+h}=l_T+hb_T\). & Series con tendencia persistente y longitud suficiente. & Puede sobreproyectar si la tendencia reciente no continúa. \\
Holt amortiguado & Usa una tendencia reducida por \(\phi\): \(\hat{y}_{T+h}=l_T+(\phi+\cdots+\phi^h)b_T\). & Tendencias que conviene moderar en horizontes largos. & Depende de estimar adecuadamente \(\phi\) y de datos suficientes. \\
Variación y log-\allowbreak variación & Pronostica cambios relativos y reconstruye el nivel. & Series positivas con dinámica más estable en cambios que en niveles. & La acumulación de errores puede crecer con el horizonte. \\
\bottomrule
\end{tabularx}
\end{table}

Con la serie real de la tabla \ref{tab:ejemplo_serie_icociv}, el pronóstico Naive a cualquier horizonte es el último valor observado, \(140,2000\). Para Drift, la pendiente promedio de las 61 observaciones fue:

\begin{equation}
\frac{y_T-y_1}{T-1}=\frac{140,2000-100,8010}{61-1}=0,656650
\end{equation}

por tanto:

\begin{equation}
\hat{y}_{T+1}=140,2000 + 1(0,656650)=140,8567
\end{equation}

En los orígenes reales usados para contrastar Holt, el ajuste lineal con \(\alpha=0,65\) y \(\beta=0,20\) produjo, desde diciembre de 2025, nivel 136,8651, tendencia 0,2895 y pronóstico a un mes 137,1545. El ajuste amortiguado con \(\phi=0,88\) produjo nivel 136,7623, tendencia 0,1471 y pronóstico 136,8917. A 18 meses, los valores fueron 137,5393 y 130,8727, respectivamente, lo que muestra el efecto acumulado de la amortiguación \parencite{holt1957,gardner_mckenzie1985}.

\paragraph{Estimación de los coeficientes de suavizamiento.} Los valores anteriores corresponden a la configuración vigente hasta agosto de 2026, en la que \(\alpha\), \(\beta^{*}\) y \(\phi\) eran \textbf{constantes fijadas internamente}. Esa parametrización se retiró.

El motivo es de fundamentación. La formulación de Holt implementada por el sistema coincide con la de la fuente, pero el \emph{procedimiento de obtención de sus parámetros} no coincidía: \textcite{hyndman2021fpp3} definen el método estimando los parámetros de suavizamiento y los estados iniciales por minimización de la suma de cuadrados de los errores de un paso dentro de muestra, y \textcite{gardner2006} es explícito al respecto: dada la disponibilidad de algoritmos de búsqueda, no queda justificación para emplear parámetros arbitrarios. Presentar el método bajo el nombre de Holt citando esas fuentes, y a la vez fijar \(\alpha=0,65\) y \(\beta^{*}=0,20\) sin ninguna que los respaldara, era una discrepancia entre lo que el sistema declaraba y lo que hacía.

En la versión vigente el sistema estima \(\alpha\), \(\beta^{*}\) y los estados iniciales \(\ell_0\) y \(b_0\) minimizando esa suma de cuadrados, y en la variante amortiguada estima además \(\phi\) dentro del rango \([0{,}80;\,0{,}98]\) que la propia fuente recomienda. Las únicas restricciones fijadas ---\(0<\alpha<1\), \(0<\beta^{*}\le\alpha\) y el rango de \(\phi\)--- proceden de la fuente. La minimización es determinista: una rejilla gruesa fija seguida de un refinamiento local desde su mejor punto, sin componente aleatoria ni dependencia de las bibliotecas del entorno, con lo que se conserva la propiedad que motivó fijar los coeficientes en su momento ---que una misma serie produzca siempre la misma cifra---. La exploración parte de varios puntos porque la superficie de respuesta del suavizamiento exponencial no es necesariamente convexa \parencite{gardner2006}.

La estimación se repite \textbf{en cada origen} de la validación temporal, empleando únicamente la historia disponible hasta ese origen, y de nuevo con toda la historia para el pronóstico final. No existe una estimación única calculada sobre la serie completa y reutilizada dentro del \emph{backtesting}, de modo que los parámetros no incorporan información posterior al origen que se evalúa.

\paragraph{Un resultado que conviene declarar.} Sobre las diez series del anexo, la solución que minimiza la suma de cuadrados satura sistemáticamente los extremos: \(\alpha\to1\) y \(\beta^{*}\to0\) en el caso lineal, y \(\phi\to0{,}80\) en el amortiguado. La lectura es directa y no es un artefacto numérico: cuando el nivel sigue de cerca a la última observación y la pendiente permanece casi constante, el método de Holt \emph{correctamente estimado} se aproxima a un paseo aleatorio con deriva sobre estas series. En una de ellas la suma de cuadrados desciende de 112,90 a 73,92 al pasar de los coeficientes fijos a los estimados, un 35\,\% menos. Esto explica por qué el método Drift ---que es precisamente ese paseo aleatorio con deriva--- domina la selección en la mayoría de las series del anexo, y es coherente con el comportamiento de un índice de precios agregado. La consecuencia práctica es que los dos candidatos de Holt aportan poca diversidad frente a Drift en este conjunto de datos; se conservan porque la selección los compara en igualdad de condiciones y porque en otras series podrían no degenerar.

\subsubsection{Backtesting walk-forward}

El backtesting implementado es de tipo temporal. En cada origen de evaluación, el modelo se entrena con datos anteriores y pronostica uno o varios horizontes futuros. Luego se compara el pronóstico contra el valor observado. Este procedimiento evita fuga de información futura y reproduce mejor el uso operativo de una herramienta de pronóstico \parencite{tashman2000}.

\begin{figure}[H]
\centering
\fbox{\begin{minipage}{0.9\textwidth}
\centering
Definir origen \(r\) \(\rightarrow\) entrenar con datos hasta \(r\) \(\rightarrow\) pronosticar horizonte \(h\) \(\rightarrow\) comparar con el valor observado \(\rightarrow\) calcular error \(e_{r,h}\) \(\rightarrow\) avanzar origen \(\rightarrow\) consolidar métricas por horizonte.
\end{minipage}}
\caption{Flujo de validación walk-forward usado por la aplicación}
\label{fig:flujo_walkforward}
\end{figure}

\begin{table}[H]
\centering
\caption{Orígenes reales de backtesting walk-forward con Drift}
\label{tab:ejemplo_walkforward}
\scriptsize
\begin{adjustbox}{max width=\textwidth}
\begin{tabular}{lllrrrr}
\toprule
\textbf{Origen} & \textbf{Datos usados} & \textbf{Modelo} & \textbf{\(h\)} & \textbf{Pronóstico} & \textbf{Real} & \textbf{Error} \\
\midrule
2025-12 & 60 observaciones & Drift & 1 & 137,3288 & 140,2000 & 2,8712 \\
2025-10 & 58 observaciones & Drift & 3 & 138,1578 & 140,2000 & 2,0422 \\
2025-07 & 55 observaciones & Drift & 6 & 139,9332 & 140,2000 & 0,2668 \\
2025-01 & 49 observaciones & Drift & 12 & 140,1873 & 140,2000 & 0,0127 \\
2024-07 & 43 observaciones & Drift & 18 & 141,4567 & 140,2000 & -1,2567 \\
\bottomrule
\end{tabular}
\end{adjustbox}
\end{table}

Los errores se agrupan por horizonte. Las filas de la tabla \ref{tab:ejemplo_walkforward} ilustran orígenes distintos y no constituyen por sí solas la métrica agregada; el backend usa todas las ventanas válidas de cada \(h\). Esta separación evita mezclar la precisión de uno, seis o dieciocho meses.

\subsubsection{Métricas de error}

El módulo de cálculo de métricas obtiene MAE, RMSE, MAPE, sMAPE, MASE, sesgo medio, \(R^2\), \(R^2\) ajustado, SSE, AIC y AICc. De ese conjunto, las que describen el desempeño fuera de muestra son MAE, RMSE, MAPE, sMAPE y MASE, porque miden magnitud absoluta, penalización de errores grandes, error porcentual y error escalado frente a una referencia naive \parencite{hyndman_koehler2006}. De todas ellas, \textbf{solo el RMSE fuera de muestra decide qué modelo se entrega}; las demás se publican como información de desempeño y no intervienen en la elección (§\ref{sec:criterios_metricas}).

\begin{align}
MAE &= \frac{1}{n}\sum_{i=1}^{n}|y_i-\hat{y}_i| \\
RMSE &= \sqrt{\frac{1}{n}\sum_{i=1}^{n}(y_i-\hat{y}_i)^2} \\
MAPE &= \frac{100}{n}\sum_{i=1}^{n}\left|\frac{y_i-\hat{y}_i}{y_i}\right| \\
sMAPE &= \frac{100}{n}\sum_{i=1}^{n}\frac{2|y_i-\hat{y}_i|}{|y_i|+|\hat{y}_i|} \\
MASE &= \frac{\frac{1}{n}\sum_{i=1}^{n}|y_i-\hat{y}_i|}{\frac{1}{T-1}\sum_{t=2}^{T}|y_t-y_{t-1}|}
\end{align}

La tabla \ref{tab:metricas_reales_por_horizonte} registra la salida reproducible para la serie Vías urbanas del anexo con corte en enero de 2026 (\(n=61\)). \(J_h\) es el número de ventanas válidas del horizonte, con \(N_0=6\).

\begin{table}[H]
\centering
\caption{Métricas por horizonte calculadas con la serie real Vías urbanas}
\label{tab:metricas_reales_por_horizonte}
\scriptsize
\begin{adjustbox}{max width=\textwidth}
\begin{tabular}{rrrrrrr}
\toprule
\textbf{\(h\)} & \textbf{\(J_h\)} & \textbf{MAE} & \textbf{RMSE} & \textbf{MAPE} & \textbf{sMAPE} & \textbf{MASE} \\
\midrule
1 & 55 & 0,7205 & 1,0874 & 0,5855\% & 0,5882\% & 1,2089 \\
3 & 53 & 1,6329 & 2,2318 & 1,3208\% & 1,3303\% & 2,8210 \\
6 & 50 & 2,5161 & 3,3456 & 2,0381\% & 2,0565\% & 4,4983 \\
12 & 44 & 3,7764 & 4,4519 & 3,0530\% & 3,0994\% & 7,2643 \\
18 & 38 & 6,6784 & 7,5202 & 5,2541\% & 5,3673\% & 13,9040 \\
\bottomrule
\end{tabular}
\end{adjustbox}
\end{table}

La prueba numérica paso a paso de estas métricas se presenta en el capítulo de pruebas, dentro del numeral de proyección y validación estadística. En la implementación, la interpretación no depende de una sola métrica. Un RMSE bajo puede ocultar sesgo; un MAPE puede no ser calculable cuando algún valor observado es cero, y entonces la evidencia se lee con las métricas que sí lo son; y MASE se conserva como señal auxiliar, no como veto automático. Ninguna métrica cancela por sí sola un pronóstico: todas se publican con su valor y su unidad, y el resultado final las presenta junto con la comparación contra benchmarks, los residuos, la estabilidad y la evidencia disponible en cada horizonte.

\subsubsection{Diagnóstico residual e intervalos}

Después de ajustar un modelo, la aplicación revisa residuos. El módulo de diagnóstico residual calcula Durbin-Watson, autocorrelación, ACF, PACF, el contraste de Jarque--Bera con implementación propia y la prueba de Ljung-Box. Esta última se apoya en la biblioteca \texttt{statsmodels} en su versión \texttt{0.14.6}, declarada como dependencia obligatoria: se importa sin ruta condicional y, si faltara, la aplicación no arranca. La función empleada es \texttt{statsmodels.stats.diagnostic.acorr\_ljungbox}, invocada con un único rezago igual a $\min(10,\,\lfloor n/5 \rfloor)$ y con \texttt{model\_df}$=0$, valor que corresponde a residuos que no provienen de un modelo autorregresivo y de media móvil ajustado. Cuando el valor-p de Ljung--Box no está disponible ---muestra de residuos insuficiente o residuos constantes--- lo que no es calculable es el diagnóstico para esa serie, no la dependencia, y así se declara en la interfaz y en los informes. Durbin-Watson se usa como diagnóstico descriptivo de autocorrelación de primer orden, mientras Ljung-Box ayuda a identificar autocorrelación residual conjunta \parencite{statsmodels_durbin_watson,statsmodels_ljungbox}. El valor-p de Jarque--Bera se evalúa contra una $\chi^2$ con dos grados de libertad, según la Ecuación~\ref{eq:jb}, y se declara no calculable con menos de ocho residuos \parencite{jarque_bera1987}. Estas alertas no sustituyen el juicio técnico, pero evitan presentar como fuerte una proyección con residuos problemáticos.

El diagnóstico residual publica cuatro contrastes con todos sus campos. La nulidad de la media se contrasta con la prueba \(t\) de una muestra, implementada con \texttt{scipy.stats.ttest\_1samp} en su forma bilateral; se reportan el número de residuos, la media, el error estándar, el estadístico, los grados de libertad y el valor-p. La constancia de la varianza se contrasta con \texttt{statsmodels.stats.diagnostic.het\_breuschpagan}, con matriz de regresores formada por la constante y el índice temporal, de modo que el estadístico del multiplicador de Lagrange tiene un grado de libertad; se reportan el estadístico, el valor-p, el número de residuos, los regresores empleados y la limitación por tamaño muestral cuando procede. La autocorrelación conjunta se contrasta con Ljung--Box, ya descrita, y la normalidad con Jarque--Bera. Cuando un contraste no es calculable ---por muestra insuficiente, dispersión nula o estadístico no finito--- se declara con su motivo, en lugar de omitirlo o de sustituirlo por una regla alternativa.

Los resultados se redactan en una de tres formas: se rechazó la hipótesis nula al nivel seleccionado, no se rechazó al nivel seleccionado, o el contraste no fue calculable. La aplicación no afirma que los residuos sean independientes, normales u homocedásticos, porque no rechazar una hipótesis nula no la demuestra. La consecuencia operativa de los cuatro contrastes es informativa: ninguno modifica el modelo seleccionado, el pronóstico, el intervalo de predicción ni el horizonte admisible.

Los errores inusuales de la validación temporal se identifican con el mismo puntaje z modificado que las variaciones de la serie, con el umbral \(|M_i|>3{,}5\) \parencite{nist_outliers}. Se publican la cantidad, la proporción, el periodo, el error absoluto y el puntaje de cada ventana. Esa información es descriptiva: la proporción de ventanas señaladas no bloquea, no degrada el horizonte y no altera el modelo ni la banda.

Las métricas de error se comunican con su valor y su unidad. La única lectura comparativa que la aplicación añade es la de MASE respecto de 1, que es el sentido con el que la métrica fue definida \parencite{hyndman_koehler2006}, junto a los cocientes de RMSE y MAE frente a los benchmarks Naive y Drift evaluados en el mismo backtesting. No se asignan categorías cualitativas a MAPE, sMAPE ni al sesgo medio.

Los intervalos de predicción del 80\,\% y del 95\,\% \textbf{se construyeron} ---como diagnóstico interno, no como salida publicada--- con los errores fuera de muestra del horizonte exacto de cada paso: la banda de \(h=6\) se calcula con errores de \(h=6\), nunca con errores de \(h=1\) reescalados. Sobre esa muestra se aplica la construccion de \textcite{hyndman2021fpp3}: el intervalo es \(\hat{y}_{T+h}\pm c\,\hat{\sigma}_h\), con \(\hat{\sigma}_h=\sqrt{\sum_i e_{h,i}^{2}/n_h}\) y \(c\) el cuantil correspondiente al nivel nominal. Como \(\hat{\sigma}_h\) se \emph{estima} a partir de \(n_h\) errores, el multiplicador exacto bajo el supuesto de normalidad de la fuente es el cuantil de una \(t\) con \(n_h\) grados de libertad, que converge al normal cuando \(n_h\) crece. La banda es \textbf{simetrica alrededor del pronostico}: cuando los errores tienen media distinta de cero, \textcite{hyndman2021fpp3} situan el remedio en el pronostico ---sumarle esa media--- y no en el intervalo, de modo que el sesgo se absorbe ensanchando la banda, porque \(\hat{\sigma}_h\) emplea \(\sum e^{2}\) y no la desviacion centrada. El intervalo del 80\,\% queda contenido en el del 95\,\% por construccion. La cobertura alcanzada se verifica retrospectivamente y se reporta: el nivel nominal no es una garantia. Si un paso no cuenta con al menos tres errores del horizonte, no se construye banda alguna para ese paso. El ancho relativo del intervalo del 95\,\% se calcula como:

\begin{equation}
\label{eq:ancho_relativo_ic95}
AnchoRelativoIC95_h=\frac{IC95_{sup,h}-IC95_{inf,h}}{\hat{y}_{T+h}}
\end{equation}

Un ancho relativo alto no cambia el valor puntual y \textbf{tampoco cambia su clasificación}. Hasta agosto de 2026, nueve cortes de amplitud ---escalonados por tramo de horizonte--- degradaban un horizonte de proyección técnica a escenario de alta incertidumbre o lo bloqueaban. Eran literales internos sin fuente que autorizara esa consecuencia, y su medición sobre el anexo de referencia registró cero activaciones. Se retiraron: el ancho relativo permanece como diagnóstico interno con su valor, y una amplitud grande se advierte sin alterar el estado del horizonte.

La cobertura se verifica mediante \emph{evaluación de origen móvil}. Los errores de cada horizonte se ordenan por origen de pronóstico y cada uno se contrasta contra el rango construido exclusivamente con los errores \emph{anteriores} del mismo paso. El error evaluado no interviene en el rango que lo evalúa y no se consulta ningún error posterior, de modo que la evaluación no incorpora información que no estuviera disponible en el momento de cada origen. El procedimiento empleado se publica junto con la cifra, de modo que el lector sepa con qué método se obtuvo y no deba suponerlo a partir del texto.

Este procedimiento sustituye a una partición temporal fija, que dividía la muestra de errores en un tramo de construcción y otro de medición. Aquella exigía que ambos tramos fueran utilizables, y por esa razón los horizontes largos quedaban sin cobertura medida: con quince errores a doce meses y nueve a dieciocho ---los tamaños que la ventana expansiva produce sobre las series disponibles--- no había muestra suficiente para partir, y eran precisamente los horizontes que se consultan para programación presupuestal anual. La evaluación de origen móvil no divide la muestra: con $n$ errores obtiene $n-2$ contrastes, donde el dos es el número de errores previos que la construcción del rango necesita para existir.

Conviene subrayar que se trata de dos cálculos distintos sobre el mismo insumo. El intervalo que se construye internamente emplea \emph{todos} los errores del paso, según la Ecuación~\ref{eq:ip_conformal}, y no cambia por efecto de esta verificación; lo que la evaluación de origen móvil produce es una medida de cuán a menudo ese procedimiento habría acertado, no una banda alternativa. Ninguno de los dos resultados llega al usuario en la versión productiva.

La evaluación de origen móvil tiene una limitación conocida: \textbf{amortigua los cambios de régimen}. Como su rango incorpora el cambio a medida que este ocurre, ante un quiebre tardío en la dispersión del error informará una cobertura más favorable que una partición fija calibrada solo con el tramo previo al quiebre. No es, por tanto, un método uniformemente superior, sino un estimador distinto: aporta más puntos de contraste y menos sensibilidad al cambio de régimen. Como en el resto del tratamiento de intervalos, no se invoca ninguna garantía de cobertura del método conformal.

La verificabilidad del paso solicitado y la clasificación del intervalo continúan rigiéndose por los criterios operativos internos ya descritos, que no se modifican con este cambio de método de evaluación.

Conviene distinguir cuatro nociones que con frecuencia se leen como una sola. El \emph{nivel nominal} es la probabilidad que el método declara, aquí el 95\,\%. La \emph{cobertura observada} es la proporción de observaciones que efectivamente quedaron dentro de la banda, medida sobre una muestra finita y publicada junto al número de errores con que se calculó, porque con pocas evaluaciones esa proporción admite muy pocos valores posibles. El \emph{tipo de banda} describe qué es el intervalo entregado, según el vocabulario que se detalla más adelante. Y el \emph{estado del horizonte} es una propiedad distinta de las tres anteriores y se comunica por separado. Que una banda se haya calculado para un nivel nominal del 95\,\% no afirma que esa cobertura se cumpla: lo que se cumple o no es la cobertura observada, que se reporta aparte.

Conviene precisar por qué el tipo de banda no distingue entre cuantil observado y respaldo paramétrico, que sería la clasificación aparentemente natural. La Ecuación~\ref{eq:ip_conformal} toma el mayor entre el cuantil de orden con corrección de muestra finita y la predicción \(t\) de Student, y sobre las series del anexo \textbf{el término paramétrico domina en los cinco horizontes evaluados}, incluso allí donde el cuantil de orden existe. Rotular esas bandas como empíricas sería una afirmación falsa, y rotularlas por el tamaño de muestra describiría una posibilidad y no la construcción efectiva. El discriminante correcto ---cuál de los dos términos dominó--- no separa nada sobre estos datos y fluctúa entre horizontes por razones ajenas a la calidad de la banda.

Una banda solo se entrega cuando existe. Si alguno de sus límites no es finito, si el límite superior resulta menor que el inferior, o si el paso no dispone de errores propios del horizonte exacto, el intervalo se declara no calculable y no se presenta como banda válida ni interviene en la clasificación del horizonte. Los límites no se intercambian para corregir el orden: un orden invertido no es un defecto de presentación, sino la señal de que el cálculo que los produjo no es válido. Esta comprobación es previa a cualquier lectura de la amplitud, porque una amplitud negativa o inexistente no admite comparación con ningún valor de referencia.

\subsubsection*{Vocabulario y criterios de la banda evaluada (regimen retirado)}

\textbf{Advertencia de vigencia.} Los tres apartados que siguen describen el régimen con que se nombraba y se calificaba la banda de predicción \emph{mientras se evaluó}. \textbf{Ninguno está vigente}: la versión final no publica el intervalo, no publica su cobertura, y la cobertura \textbf{no decide} el estado de ningún horizonte. Se conservan porque documentan la decisión de retirarlos y la evidencia que la sostiene, que es lo que hace auditable una omisión; deben leerse en pasado.

El identificador con que se nombra cada banda viajaba a la interfaz, a los informes y a las exportaciones. Tras el retiro del intervalo ninguno se publica: los tres viven en el resultado como diagnóstico interno, y se documentan aquí porque de esa tipificación depende la trazabilidad de la decisión. Su redacción no era un detalle de estilo: es lo que el lector retiene. La aplicación emplea tres identificadores, y ninguno afirma que la cobertura se cumpla:

\begin{itemize}
  \item \textbf{Banda calculada para un nivel nominal del 95\,\%}: la banda existe, sus límites son válidos y su cobertura se evaluó sobre el paso solicitado. La cifra observada se publica al lado.
  \item \textbf{Rango de referencia}: la banda existe matemáticamente, pero se sostiene en evidencia limitada o en un respaldo cuya cobertura no quedó confirmada en el paso entregado. Se muestra como referencia, no como cobertura verificada.
  \item \textbf{Banda no calculable}: la banda no existe, por las tres imposibilidades matemáticas ya descritas.
\end{itemize}

Se retiraron los identificadores anteriores ---\emph{nominal}, \emph{admisible con advertencia}, \emph{cobertura insuficiente} y \emph{no verificable}---, porque tres de los cuatro emiten un juicio sobre la cobertura observada y el primero se lee como cobertura confirmada del 95\,\%. La expresión ``nivel nominal del 95\,\%'' se conserva, pero calificando el \emph{nivel} para el que la banda se construyó, nunca el resultado obtenido: es un hecho de la construcción, no una propiedad verificada. Esta separación entre declarar el nivel y el procedimiento, por un lado, y afirmar una propiedad del resultado, por otro, es la que exige la Guía ISO/IEC 98-3 en su apartado 6.

Es importante subrayar qué \emph{no} significa el primer identificador. ``Banda calculada para un nivel nominal del 95\,\%'' no significa que la cobertura observada sea del 95\,\%, ni que se haya verificado, ni que esté garantizada: significa que la banda se calculó para ese nivel y que su cobertura pudo medirse. El valor medido puede quedar por debajo del nivel declarado, y cuando así ocurre se informa explícitamente junto al número de contrastes que sostienen la cifra.

La cobertura observada \textbf{no se publica en esta versión}: al retirarse el intervalo, la calidad de una banda que el usuario no recibe deja de ser un resultado y queda como diagnóstico interno. Cuando se publicaba, se hacía siempre como \textbf{recuento y proporción}: ``16/22'' junto a ``0,727''. La proporción por sí sola induce a error cuando las evaluaciones son pocas, porque admite muy pocos valores posibles: 1,000 sobre siete contrastes y 1,000 sobre veintiuno se leen igual y no significan lo mismo. Junto a ambas se publica la \textbf{diferencia frente al nivel nominal} en puntos porcentuales, calculada como \(100\times(p-0{,}95)\); así, una cobertura de 0,727 muestra \(-22{,}3\) puntos y una de 1,000 muestra \(+5{,}0\). Esa diferencia es \textbf{estrictamente descriptiva}: informa a qué distancia quedó lo observado del nivel para el que la banda se calculó, y no aprueba, no degrada, no bloquea ni sustituye a ninguno de los criterios operativos, que siguen decidiendo el estado del horizonte.

El criterio que regía esa publicación, hoy sin efecto por el retiro, era publicarla \textbf{siempre que existiera una evaluación real}, con independencia de cuántos contrastes la sostengan. Conviene separar tres cosas que es fácil confundir: que la cobertura se haya \emph{medido}, que se \emph{publique} y que pueda \emph{usarse} en la clasificación del horizonte. Un horizonte a doce meses reúne quince errores fuera de muestra sobre las series disponibles, que por origen móvil producen trece contrastes: una cobertura de 10/13 es una medición real y se informa como tal, junto a su proporción y su diferencia frente al nivel nominal. Que esos trece contrastes no alcancen el criterio operativo interno vigente afecta a si esa cifra puede intervenir en la clasificación, \textbf{no a si se muestra}. La aplicación lo declara explícitamente: publica el número de evaluaciones disponibles y advierte de que el criterio vigente exige un mínimo mayor para utilizar esa cobertura al clasificar. Ese mínimo es un criterio operativo del proyecto, no un requisito estadístico universal ni una garantía.

Conviene precisar un caso extremo que el sistema debe saber distinguir. Una cobertura observada de \textbf{0,000} ---ninguna de las observaciones evaluadas cayó dentro de la banda--- es un \emph{resultado válido}, y se publica como ``0/n'' con su diferencia de \(-95{,}0\) puntos. No debe confundirse con la \emph{ausencia} de medición, que es lo que ocurre cuando el paso no reúne contrastes suficientes. Son situaciones distintas: la primera dice que la banda falló siempre, la segunda que no se pudo comprobar. Por la misma razón, un horizonte cuya cobertura sea 0,000 entra con normalidad en el cálculo de la cobertura mínima de la trayectoria y en la advertencia de consistencia entre horizontes, en lugar de quedar excluido de ellos.

\subsubsection*{Qué evidencia decide el estado del horizonte solicitado}

En aquel régimen, el estado del horizonte solicitado se decidía con la evidencia de \emph{ese mismo horizonte}: su cobertura observada, su número de contrastes fuera de muestra y la validez de su banda. La cobertura de los demás pasos de la trayectoria no lo degradaba. Hoy ninguna de esas tres magnitudes interviene en el estado del horizonte.

El motivo es estadístico y no de conveniencia. \textcite{christoffersen1998} distingue la cobertura incondicional de la condicional y sostiene evaluar cada horizonte por separado; \textcite{diebold1995} y \textcite{clark2013} muestran que los errores de horizontes distintos comparten historia de entrenamiento y están correlacionados entre sí, de modo que el mínimo entre ellos no es un estadístico con distribución conocida. Usar ese mínimo para degradar un paso concreto carecía de respaldo.

La cobertura mínima de la trayectoria no se descartaba: se conservaba, se publicaba, se localizaba en el horizonte donde ocurría y se acompañaba del número de contrastes que la sostenían. Cuando quedaba por debajo del umbral de aceptación en un horizonte distinto del solicitado, la aplicación emitía una \emph{advertencia de consistencia entre horizontes} que declara tres cosas: que otro horizonte muestra menor cobertura, que ello no invalida automáticamente el horizonte solicitado, y que la lectura conjunta de la trayectoria requiere cautela porque los pasos intermedios no cubren igual que el paso entregado.

El caso que ilustra la distinción es una serie de insumos evaluada a seis meses. Su banda cubrió las veintiuna observaciones contrastadas ---cobertura observada de 1,000--- y aun así se presentaba como escenario, porque la banda del paso a cuatro meses había cubierto 0,762. Al ingeniero que consultaba seis meses se le entregaba un pronóstico cuya banda no falló nunca en la evaluación, rotulado por el comportamiento de un paso que no había pedido. Con la regla vigente, ese horizonte se entrega como proyección técnica y la cobertura menor del paso a cuatro meses aparece, con su horizonte y su conteo, como advertencia separada.

Sobre el conjunto de cincuenta escenarios de verificación, aquel criterio modificaba el estado de \textbf{tres}, sin alterar ningún pronóstico puntual, ningún límite de banda, ninguna métrica, ningún modelo seleccionado, el horizonte máximo recomendado ni el ajuste de cambio de año.

\subsubsection*{Alcance declarado de los criterios de cobertura}

De los tres criterios operativos internos ---mínimo de dieciséis contrastes para considerar verificable un paso, y umbrales de 0,90 y 0,80 sobre la cobertura observada---, \textbf{ninguno decide ya el estado del horizonte}. Los tres conservan su valor y su presencia en el texto de las advertencias; ninguno cambia un modelo, un horizonte ni una clasificación.

La razón es la misma para los tres y es de fondo: \textbf{no se identificó fuente que los sustente}. Ninguna referencia fija una cobertura observada mínima por debajo de la cual un intervalo deje de poder comunicarse, ni un número de errores por debajo del cual una medición válida deba retirarse de la salida. Una muestra corta limita la \emph{precisión} de la cobertura medida; no impide medirla ni proyectar.

El umbral de \textbf{0,90 se convirtió en comunicación descriptiva}. La medición que lo motivó es directa: sobre los cincuenta escenarios de verificación, aislado del corte de 0,80 ---con el que se confundía por el orden en que se comparaban---, su efecto propio fue de \textbf{cero} cambios de estado, \textbf{cero} cambios de tipo de banda y \textbf{cero} cambios numéricos. La razón es visible en la propia regla: las dos categorías que 0,90 separaba compartían identificador de banda, etiqueta visible y consecuencia sobre el horizonte, de modo que lo único que elegía era la redacción del mensaje. Conserva su valor ---no se sustituyó por 0,95 ni por ningún otro corte--- y perdió su función decisoria, que ahora se resuelve fuera de la escalera de clasificación para que no pueda recuperarla por un reordenamiento accidental de las comparaciones.

Lo que sustituye a esa categoría como canal de comunicación no es otro umbral, sino las magnitudes que sostienen la cifra: recuento \(x/y\), proporción, nivel nominal declarado, distancia en puntos porcentuales, número de errores fuera de muestra del paso y método con el que se evaluó la cobertura. Seis cantidades verificables en lugar de una etiqueta que emitía un juicio.

La objeción que en su momento retuvo estos cortes era razonable y conviene enunciarla, porque su respuesta es la que sostiene la regla vigente: retirarlos haría que una banda que cubre 0,727 recibiera el mismo identificador que otra que cubre 1,000, perdiéndose la única señal categórica que las distinguía. La respuesta es que esa señal se sustituye, no se elimina, y por seis cantidades verificables en lugar de una etiqueta: recuento \(x/y\), proporción, nivel nominal declarado, distancia en puntos porcentuales, número de errores fuera de muestra del paso y método de evaluación. Una banda que cubre 0,727 y otra que cubre 1,000 se distinguen ahora por su cifra y por su advertencia, no por un rótulo derivado de un umbral sin sustento.

Sobre los cincuenta escenarios de verificación, la retirada de estos tres cortes ---junto con la de los demás criterios sin fuente--- promueve veintisiete horizontes a proyección técnica y no degrada ninguno. Cada promoción es trazable a la retirada de un criterio concreto, y el pronóstico puntual y las métricas de esos horizontes son exactamente los que ya se calculaban: ninguna promoción alteró una cifra.

Lo único que sigue impidiendo entregar un horizonte es que \textbf{no se pueda calcular}, y tras la revisión independiente de agosto de 2026 eso se reduce a dos causas. La primera es que el propio pronóstico puntual no sea un número finito. La segunda es que el horizonte exceda la cota de existencia de la evidencia: con ventana expansiva, evaluar el horizonte \(h\) requiere \(n-N_0-h+1\ge 1\) ventanas, de donde \(h\le n-N_0\); por encima de esa cota no hay ningún error fuera de muestra, no porque el resultado sea malo sino porque el dato no existe. Los estados de la banda dejaron de bloquear: describen un intervalo que esta versión no publica y que, conforme al principio de que una deficiencia del intervalo no invalida el valor puntual, se registran como diagnóstico. Son imposibilidades aritméticas, no juicios de calidad.

Hasta agosto de 2026 figuraba una cuarta causa: que el MAPE o el sMAPE no fueran finitos. Se retiró porque no era una imposibilidad de la misma naturaleza. El MAPE no está definido cuando algún valor observado es cero y el sMAPE cuando se anula la suma de los valores absolutos del observado y el pronóstico; en ambos casos lo indefinido es \emph{la razón porcentual}, no el pronóstico, cuyo error cuadrático medio, error absoluto medio y error escalado siguen existiendo. No se identificó fuente que sustente cancelar una proyección porque una métrica porcentual carezca de denominador, y el criterio se apoyaba en presentar esa carencia como falta de evidencia. Sobre las diez series del anexo evaluadas en veinticuatro horizontes el criterio no se activó ninguna vez, de modo que su retirada no relaja ningún caso observado: elimina una regla sin sustento. Cuando la métrica no es calculable se publica como tal y la evidencia se lee con las métricas que sí lo son.

Se mantiene, como en todo el tratamiento de intervalos, que \textbf{no se invoca la garantía de cobertura del método conformal}: esa garantía exige intercambiabilidad, y los errores obtenidos por validación con ventana expansiva no la cumplen. Los umbrales son criterios operativos internos del proyecto y no reglas estadísticas universales; su sensibilidad está medida y documentada.

\paragraph{Por qué esta subsección no incluye una tabla de intervalos.} Versiones anteriores de este documento presentaban aquí una tabla titulada «Intervalos de predicción reproducibles de la serie Vías urbanas», con los extremos inferior y superior y el ancho relativo para \(h = 1, 3, 6, 12\) y 18. Esa tabla \textbf{se retiró}, y su retirada es coherente con la decisión que se documenta a continuación: si el intervalo no forma parte de la salida de la aplicación, publicar sus extremos exactos para la serie de referencia sería entregarlo por otra vía. Tampoco se publican el semiancho, el ancho relativo ni la desviación con que se construía, porque cualquiera de ellos, junto al pronóstico puntual ---que sí se publica---, reconstruye la banda.

Lo que sí se conserva es la \emph{evidencia que motivó el retiro}, porque es el fundamento de la decisión y debe poder auditarse. Las mediciones de cobertura que siguen describen, por tanto, una construcción \textbf{ya retirada del producto}, y no una propiedad de la salida actual. La Tabla~\ref{tab:cobertura_empirica_icociv} recoge la medición sobre la serie de referencia.

\begin{table}[H]
\centering
\caption{Cobertura empírica del método de intervalos evaluado y retirado (no publicado), serie Vías urbanas}
\label{tab:cobertura_empirica_icociv}
\scriptsize
\begin{adjustbox}{max width=\textwidth}
\begin{tabular}{rrrr}
\toprule
\textbf{\(h\)} & \textbf{Ventanas de calibración} & \textbf{Ventanas de prueba} & \textbf{Cobertura al 95\,\%} \\
\midrule
1 & 12 & 13 & 0,92 \\
2 & 12 & 12 & 1,00 \\
3 & 11 & 12 & 0,92 \\
4 & 11 & 11 & 1,00 \\
5 & 10 & 11 & 1,00 \\
6 & 10 & 10 & 1,00 \\
\bottomrule
\end{tabular}
\end{adjustbox}
\end{table}

En esta serie la cobertura al 95\,\% se situó entre 0,92 y 1,00, dentro de lo que la resolución de la medición permite distinguir del nivel nominal: con diez a trece ventanas de prueba esa resolución es de aproximadamente ocho puntos porcentuales, de modo que el resultado tampoco \emph{confirma} el nivel nominal.

Ese resultado, sin embargo, no puede generalizarse. La verificación se extendió a una muestra de ocho series de distinto nivel de agregación ---agregado nacional, grupo de obra, subclase, dos tipologías, capítulo constructivo y dos insumos--- evaluadas en los horizontes de uno a seis meses, con veinticuatro combinaciones en total. La cobertura media del intervalo del 95\,\% fue de 0,90, con un mínimo de 0,64 y un máximo de 1,00; solo el 62\,\% de las combinaciones alcanzó o superó 0,90. Los casos de menor cobertura corresponden a series más volátiles, en particular el capítulo constructivo (0,64 a seis meses) y el grupo de obra (0,75 a tres meses), mientras que las series agregadas y suaves se mantuvieron entre 0,85 y 1,00.

La conclusión honesta es que el procedimiento aproximaba el nivel nominal del 95\,\% en las series agregadas y estables, pero \textbf{no lo alcanzaba de manera uniforme} en las series más volátiles. Esa heterogeneidad, unida a la ausencia de una construcción completa sustentada en el sentido que exigen los requisitos metodológicos del proyecto, es lo que condujo a \textbf{retirar el intervalo de todas las salidas}. La banda nominal del 80\,\%, que la aplicación calculaba de forma auxiliar, se había retirado antes por la misma razón y con evidencia aún más desfavorable: cobertura media de 0,77 y mínima de 0,55, es decir, no ofrecía la cobertura que su nombre anunciaba.

\textbf{Estado vigente.} Ninguna de las dos bandas se publica. No se publican sus extremos, su ancho, su ancho relativo, su desviación, sus cuantiles, su método de construcción ni su cobertura; tampoco el vector completo de errores fuera de muestra, que junto al pronóstico puntual permitiría reconstruirlos. La cobertura dejó además de \emph{decidir}: no degrada el horizonte, no lo clasifica y no genera advertencias al usuario. Ambos cálculos permanecen dentro de la aplicación como diagnóstico interno que no interviene en ningún resultado, de modo que la decisión de retirarlos sigue siendo auditable. De la evaluación fuera de muestra sí se publican los \emph{agregados} ---RMSE, MAE, MASE, MAPE y sMAPE cuando son calculables, sesgo medio y número de ventanas---, que sostienen la lectura del error sin reconstruir banda alguna.

Retirar el intervalo no afirma que la incertidumbre no exista: afirma que esta versión no puede acotarla con un método cuya construcción completa esté sustentada. Es una limitación declarada, y el pronóstico puntual se publica cuando es calculable.

Estas cifras no fueron un objetivo de ajuste: los umbrales del procedimiento no se calibraron contra la muestra de prueba, pues hacerlo habría convertido la verificación en una tautología. Debe advertirse además que la medición era conservadora por construcción, ya que calibraba con la mitad de las ventanas disponibles mientras la banda calculada usaba todas ellas.

\subsubsection{Ajuste de cambio de año}

El componente de cambio de año \textbf{mide y publica} el patrón descrito en el marco teórico; desde el 12 de agosto de 2026 \textbf{no lo trata}. El diagnóstico empírico sobre el anexo vigente muestra que el patrón es prácticamente universal en los niveles agregados del índice ---en la tabla de tipos de obra, 46 de 46 series presentan evidencia, con una razón mediana de 7,83 y un salto mediano de 2,57 \%--- y mayoritario, aunque no universal, en insumos, donde 7\,585 de 11\,658 series superan el criterio. La serie Vías urbanas presenta un salto mediano de 2,19 \% frente a un movimiento mensual típico de 0,24 \%, mientras el acero, con un salto de \(-0{,}57\) \% frente a un movimiento típico de 1,58 \%, no supera el criterio: su variación está dominada por el mercado internacional y no por el calendario. Esta medición es el resultado que el componente entrega hoy.

Mientras el tratamiento estuvo activo, no se aplicaba por el solo hecho de detectarse: se exigían dos condiciones acumulativas ---que la serie presentara evidencia según el criterio robusto y que sobre las ventanas walk-forward reales el ajuste no deteriorara ni el MAE ni el RMSE---, y para esta última verificación \(\gamma\) se reestimaba en cada origen usando únicamente la historia anterior, de modo que la comparación no incorporaba información futura. La validación se agregaba sobre un conjunto fijo de horizontes de referencia ---uno, tres y seis meses--- para que la decisión fuera una propiedad de la serie y del modelo. \textbf{Ese mecanismo completo se retiró}: hoy la variable que registra la aplicación del ajuste vale siempre \emph{falso}, y ninguna de estas puertas llega a evaluarse sobre el pronóstico.

Una versión intermedia exigió una tercera condición ---que el horizonte solicitado contuviera al menos un enero--- y produjo una inconsistencia detectada en la validación previa al empaquetado: el factor se evaluaba por paso, pero la condición se comprobaba sobre el horizonte \emph{total}, de modo que una solicitud de doce meses activaba el ajuste en todos los pasos mientras que una de tres lo desactivaba por completo, y el valor proyectado para un mismo mes cambiaba según lo que se pidiera. Como la información histórica disponible es idéntica en ambos casos, la estimación para un mes dado no puede depender de la longitud consultada. La condición se eliminó primero, y el tratamiento entero después; el registro se conserva porque explica por qué el perfil se mide por paso y no por horizonte solicitado.

La razón de la retirada no está en este comportamiento, que era coherente ---en un ciclo completo de doce meses el factor valía exactamente uno, de modo que el crecimiento anual del modelo permanecía inalterado---, sino en que ninguno de sus componentes contaba con sustento completo: la mediana de los saltos pasados como estimador del salto futuro, la forma del factor con su término \(j/12\) ---falso para el modelo ingenuo---, los tres cortes de detección calibrados sobre el propio anexo de aplicación y el conjunto de validación \((1,3,6)\), literal sin fuente. El fenómeno está medido y es contundente; lo que falta es un método publicado para tratarlo. No se introdujo reemplazo, porque una heurística nueva reproduciría el mismo defecto.

La Tabla~\ref{tab:validacion_ajuste_calendario} conserva la validación retrospectiva que se realizó entonces, porque documenta el comportamiento del tratamiento retirado y sostiene el juicio sobre él: el valor exactamente neutro en \(h=12\) corrobora la propiedad de la fórmula, el acero no cambia porque no llegaba a evaluarse, y el concreto simple ---pese a presentar evidencia de patrón--- se deterioraba en \(h=2\) y \(h=3\), de modo que el sistema lo rechazaba. Ninguna de estas cifras describe el cálculo vigente, en el que el ajuste no interviene.

\begin{table}[H]
\centering
\caption{Mejora porcentual en MAE por el ajuste de cambio de año, validación walk-forward}
\label{tab:validacion_ajuste_calendario}
\scriptsize
\begin{adjustbox}{max width=\textwidth}
\begin{tabular}{lrrrrrr}
\toprule
\textbf{Serie} & \(h=1\) & \(h=2\) & \(h=3\) & \(h=6\) & \(h=12\) & \(h=18\) \\
\midrule
Vías urbanas & 32,7 & 36,3 & 38,6 & 40,9 & 0,0 & 28,5 \\
Caminos vecinales & 31,1 & 39,6 & 46,5 & 37,4 & 0,0 & 4,3 \\
Carreteras y calles & 30,7 & 33,8 & 38,7 & 35,8 & 0,0 & 12,9 \\
Cemento & 3,2 & 13,7 & 12,5 & 15,4 & 0,0 & 3,6 \\
Concreto simple & 2,7 & \(-0{,}5\) & \(-6{,}0\) & 2,2 & 0,0 & 2,0 \\
Acero & 0,0 & 0,0 & 0,0 & 0,0 & 0,0 & 0,0 \\
\bottomrule
\end{tabular}
\end{adjustbox}
\end{table}

La trazabilidad del perfil queda registrada en el resultado, en la interfaz y en las exportaciones: si hay evidencia, número de transiciones observadas, salto mediano, movimiento mensual típico, razón entre ambos, consistencia de signo y eneros dentro del horizonte. Todos son medidas de la serie observada. El campo de aplicación del tratamiento viaja igualmente y su valor es siempre \emph{no aplicado}, porque desde agosto de 2026 el ajuste está retirado: la trayectoria que se publica es la del modelo base, sin factor de calendario. Esta versión tampoco publica intervalo de predicción, de modo que no hay banda que construir sobre trayectoria alguna.

\subsubsection{Horizonte viable, advertencias y bloqueo}

El horizonte solicitado es el número de meses que el usuario quiere proyectar. El horizonte evaluado es cada \(h\) que la aplicación somete a backtesting y diagnóstico. El máximo recomendado es el mayor horizonte con evidencia suficiente para uso técnico, y el máximo admisible incluye además los aceptados como escenario con advertencia. Desde agosto de 2026 \textbf{ninguno de los dos exige contigüidad}: si un horizonte intermedio no se sostiene, los posteriores que sí lo hacen se publican igualmente, y el hueco se informa como tal. Ninguna fuente exige que los horizontes válidos formen un prefijo consecutivo, y el encadenamiento anterior hacía que una sola falla temprana cancelara toda la trayectoria posterior. En la configuración interna de criterios, los horizontes rápidos de interfaz son 1, 3, 6, 12 y 18 meses; ese valor de 18 meses es una opción operativa, no un máximo estadístico fijo.

El objetivo de limitar proyecciones no es impedir el cálculo de corto plazo, sino evitar que el usuario extienda la serie hacia periodos sin evidencia suficiente para respaldar esos valores como insumo técnico. Por eso la clasificación distingue \emph{no poder medir} de \emph{haber medido mal}, mediante dos reglas complementarias:

\begin{itemize}
    \item \textbf{Número de ventanas walk-forward disponibles (dato descriptivo).} La aplicación informa cuántas ventanas de validación reúne cada horizonte, porque de ese número depende cuánta evidencia sostiene la comparación. Los cortes que en su momento convirtieron ese recuento en un veto ---seis ventanas para admitir el horizonte, tres para evaluarlo--- carecían de fuente que autorizara la consecuencia y \textbf{fueron retirados}: hoy una evidencia \emph{escasa} se comunica como advertencia y no cancela por sí sola el pronóstico puntual, conforme al principio de que una deficiencia de la evidencia no invalida automáticamente el valor puntual. Conviene distinguir esa escasez de la \emph{inexistencia}: mientras $W \geq 1$ el horizonte tiene evidencia ---por limitada que sea--- y se entrega con su advertencia; cuando $W = 0$ no hay ninguna ventana que medir, y eso no es un umbral de calidad sino la ausencia del dato, que sí impide evaluar el horizonte.
    \item \textbf{Bandas del ancho relativo del IC95 (diagnóstico interno, no publicado).} Cada tramo de horizonte tiene tres bandas: aceptable, escenario con advertencia y bloqueo (por ejemplo, para \(h\leq3\): hasta 10\% aceptable, de 10\% a 30\% escenario y más de 30\% bloqueo; los límites crecen con el horizonte hasta 45\%--75\% para \(h>18\)). Estos cortes son criterios operativos internos y, con el intervalo retirado de las salidas, quedan como diagnóstico no publicado. El encadenamiento consecutivo que propagaba el bloqueo a los horizontes posteriores también se retiró: ninguna fuente exige que los horizontes válidos formen un prefijo.
\end{itemize}

Con estas reglas, el resultado de cada horizonte queda en uno de tres estados: permitido sin advertencias, permitido con advertencias registradas, o no evaluable. Una clasificación de tres niveles previa incluía un cuarto nivel intermedio ---``escenario de alta incertidumbre'', entre la advertencia y la no evaluabilidad---, pero la comparación que lo distinguía dejó de coincidir con ningún valor real desde el cierre del 08-08-2026, de modo que ese cuarto nivel había quedado inalcanzable; el código correspondiente se retiró en esta revisión. El estado \emph{bloqueado por desempeño} desapareció: ningún valor de una métrica cancela ya un horizonte, porque los cortes que lo hacían ---MAPE, sMAPE, proporción de errores inusuales, cocientes frente a los benchmarks y amplitud del intervalo--- carecían de fuente que autorizara esa consecuencia y hoy se publican con su valor. Cuando se niega un horizonte, el mensaje identifica el horizonte afectado y la causa, que solo puede ser una de dos, ambas imposibilidades y no juicios de calidad: que el pronóstico puntual no sea finito, o que no exista el dato con el que evaluarlo ---es decir, $W = 0$, lo que ocurre exactamente cuando $h > n - N_0$---. Ninguna magnitud de error, ninguna amplitud y ningún recuento intermedio de ventanas niegan un horizonte.

\begin{table}[H]
\centering
\caption{Determinación dinámica del horizonte en la serie real Vías urbanas}
\label{tab:ejemplo_horizonte_viable}
\scriptsize
\begin{adjustbox}{max width=\textwidth}
\begin{tabular}{rrlrrl}
\toprule
\textbf{\(h\)} & \textbf{\(J_h\)} & \textbf{Modelo} & \textbf{RMSE} & \textbf{MAPE} & \textbf{Uso comunicado} \\
\midrule
1 & 55 & Drift & 1,0874 & 0,5855\% & Proyección operativa \\
3 & 53 & Drift & 2,2318 & 1,3208\% & Proyección operativa \\
6 & 50 & Drift & 3,3456 & 2,0381\% & Proyección técnica de corto/mediano plazo \\
12 & 44 & Drift & 4,4519 & 3,0530\% & Proyección extendida con cautela \\
18 & 38 & Drift & 7,5202 & 5,2541\% & Escenario estadístico extendido \\
\bottomrule
\end{tabular}
\end{adjustbox}
\end{table}

Conviene distinguir dos columnas que se leen fácilmente como una sola. El \emph{uso comunicado} (última columna) depende únicamente de la longitud del horizonte y crece de ``proyección operativa'' a ``escenario estadístico extendido'' según la tabla \ref{tab:matriz_conceptual}; no es una clasificación estadística ni una puerta. La \emph{clasificación estadística} interna, en cambio, es la misma en las cinco filas: los cinco horizontes están permitidos con advertencias registradas (``técnica con cautela''), porque el modelo seleccionado es Drift y esa condición por sí sola añade una advertencia descriptiva ---``el método seleccionado es un benchmark/escenario simple''---, sin degradar ni limitar ningún horizonte. En este caso real, el máximo recomendado y el máximo admisible coinciden en 55 meses: con \(n=61\) y \(N_0=6\), la cota de existencia \(h\leq n-N_0\) es exactamente ese valor, y ningún horizonte dentro de ella resulta no recomendable. El ejemplo demuestra que el máximo no se fija a priori ni depende del tipo de modelo ganador: depende únicamente de la evidencia fuera de muestra disponible para cada horizonte. La interfaz debe comunicar el estado exacto del horizonte solicitado y no convertir la advertencia de un horizonte en bloqueo global de los demás.

\subsubsection{Tratamiento y presentación de resultados}

La reproducibilidad se comprobó conservando periodo, valor proyectado, modelo, horizonte, intervalos y clasificación. La tabla \ref{tab:reproducibilidad_real_icociv} muestra un fragmento de la salida asociada a los seis primeros meses posteriores al último dato real.

\begin{table}[H]
\centering
\caption{Fragmento reproducible de la proyección de la serie Vías urbanas}
\label{tab:reproducibilidad_real_icociv}
\scriptsize
\begin{adjustbox}{max width=\textwidth}
\begin{tabular}{lrlrrrl}
\toprule
\textbf{Periodo} & \textbf{Proyectado} & \textbf{Modelo} & \textbf{\(h\)} & \textbf{Factor de actualización} & \textbf{\(W\)} & \textbf{Uso comunicado} \\
\midrule
2026-02 & 140,8567 & Drift & 1 & 1,0047 & 55 & Proyección operativa \\
2026-03 & 141,5133 & Drift & 2 & 1,0094 & 54 & Proyección operativa \\
2026-04 & 142,1700 & Drift & 3 & 1,0141 & 53 & Proyección operativa \\
2026-05 & 142,8266 & Drift & 4 & 1,0187 & 52 & Proyección técnica de corto/mediano plazo \\
2026-06 & 143,4832 & Drift & 5 & 1,0234 & 51 & Proyección técnica de corto/mediano plazo \\
2026-07 & 144,1399 & Drift & 6 & 1,0281 & 50 & Proyección técnica de corto/mediano plazo \\
\bottomrule
\end{tabular}
\end{adjustbox}
\end{table}

El resultado estadístico se almacena como un diccionario estructurado que incluye serie, validación, análisis, modelo seleccionado, candidatos evaluados, ranking, backtesting, métricas, horizonte, intervalos, advertencias y proyecciones. Esta estructura alimenta la interfaz y los reportes. En el DOCX, el generador de reportes incorpora resumen ejecutivo, ruta jerárquica, serie compacta, modelos evaluados, backtesting, determinación de horizonte, parámetros reproducibles, gráfica, residuos, limitaciones y referencias estadísticas.

La tabla de salida no oculta modelos descartados. Si Drift gana frente a Holt o si un modelo condicional no se ejecuta por longitud insuficiente, esa razón debe aparecer como parte de la trazabilidad. De esta forma, el usuario puede reconstruir el flujo completo: ruta \ICOCIV{} seleccionada, vector mensual, validación, modelos probados, métricas calculadas, modelo elegido, horizonte aceptado y resultado final.

\subsection{Módulo de empalme ICCP--ICOCIV}

\begin{figure}[H]
\centering
\fbox{\begin{minipage}{0.92\textwidth}
\centering
Fecha inicial y fecha final \(\rightarrow\) serie \ICCP{} seleccionada y ruta \ICOCIV{} equivalente \(\rightarrow\) \(I_0\) \ICCP{} e \(I\) \ICCP{} \(\rightarrow\) \(R_1\) \(\rightarrow\) \(I_0\) \ICOCIV{} e \(I\) \ICOCIV{} \(\rightarrow\) \(R_2\) \(\rightarrow\) \(R=R_1+R_2\) \(\rightarrow\) valor actualizado \(\rightarrow\) tablas y exportación.
\end{minipage}}
\caption{Flujo de cálculo del empalme ICCP--ICOCIV}
\label{fig:flujo_empalme_iccp_icociv}
\end{figure}

\subsubsection{Anexo 10 ICCP histórico}

El Anexo 10 \ICCP{} histórico se carga como fuente interna de índices \ICCP. Esto evita exigir al usuario cargar manualmente el archivo en cada cálculo y permite separar las series disponibles de acuerdo con su naturaleza. La aplicación utiliza el periodo de transición de diciembre de 2021 para decidir si el cálculo usa solamente \ICCP, solamente \ICOCIV{} o ambos tramos.

\subsubsection{Selección de Total ICCP, Canasta general y Grupo de obra}

La aplicación diferencia tres tipos de serie \ICCP: Total \ICCP, Canasta general y Grupo de obra. Esta separación evita mezclar categorías. Cuando el usuario selecciona Total \ICCP, la serie disponible corresponde únicamente al índice total. Cuando selecciona Canasta general, se presentan series como equipos, materiales, transporte, mano de obra y costos indirectos. Cuando selecciona Grupo de obra, se presentan grupos como obras de explanación, subbases y bases, transporte de materiales, aceros y elementos metálicos, concretos, morteros y obras varias, entre otros grupos disponibles en el anexo.

La selección se guarda con tipo de serie, serie seleccionada, ruta \ICCP{} completa, índices inicial y final, factor del tramo y observaciones de trazabilidad.

\subsubsection{Selección de ruta ICOCIV equivalente}

El usuario selecciona una ruta \ICOCIV{} equivalente mediante el selector jerárquico. La equivalencia no se asume automáticamente; debe justificarse técnicamente. La tabla de equivalencias muestra el insumo contractual, la serie \ICCP{} realmente seleccionada y el insumo o ruta \ICOCIV{} equivalente. Esta presentación permite revisar si el criterio usado es coherente con el insumo analizado.

\subsubsection{Unidad y datos contractuales}

El formulario de empalme solicita el insumo contractual, la unidad, el precio base y los valores necesarios para aplicar las fórmulas. La unidad se selecciona desde una lista cerrada con las opciones: m lineal, m2, m3, kg, Unidad y Global. La opción ``Sin selección'' funciona únicamente como marcador visual y no se acepta como dato válido para calcular.

\subsubsection{Fórmula general de transición}

La fórmula general implementada corresponde a las ecuaciones \ref{eq:r1_general}, \ref{eq:r2_general}, \ref{eq:r_total} y \ref{eq:valor_actualizado_general}. En este caso, \(R_1\) representa el primer valor parcial ajustado con \ICCP, \(R_2\) representa el segundo valor parcial ajustado con \ICOCIV{} y \(R\) corresponde a la suma \(R_1+R_2\). El valor actualizado se obtiene sumando \(R\) a la base ajustable.

\subsubsection{Fórmula especial para acero}

La fórmula especial para acero corresponde a las ecuaciones \ref{eq:r1_acero}, \ref{eq:r2_acero}, \ref{eq:r_acero} y \ref{eq:z_acero}. Este flujo usa \(P_0\), \(I_x\) y \(q\) para calcular el valor adicional por fluctuación del acero. La aplicación separa esta ruta del cálculo general para evitar mezclar variables y para conservar la trazabilidad propia del caso especial.

\subsubsection{Trazabilidad de los cálculos}

Los ejemplos numéricos paso a paso del cálculo general y del cálculo especial de acero se presentan en el capítulo de pruebas, dentro de los numerales específicos del módulo \ICCP--\ICOCIV. En la app real, los índices \(I\) e \(I_0\) se toman de la serie \ICCP{} seleccionada y de la ruta \ICOCIV{} equivalente, o de la proyección \ICOCIV{} cuando la fecha final supera el último dato real disponible.

\subsubsection{Tablas de equivalencias y cálculo}

La tabla \ref{tab:equivalencias_modelo} presenta el modelo de equivalencias usado por la aplicación. En la interfaz, cada cálculo válido agrega una fila a esta tabla, de manera que varios insumos pueden documentarse dentro de la misma sesión.

\begin{table}[H]
\centering
\caption{Equivalencias entre insumos contractuales, grupos ICCP e insumos ICOCIV}
\label{tab:equivalencias_modelo}
\scriptsize
\begin{adjustbox}{max width=\textwidth}
\begin{tabular}{lll}
\toprule
\textbf{Insumo contractual} & \textbf{Grupo de obra equivalente ICCP} & \textbf{Insumo equivalente ICOCIV} \\
\midrule
MEZCLA CONCRETO 280 MPa & Concretos, morteros y obras varias & Concreto simple \\
HIERRO DE 60.000 PSI 420 MPa & Aceros y elementos metálicos & Acero corrugado \\
RECEBO & Subbases y bases & Recebo \\
\bottomrule
\end{tabular}
\end{adjustbox}
\end{table}

La tabla \ref{tab:calculo_valor_ajustado_modelo} muestra la estructura de cálculo del valor ajustado. Los encabezados de índices incorporan los periodos seleccionados por el usuario, de forma que el resultado no queda desligado de sus fechas de cálculo.

\afterpage{%
\begin{landscape}
\begin{table}[H]
\centering
\caption{Modelo de tabla de cálculo del valor ajustado}
\label{tab:calculo_valor_ajustado_modelo}
\scriptsize
\begin{adjustbox}{max width=\linewidth}
\begin{tabular}{lllllllllll}
\toprule
\textbf{Insumo} & \textbf{Unidad} & \textbf{Precio APU} & \textbf{\ICCP{} \(I_0\)} & \textbf{\ICCP{} \(I\)} & \textbf{\(R_1\)} & \textbf{\ICOCIV{} \(I_0\)} & \textbf{\ICOCIV{} \(I\)} & \textbf{\(R_2\)} & \textbf{\(R\)} & \textbf{Valor ajustado} \\
\midrule
MEZCLA CONCRETO 280 MPa & m3 & \$ 1.000.000,00 & 100,00 & 105,00 & \$ 50.000,00 & 100,00 & 110,00 & \$ 105.000,00 & \$ 155.000,00 & \$ 1.155.000,00 \\
\bottomrule
\end{tabular}
\end{adjustbox}
\end{table}
\end{landscape}%
}

\subsubsection{Integración con proyección ICOCIV}

Cuando la fecha final solicitada en el empalme está dentro de los periodos reales del archivo \ICOCIV{} cargado, la aplicación usa el índice real de la ruta seleccionada. Cuando la fecha final supera el último periodo disponible, el empalme no genera una extrapolación propia: reutiliza el módulo de proyección existente para obtener el índice \ICOCIV{} final. Esa integración conserva el modelo seleccionado, el estado del horizonte, las advertencias y las salidas de la pestaña de proyección.

Si el índice final usado es proyectado, el resultado debe quedar identificado como tal en el detalle técnico, la tabla de cálculo y la trazabilidad. Esta marca es necesaria porque un índice proyectado no equivale a un dato oficial publicado.

\subsection{Exportación de resultados}

\subsubsection{Informe DOCX}

El informe DOCX documenta resultados de análisis y proyección cuando el flujo de usuario lo requiere. El generador de reportes incorpora identificación de la ruta \ICOCIV, serie histórica, resumen estadístico, modelo seleccionado, métricas de validación, horizonte, resultados de proyección, gráfica, fila fuente, diagnóstico, metodología y observaciones. Este reporte permite conservar una evidencia externa al entorno de la aplicación.

\subsubsection{Excel de empalme y trazabilidad}

El Excel de empalme exporta resultados del cálculo, equivalencias, índices, periodos, fórmulas y trazabilidad. Conforme al requisito RF-10, la implementación genera una sola hoja con secciones consecutivas: información general, trazabilidad del cálculo, equivalencias \ICCP--\ICOCIV, cálculo del valor ajustado con fórmulas verificables, metodología general, metodología y detalle del caso especial de acero cuando existen cálculos de ese tipo, y observaciones. Las columnas \(R_1\), \(R_2\), \(R\) y valor actualizado se exportan como fórmulas de Excel, y el detalle de acero referencia por fórmula la fila del mismo cálculo dentro de la misma hoja, lo que permite reconstruir y auditar cada resultado.

\subsection{Validaciones, mensajes y manejo de errores}

Las validaciones incluyen fecha inicial y final, selección de serie \ICCP, ruta \ICOCIV, unidad, base ajustable, disponibilidad de índices y viabilidad de proyección. Los mensajes explican el problema para evitar fallos silenciosos. Entre los casos controlados se encuentran la falta de unidad válida, el uso de ``Sin selección'', la ausencia de serie \ICCP, la falta de ruta \ICOCIV{} equivalente, fechas fuera del rango disponible y proyecciones no viables.

\subsection{Consideraciones de usabilidad}

La interfaz busca reducir sobrecarga visual, separar módulos, mostrar resultados principales y permitir detalle técnico cuando el usuario lo requiera. Los botones de cálculo, limpieza, detalle, edición, eliminación y exportación se organizaron para que el flujo de trabajo sea acumulativo: el usuario puede calcular varios insumos, revisar cada uno, corregirlos y exportar el conjunto.

\section{PRUEBAS DE SOFTWARE}

\subsection{Diseño general de las pruebas}

Las pruebas se organizaron como verificación funcional, estadística y documental de los módulos principales. La estrategia combinó: (i) inspección del código que carga, filtra y proyecta; (ii) casos reproducibles construidos con el Anexo de series del \ICOCIV{} publicado por el \DANE; (iii) pruebas unitarias existentes para series válidas y casos adversos; y (iv) revisión de las salidas de interfaz y exportación. Cada caso se documentó con los campos de la tabla \ref{tab:protocolo_pruebas}, de modo que un resultado esperado no se confundiera con un resultado efectivamente obtenido.

\begin{table}[H]
\centering
\caption{Protocolo de documentación de pruebas}
\label{tab:protocolo_pruebas}
\small
\begin{tabularx}{\textwidth}{@{}p{3.2cm}Y@{}}
\toprule
\textbf{Campo} & \textbf{Contenido exigido} \\
\midrule
Objetivo & Comportamiento o cálculo que se busca comprobar. \\
Entrada & Archivo, ruta, periodos, valores o acción ejecutada. \\
Procedimiento & Secuencia reproducible y componente de la aplicación involucrado. \\
Criterio de aceptación & Condición observable que permite clasificar el caso. \\
Resultado y evidencia & Valor obtenido, tabla, mensaje, función o archivo que permite revisarlo. \\
Interpretación y limitación & Alcance técnico del resultado y asuntos que quedan pendientes. \\
\bottomrule
\end{tabularx}
\end{table}

Los estados usados son \textit{Conforme}, cuando existe resultado reproducible o trazabilidad suficiente; \textit{Conforme por inspección}, cuando se comprobó la lógica en código pero la acción gráfica no se repitió en esta revisión; y \textit{No conforme}, cuando la salida real contradice un requisito. Esta validación no sustituye una auditoría institucional ni convierte una proyección en dato oficial.

\subsection{Pruebas funcionales del módulo ICOCIV}

\textbf{Objetivo.} Verificar que el cargador leyera el anexo oficial, detectara las tablas de la hoja \textit{Anexo 16}, identificara los encabezados mensuales y conservara las columnas descriptivas necesarias para la selección jerárquica.

\textbf{Entrada y procedimiento.} Se utilizó el Anexo de series del \ICOCIV{} publicado por el \DANE. Se revisó el flujo del cargador que localiza los bloques A16/A16.x, separa columnas fijas y periodos, convierte los valores de índice y ordena los meses. El resultado detectó doce tablas: \texttt{T\_16}, \texttt{T\_16\_1}, \texttt{T\_16\_2}, \texttt{T\_16\_3}, \texttt{T\_16\_6}, \texttt{T\_16\_7}, \texttt{T\_16\_8}, \texttt{T\_16\_9}, \texttt{T\_16\_10}, \texttt{T\_16\_11}, \texttt{T\_16\_12} y \texttt{T\_16\_13}.

\textbf{Criterio y resultado.} La prueba se consideró conforme si las tablas y periodos quedaban disponibles sin mezclar campos descriptivos con valores mensuales. La ruta de prueba de la tabla \texttt{T\_16\_2} produjo una serie de 61 observaciones, desde enero de 2021 hasta enero de 2026. Los extremos 100,8010 y 140,2000 coincidieron con la evidencia tabular generada a partir del archivo. La prueba confirma la capacidad de carga para este anexo; no garantiza compatibilidad automática con futuros archivos cuya estructura cambie.

\subsection{Pruebas del selector jerárquico}

\textbf{Objetivo.} Comprobar que cada lista dependiera de la selección superior y que la fila fuente, la ruta visible, la serie y el reporte conservaran la misma identidad.

\textbf{Entrada.} Se seleccionó el grupo 530201, ``Carreteras, calles, vías férreas y pistas de aterrizaje, puentes, carreteras elevadas y túneles''; la subclase 53211, ``Carreteras (excepto carreteras elevadas); calles''; y la tipología 53211\_07, ``Vías urbanas'', dentro de \texttt{T\_16\_2}.

\textbf{Procedimiento y evidencia.} Se inspeccionó el controlador que filtra cada nivel con los valores ya elegidos y resuelve una sola fila al completar la ruta. También se revisó que un cambio en un nivel superior ejecute el reinicio de todas las listas inferiores y actualice la etiqueta de ruta. La serie recuperada contiene los mismos 61 valores usados en las pruebas estadísticas.

\textbf{Criterio y resultado.} La prueba fue conforme por inspección: la selección superior restringe las opciones inferiores, los niveles descendientes se limpian al cambiar el padre y la fila resuelta alimenta serie, gráfica y trazabilidad. Este control es crítico porque un fallo de ruta puede producir un índice numéricamente válido pero asociado a una obra distinta.

\subsection{Pruebas de proyección y validación estadística}

La validación estadística se desarrolló con la ruta real de Vías urbanas descrita anteriormente. El periodo disponible fue enero de 2021--enero de 2026 y el último índice fue 140,2000. El propósito no fue encontrar un modelo universal, sino verificar que la aplicación preparara la serie, comparara candidatos con orden temporal, cuantificara error e incertidumbre y comunicara los límites del horizonte.

\subsubsection{Preparación, continuidad y casos adversos}

\textbf{Procedimiento.} La serie se convirtió a frecuencia mensual, se ordenó y se verificaron duplicados, meses faltantes, valores no numéricos y valores no positivos. El caso real obtuvo 61 observaciones, cero meses faltantes, cero duplicados, cero valores no numéricos y cero valores no positivos; por tanto, quedó habilitado para transformaciones logarítmicas, backtesting y proyección.

La batería de pruebas del proyecto también define cuatro casos adversos: duplicado mensual, hueco temporal, serie de ocho observaciones y salto atípico. El criterio es que el duplicado produzca error crítico; el hueco quede marcado para revisión; el salto se comunique sin borrar el dato; y la serie corta \textbf{entregue el pronóstico puntual declarando que la evidencia es limitada}. Este último criterio cambió en agosto de 2026: antes se exigía que la serie corta \emph{bloqueara} la proyección, exigencia que se retiró junto con el mínimo de 8 observaciones, de modo que hoy la prueba comprueba que el punto se entregue, que sea finito, que el número de ventanas se declare y que no reaparezca el veto. La lógica de estas validaciones fue conforme por inspección del archivo de pruebas; su ejecución completa no se repitió en el entorno de compilación porque este no incluye \texttt{scikit-learn}.

\subsubsection{Descriptivos, variaciones y atípicos}

\textbf{Resultado reproducible.} La media fue 120,2097, la mediana 123,5700, la desviación estándar 12,5135 y el coeficiente de variación 10,41\%. El índice pasó de 100,8010 a 140,2000, equivalente a una variación acumulada de 39,0859\%. Para febrero de 2021 la variación mensual fue

\begin{equation}
100\left(\frac{100,9550}{100,8010}-1\right)=0,1528\%.
\end{equation}

El control robusto identificó enero de 2023 como observación de atención: variación 4,9265\%, mediana mensual 0,2507\%, MAD 0,2427 y puntaje robusto 12,9966. El criterio de aceptación fue conservar la observación oficial y emitir advertencia; no imputarla ni eliminarla automáticamente.

\subsubsection{Modelos de referencia y suavizamiento}

\textbf{Objetivo.} Verificar que los candidatos produjeran pronósticos trazables y que las diferencias entre tendencia completa y amortiguada fueran visibles. Con origen enero de 2026, Naive mantuvo 140,2000 en todos los horizontes. Drift estimó una pendiente mensual de 0,656650 y produjo 140,8567, 142,1700, 144,1399, 148,0798 y 152,0197 para 1, 3, 6, 12 y 18 meses.

Para un origen en diciembre de 2025, Holt lineal con \(\alpha=0,65\) y \(\beta=0,20\) obtuvo nivel 136,8651, tendencia 0,2895 y pronóstico a un mes 137,1545. Holt amortiguado con \(\phi=0,88\) obtuvo nivel 136,7623, tendencia 0,1471 y pronóstico 136,8917. A 18 meses las proyecciones fueron 137,5393 y 130,8727, respectivamente. El resultado demuestra por qué el modelo debe elegirse con evidencia fuera de muestra y no por extensión visual de la tendencia.

\subsubsection{Verificación del procedimiento walk-forward}

\textbf{Objetivo y procedimiento.} Verificar que cada origen entrenara solo con información anterior al dato evaluado. Para cada horizonte \(h\), el error fue \(e_{r,h}=y_{r+h}-\hat{y}_{r+h}\). La tabla \ref{tab:prueba_walkforward_real} presenta una traza reproducible del modelo Drift.

\begin{table}[H]
\centering
\caption{Verificación walk-forward con la serie real de Vías urbanas}
\label{tab:prueba_walkforward_real}
\scriptsize
\begin{adjustbox}{max width=\textwidth}
\begin{tabular}{lrrrrr}
\toprule
\textbf{Origen} & \textbf{Observaciones} & \textbf{Horizonte} & \textbf{Pronóstico} & \textbf{Observado} & \textbf{Error} \\
\midrule
2025-12 & 60 & 1 & 137,3288 & 140,2000 & 2,8712 \\
2025-10 & 58 & 3 & 138,1578 & 140,2000 & 2,0422 \\
2025-07 & 55 & 6 & 139,9332 & 140,2000 & 0,2668 \\
2025-01 & 49 & 12 & 140,1873 & 140,2000 & 0,0127 \\
2024-07 & 43 & 18 & 141,4567 & 140,2000 & -1,2567 \\
\bottomrule
\end{tabular}
\end{adjustbox}
\end{table}

\textbf{Resultado.} Ningún origen utilizó el dato de enero de 2026 como entrenamiento antes de predecirlo. Los errores se conservaron separados por horizonte, requisito necesario para no atribuir a seis o doce meses la precisión observada a un mes. El procedimiento fue conforme con el criterio de validación temporal fuera de muestra \parencite{tashman2000}.

\subsubsection{Verificación del cálculo de métricas}

\textbf{Objetivo.} Comprobar MAE, RMSE, MAPE, sMAPE y MASE con los errores fuera de muestra de la ruta real. La tabla \ref{tab:prueba_metricas_reales} resume los resultados de Drift por horizonte.

\begin{table}[H]
\centering
\caption{Métricas reproducibles por horizonte para Vías urbanas}
\label{tab:prueba_metricas_reales}
\scriptsize
\begin{tabular}{rrrrrrr}
\toprule
\textbf{Horizonte} & \textbf{Ventanas} & \textbf{MAE} & \textbf{RMSE} & \textbf{MAPE} & \textbf{sMAPE} & \textbf{MASE} \\
\midrule
1 & 25 & 0,8284 & 1,0617 & 0,6261\% & 0,6270\% & 1,1795 \\
3 & 23 & 1,8120 & 2,0780 & 1,3661\% & 1,3680\% & 2,5639 \\
6 & 20 & 2,1063 & 2,4779 & 1,5936\% & 1,5789\% & 2,8980 \\
12 & 14 & 2,0001 & 2,4025 & 1,4899\% & 1,4740\% & 2,7233 \\
18 & 8 & 4,4664 & 4,9492 & 3,2713\% & 3,2069\% & 5,9650 \\
\bottomrule
\end{tabular}
\end{table}

Por ejemplo, en \(h=1\) la suma de los 25 errores absolutos fue \(25\times0,8284=20,7100\), por lo que \(MAE=20,7100/25=0,8284\). En el origen diciembre de 2025, el error absoluto de 2,8712 y la escala naive interna de 0,6639 produjeron

\begin{equation}
MASE=\frac{2,8712}{0,6639}=4,3249.
\end{equation}

\textbf{Interpretación.} Los errores porcentuales son bajos en la escala de esta ruta, pero MASE mayor que uno indica que Drift no supera de manera uniforme la referencia naive escalada. Por tanto, el ranking no puede basarse solo en MAPE. El resultado es conforme porque las métricas se calculan por horizonte, controlan denominadores y permanecen disponibles para comparar modelos.

\subsubsection{Diagnóstico residual e intervalos}
\label{sec:diagnostico_residual_intervalos}

\textbf{Procedimiento.} Cuando un candidato produce residuos suficientes, la aplicación calcula sesgo, Durbin--Watson y diagnósticos de autocorrelación; los hallazgos generan advertencias y no se usan para borrar observaciones. La inspección del módulo confirmó que estos diagnósticos alimentan la explicación del resultado.

El criterio de aceptación de esta comprobación interna fue que el pronóstico quedara contenido en la banda calculada y que la construcción se aplicara por horizonte exacto. \textbf{No se publican aquí los extremos de esa banda para la serie de referencia}: el intervalo no forma parte de ninguna salida de la versión final, de modo que tabular sus límites exactos ---o su ancho relativo, que junto al pronóstico los reconstruye--- equivaldría a entregarlo por otra vía. Lo que la comprobación acredita es que el procedimiento evaluado se comportó como su formulación predice; no que exista una banda que el usuario reciba.


\textbf{Resultado.} El ancho de la banda calculada internamente creció de forma monótona con el horizonte, como corresponde a una dispersión estimada con los errores de cada paso, y el pronóstico quedó contenido en ella en todos los horizontes evaluados. Ese comportamiento es coherente con la formulación, pero \textbf{no bastó para adoptarla}: la construcción completa no alcanzó el sustento exigido y la banda se retiró del producto (véanse la discusión de cobertura y el apartado de limitaciones).

\subsubsection{Verificación de horizonte viable}

\textbf{Objetivo.} Verificar que el sistema evaluara cada horizonte con su propia evidencia fuera de muestra y publicara la magnitud que sostiene esa evaluación, en lugar de aplicar un límite fijo sin relación con la serie. La tabla \ref{tab:prueba_horizonte_viable} contiene el resultado real sobre la ruta de Vías urbanas del anexo con corte en enero de 2026 ($n = 61$), con Drift como modelo seleccionado.

\begin{table}[H]
\centering
\caption{Evaluación por horizonte de la ruta Vías urbanas, con el número de ventanas fuera de muestra que sostiene cada fila}
\label{tab:prueba_horizonte_viable}
\scriptsize
\begin{adjustbox}{max width=\textwidth}
\begin{tabular}{rrrrrrl}
\toprule
\textbf{Horizonte} & \textbf{$W$} & \textbf{MAE} & \textbf{RMSE} & \textbf{MAPE} & \textbf{MASE} & \textbf{Uso comunicado} \\
\midrule
1 mes & 55 & 0,7205 & 1,0874 & 0,5855\% & 1,2089 & Proyección operativa \\
3 meses & 53 & 1,6329 & 2,2318 & 1,3208\% & 2,8210 & Proyección operativa \\
6 meses & 50 & 2,5161 & 3,3456 & 2,0381\% & 4,4983 & Proyección técnica de corto/mediano plazo \\
12 meses & 44 & 3,7764 & 4,4519 & 3,0530\% & 7,2643 & Proyección extendida con cautela \\
18 meses & 38 & 6,6784 & 7,5202 & 5,2541\% & 13,9040 & Escenario estadístico extendido \\
\bottomrule
\end{tabular}
\end{adjustbox}
\end{table}

Los cinco horizontes quedaron permitidos y ninguno fue negado: con $n = 61$ y $N_0 = 6$, la cota de existencia $h \leq n - N_0$ sitúa en $h = 55$ el último horizonte con evidencia fuera de muestra, de modo que los cinco de la tabla la tienen holgadamente. La columna $W$ es el número de ventanas de origen móvil realmente disponibles para ese horizonte, y coincidió exactamente con el recuento del backtesting en las cinco filas. La última columna es una etiqueta de comunicación que depende únicamente de la longitud del horizonte solicitado; no es una clasificación estadística ni una puerta, y no interviene en ninguna decisión. El deterioro monótono de las cinco métricas con el horizonte es la información que el usuario necesita para juzgar el alcance: MASE pasa de $1{,}21$ en un mes a $13{,}90$ en dieciocho, lo que indica que a horizontes largos el modelo no supera a la referencia naive escalada. Esa magnitud se publica; no se convierte en un veto.

El ancho relativo de la banda del 95\,\% ---definido en la Ecuación \ref{eq:ancho_relativo_ic95}--- se calculó también para esta ruta como diagnóstico interno, y no intervino en ninguna de las cinco filas: con el intervalo retirado de las salidas, esa magnitud no clasifica, no degrada y no se publica.

\textbf{Resultado e interpretación.} La salida no autoriza uso automático, pero tampoco lo niega por la magnitud de una métrica. Ninguno de los cinco horizontes fue bloqueado: los cinco se entregan, cada uno con su propia evidencia y con las advertencias que le correspondan. Lo que la tabla acredita es que la evaluación es \emph{por horizonte} y que la magnitud que la sostiene ---\(W\)--- y las que describen el error se publican con su valor, de modo que el ingeniero responsable pueda juzgar hasta dónde le sirve la proyección. El MASE creciente es la advertencia más informativa del caso, y es exactamente eso: una advertencia. La interfaz conserva el estado de cada horizonte, el máximo recomendado y las advertencias, sin convertir ninguna de ellas en una puerta.

\subsubsection{Reproducibilidad de la proyección}

Con último dato 140,2000 en enero de 2026 y modelo Drift, los primeros seis periodos fueron 140,8567; 141,5133; 142,1700; 142,8266; 143,4832 y 144,1399 para febrero--julio de 2026. Repetir el cálculo con la misma ruta, corte, modelo y versión de datos debe reproducir esos valores dentro de la tolerancia de redondeo. La evidencia completa incluye ruta, fila fuente, origen, parámetros, horizontes y métricas; sin esos metadatos, un número proyectado aislado no cumple el criterio de trazabilidad.

\subsubsection{Auditoría de fórmulas contra valores calculados a mano}

\textbf{Objetivo y criterio.} Verificar que cada fórmula estadística implementada coincide con su definición teórica, fijando el resultado esperado de forma independiente del código. Las pruebas automatizadas de fórmulas usan valores calculados a mano o propiedades matemáticas cerradas.

\textbf{Procedimiento y resultado.} Con observados \((100; 102; 104; 106)\) y predichos \((101; 101; 105; 104)\) ---errores \(-1, +1, -1, +2\)--- se verifica MAE \(=1{,}25\), RMSE \(=\sqrt{1{,}75}\), sesgo \(=+0{,}25\), y los valores manuales de MAPE y sMAPE término a término; el MAPE excluye correctamente los observados en cero y el sMAPE respeta su cota teórica de 200. El MASE usa la escala naive del entrenamiento: con entrenamiento \((90; 92; 96)\), escala \(=3\) y MASE \(=1{,}25/3\); cambiar el tramo de prueba no altera la escala, lo que confirma la ausencia de fuga. Durbin--Watson supera 3,5 con residuos alternantes y cae bajo 0,5 con residuos persistentes. Naive y Drift reproducen sus fórmulas cerradas sobre una serie lineal perfecta. En el walk-forward, cada predicción de Drift coincide con la fórmula analítica calculada solo con la historia previa al origen, y un salto artificial de 500 puntos en el último dato no altera ninguna predicción anterior: no existe fuga de información futura. Finalmente, sobre la banda que se conserva como diagnóstico interno se comprueba que el valor puntual queda contenido en ella y que su ancho no decrece con el horizonte ---comprobación de coherencia de la construcción evaluada, no de una salida que el usuario reciba---. Las nueve pruebas pasan.

\subsubsection{Medición de la combinación de pronósticos}

\textbf{Objetivo y criterio.} Decidir con evidencia, y no por criterio, si combinar varios modelos aporta frente a proyectar con un único modelo por serie. El criterio de aceptación fue que la combinación mejorara el error fuera de muestra sin romper la consistencia de la trayectoria.

\textbf{Procedimiento.} Se tomó una muestra aleatoria estratificada de 40 series del anexo vigente, repartidas entre las cinco tablas del \ICOCIV, y se evaluaron los horizontes de 1 a 18 meses, con 720 comparaciones en total. Para evitar que la combinación se midiera sobre los mismos datos con los que se calibraron sus pesos, los pesos y la elección del modelo único se estimaron con la primera mitad de los orígenes walk-forward, y la comparación se midió en la segunda mitad, que ninguno de los dos había visto.

\textbf{Resultado.} Los pesos calculados por horizonte producían una mejora aparente del 13,4 \% en RMSE, pero esa cifra desaparecía al separar calibración y medición: la ganancia provenía de ajustar los pesos sobre los mismos errores con que luego se evaluaban. Medida de forma limpia, la combinación por horizonte perdía frente al modelo único en los horizontes cortos, con un deterioro del 8,9 \% en RMSE y del 16,9 \% en MAE para \(h=1\), y solo aportaba desde \(h\geq6\). La variante de pesos fijos por serie, única compatible con la consistencia de la trayectoria, ganó en apenas 312 de 720 comparaciones y empeoró el RMSE un 9,1 \% en promedio, con ventaja únicamente en la banda estrecha de 4 a 6 meses.

\textbf{Decisión.} No existiendo evidencia de que la combinación de pronósticos aporte bajo el esquema de modelo único por serie, se retiró del sistema junto con su criterio operativo asociado. La aplicación proyecta siempre con un único modelo por serie.

\subsubsection{Evaluación de modelos y justificación de exclusiones}
\label{sec:decisiones_alcance_modelos}

Este apartado reúne las decisiones de alcance sobre los modelos que la literatura describe y que el sistema no incorporó. Se documentan aquí, y no en el marco teórico, porque son resultados de la evaluación realizada en este proyecto y no propiedades generales de los métodos.

\textbf{Suavizamiento exponencial ETS.} La familia ETS descrita en la sección \ref{sec:teoria_ets} se consideró como candidato adicional del catálogo. La evaluación comparativa mostró que, en el entorno de ejecución de la aplicación, la implementación disponible reproducía la trayectoria de Holt con tendencia amortiguada ---modelo ya presente en el catálogo--- de modo que no ampliaba el espacio de hipótesis evaluadas sino que duplicaba un candidato existente. A ello se añadieron dos consideraciones de alcance: la estimación conjunta de parámetros de suavizamiento y estados iniciales exige más historia que la disponible en buena parte de las series del \ICOCIV, cuya longitud típica es de sesenta a setenta observaciones mensuales, y la componente estacional ---principal aportación de ETS frente a Holt--- requeriría varios ciclos anuales completos para identificarse con solvencia. En consecuencia, ETS no forma parte del catálogo de modelos implementados y no figura entre los modelos que la aplicación ejecuta o evalúa.

Esta decisión es de alcance de SAVIP y no una conclusión general sobre el método: ETS conserva plena validez en la literatura y resulta apropiado en series con estacionalidad marcada e historia suficiente. Su incorporación con estimación completa de los parámetros de espacio de estados queda registrada como trabajo futuro, condicionada a disponer de una implementación optimizada y a demostrar mejora del desempeño fuera de muestra frente a Holt amortiguado.

\textbf{Combinación de pronósticos.} La medición descrita en el apartado anterior sostiene la decisión de no combinar. Bajo el esquema de un único modelo por serie, la variante de pesos por horizonte es incompatible con el criterio de trayectoria consistente ---los primeros valores proyectados cambiarían según el horizonte solicitado--- y la variante de pesos fijos por serie, única compatible con ese criterio, no mostró ventaja al medirse sin fuga de información. Se añade una consideración de interpretabilidad propia del uso previsto: el resultado debe poder justificarse ante un tercero mediante una ecuación y unos parámetros identificables, condición que una combinación ponderada no satisface. La decisión es igualmente de alcance: la combinación de pronósticos sigue siendo una práctica válida y su desempeño podría revisarse con un catálogo de modelos más diverso o con series de mayor longitud.

\textbf{Efecto sobre el sistema.} Ninguna de las dos exclusiones altera las cifras que la aplicación produce: ni ETS ni la combinación intervienen en la selección, el pronóstico, el intervalo, el horizonte admisible ni la salvaguarda. Su efecto es de alcance ---acotan el catálogo a modelos interpretables y trazables--- y su limitación es la contrapartida: el sistema no ofrece hoy una alternativa con estacionalidad explícita ni un mecanismo de reducción de varianza por agregación de modelos.

\subsubsection{Verificación del ajuste de cambio de año}

\textbf{Objetivo y criterio.} Comprobar que la \emph{detección} del patrón de cambio de año recupera correctamente su magnitud y que el \emph{tratamiento} no se aplica en ningún caso, de modo que la trayectoria entregada es siempre la del modelo base. Las pruebas automatizadas usan series sintéticas con salto controlado en enero y las series reales del anexo.

\textbf{Procedimiento y resultado.} Con una serie sintética de seis años y salto de 5 \% en cada enero, la detección recupera el valor con error menor a \(10^{-6}\) e identifica cinco transiciones: la \emph{medición} del fenómeno es correcta. Sobre las series reales, Vías urbanas presenta el patrón con claridad, el concreto simple también lo presenta y el acero no supera el criterio de detección. En las tres, y en cualquier otra, el tratamiento queda \emph{no aplicado} y la trayectoria proyectada coincide exactamente con la del modelo base; así lo verifican las pruebas del archivo. Las exportaciones DOCX, PDF, HTML y CSV se generan con la trazabilidad del perfil medido y con el estado de no aplicación declarado.

\subsection{Pruebas del módulo ICCP--ICOCIV}

\subsubsection{Selección de fechas}

\textbf{Objetivo y criterio.} Comprobar periodos antes, durante y después del corte de diciembre de 2021. El sistema debe resolver solo \ICCP, solo \ICOCIV{} o dos tramos según las fechas, y los controles de año/mes deben incrementar y decrementar dentro de su rango. La inspección del evento de cálculo y de los controles confirmó esta separación; la respuesta visual de todos los botones se clasifica como conforme por inspección y requiere repetición manual en la entrega instalada.

\subsubsection{Selección de series ICCP}

\textbf{Entrada y procedimiento.} Se cambió conceptualmente entre Total \ICCP, Canasta general y Grupo de obra y se contrastaron las listas construidas desde el Anexo 10. Total solo admite ``Total ICCP''; Canasta general contiene Equipos, Materiales, Transporte, Mano de obra y Costos indirectos; Grupo de obra contiene exclusivamente los grupos de obra del anexo. Al cambiar de tipo, la selección previa se invalida.

\textbf{Resultado.} La estructura separada y la limpieza del valor dependiente fueron conformes por inspección. ``Sin selección'' funciona como marcador visual y no se acepta como dato de cálculo.

\subsubsection{Selección de ruta ICOCIV equivalente}

El criterio fue que la ruta \ICOCIV{} seleccionada se guardara completa, que su último nivel apareciera como insumo equivalente y que la misma fila suministrara los índices. El controlador y la estructura del registro de historial cumplen esta relación por inspección. La aplicación conserva la equivalencia elegida, pero la pertinencia contractual continúa bajo responsabilidad profesional.

\subsubsection{Unidad obligatoria}

Se verificó que el desplegable contuviera exclusivamente ``Sin selección'', m lineal, m2, m3, kg, Unidad y Global. El criterio fue bloquear el cálculo y mostrar ``Seleccione una unidad válida.'' cuando no existiera una unidad aceptada. La lista cerrada y la validación son conformes por inspección; ninguna unidad adicional forma parte del requisito.

\subsubsection{Cálculo general}

\textbf{Objetivo de la prueba.} Verificar que la aplicación aplique correctamente el cálculo general de transición \ICCP--\ICOCIV, separando el tramo \ICCP, el tramo \ICOCIV{} y la suma \(R=R_1+R_2\).

\textbf{Datos de entrada.}

\begin{table}[H]
\centering
\caption{Datos de entrada para verificación del cálculo general}
\label{tab:prueba_empalme_general_datos}
\scriptsize
\begin{tabular}{ll}
\toprule
\textbf{Variable} & \textbf{Valor} \\
\midrule
\(P\) & \$1.000.000 \\
\(A\) & \$100.000 \\
\(I_0\) \ICCP{} & 100,00 \\
\(I\) \ICCP{} & 108,00 \\
\(I_0\) \ICOCIV{} & 100,00 \\
\(I\) \ICOCIV{} & 112,00 \\
\bottomrule
\end{tabular}
\end{table}

\textbf{Fórmula usada y sustitución numérica.} La base ajustable fue:

\begin{equation}
P-A=1.000.000-100.000=900.000
\end{equation}

El primer tramo, calculado con \ICCP, fue:

\begin{equation}
R_1=(1.000.000-100.000)\left[\left(\frac{108}{100}\right)-1\right]=72.000
\end{equation}

El segundo tramo, calculado con \ICOCIV{} sobre la base ya ajustada por \(R_1\), fue:

\begin{equation}
R_2=(900.000+72.000)\left[\left(\frac{112}{100}\right)-1\right]=116.640
\end{equation}

Por tanto:

\begin{equation}
R=R_1+R_2=72.000+116.640=188.640
\end{equation}

\begin{equation}
\text{Valor actualizado}=(P-A)+R=900.000+188.640=1.088.640
\end{equation}

\textbf{Criterio, resultado e interpretación.} El resultado esperado fue \(R_1=\$72.000\), \(R_2=\$116.640\), \(R=\$188.640\) y valor actualizado \(\$1.088.640\). La sustitución reproduce exactamente \(R=R_1+R_2\) conforme al lineamiento adjunto. En este procedimiento \(R_1\), \(R_2\) y \(R\) son valores monetarios de ajuste y deben conservar formato de moneda; los cocientes de índices son factores adimensionales y no deben confundirse con ellos.

\subsubsection{Cálculo especial de acero}

\textbf{Objetivo de la prueba.} Verificar que el cálculo especial de acero use \(P_0\), \(I_x\), \(q\), \(R_1\), \(R_2\), \(R\) y \(Z\) sin mezclarlos con los campos del cálculo general.

\textbf{Datos de entrada.}

\begin{table}[H]
\centering
\caption{Datos de entrada para verificación del cálculo especial de acero}
\label{tab:prueba_acero_datos}
\scriptsize
\begin{tabular}{ll}
\toprule
\textbf{Variable} & \textbf{Valor} \\
\midrule
\(P_0\) & \$5.000.000 \\
\(I_0\) \ICCP{} & 100,00 \\
\(I\) \ICCP{} & 110,00 \\
\(I_0\) \ICOCIV{} & 100,00 \\
\(I\) \ICOCIV{} & 115,00 \\
\(I_x\) & \$6.500/kg \\
\(q\) & 1.000 kg \\
\bottomrule
\end{tabular}
\end{table}

\textbf{Fórmula usada y sustitución numérica.} El primer ajuste del acero con \ICCP{} fue:

\begin{equation}
R_1=5.000.000\left[\left(\frac{110}{100}\right)-1\right]=500.000
\end{equation}

El segundo ajuste con \ICOCIV{} fue:

\begin{equation}
R_2=(5.000.000+500.000)\left[\left(\frac{115}{100}\right)-1\right]=825.000
\end{equation}

La suma del ajuste fue:

\begin{equation}
R=R_1+R_2=500.000+825.000=1.325.000
\end{equation}

El valor facturado del acero fue:

\begin{equation}
I_x \times q=6.500 \times 1.000=6.500.000
\end{equation}

Finalmente:

\begin{equation}
Z=(I_x \times q)-(R+P_0)=6.500.000-(1.325.000+5.000.000)=175.000
\end{equation}

\textbf{Resultado e interpretación.} El resultado esperado fue \(R_1=\$500.000\), \(R_2=\$825.000\), \(R=\$1.325.000\) y \(Z=\$175.000\). La prueba confirma que el flujo de acero conserva la comparación final entre el valor facturado \(I_x \times q\) y el valor ofertado ajustado \(P_0+R\). En la aplicación, estas variables deben aparecer en el detalle técnico y en la exportación para que el cálculo pueda reconstruirse.

\textbf{Alcance de la reproducibilidad y limitación declarada.} Conviene distinguir tres niveles que no son equivalentes. El primero es la \emph{reproducibilidad matemática}: dadas las entradas, la sustitución en las ecuaciones \ref{eq:r1_acero} a \ref{eq:z_acero} produce siempre los mismos valores, y una verificación independiente con aritmética decimal y de punto flotante reprodujo \(R_1\), \(R_2\), \(R\) y \(Z\) sin diferencia alguna, tanto en este caso de prueba como en el caso con valores de magnitud contractual empleado en la validación. El segundo es la \emph{trazabilidad institucional} de los índices: \(I_0\) e \(I\) proceden del anexo oficial del \DANE{} incluido en el paquete, con su hash, de modo que su procedencia es verificable. El tercero es la \emph{trazabilidad contractual} de las entradas propias del contrato ---\(P_0\), \(I_x\) y \(q\)---, que no está cubierta.

Estas tres variables son datos aportados por quien realiza el ejercicio: el valor base del insumo, el precio facturado por kilogramo y la cantidad ejecutada. En este trabajo se emplearon como valores de entrada de un ejercicio de verificación, y no se acompañan de la fuente primaria que los sustentaría en un caso real ---acta contractual, factura, orden de pago o documento equivalente---, que no forma parte de la información disponible para el proyecto. En consecuencia, el caso del acero debe leerse como \textbf{numéricamente reproducible} y \textbf{parcialmente reproducible como evidencia contractual}: sostiene que la fórmula está correctamente implementada, no que un ajuste concreto haya sido liquidado con esos valores. Presentarlo de otro modo confundiría la verificación de un procedimiento con la validación de un expediente. Aportar la fuente primaria de \(P_0\), \(I_x\) y \(q\) queda registrado como condición previa a cualquier uso del resultado con efectos contractuales.

\subsubsection{Tablas acumulativas}

\textbf{Objetivo y procedimiento.} Verificar que dos cálculos válidos consecutivos conservaran el primero y agregaran una fila a cada tabla. La revisión del manejador mostró que cada registro aceptado se incorpora al historial y que ambas tablas se reconstruyen desde ese historial; limpiar el formulario no elimina los registros aceptados.

\textbf{Resultado.} La acumulación es conforme por inspección de código, pero no se presenta como nueva ejecución gráfica en esta revisión documental. Antes de la entrega institucional debe conservarse evidencia visual de dos cálculos consecutivos con conteos \(1\rightarrow2\) en ambas tablas.

\subsubsection{Integración con la proyección}

El criterio fue usar un índice real cuando la fecha final no superara el último periodo cargado y reutilizar el servicio de proyección existente cuando lo superara. La inspección confirmó la bifurcación y la marca de valor proyectado en el registro. Si el servicio devuelve un horizonte no viable, el cálculo debe bloquearse; no existe autorización metodológica para inventar otro índice. La prueba integral con interfaz queda clasificada como conforme por inspección hasta repetir ambos escenarios en el entorno instalado con sus dependencias completas.

\subsubsection{Exportación a Excel}

\textbf{Requisito.} RF-10 exige una sola pestaña con resultados, fórmulas y trazabilidad. \textbf{Resultado obtenido.} El exportador genera una única hoja con secciones consecutivas: información general, trazabilidad, equivalencias, cálculo del valor ajustado con fórmulas, metodología y detalle de acero cuando aplica. La prueba automática de exportación verifica que el libro contenga exactamente una hoja, que las columnas \(R_1\), \(R_2\), \(R\) y valor actualizado sean fórmulas de Excel por cada fila y que el detalle de acero referencie por fórmula la fila del mismo cálculo dentro de la misma hoja. Por tanto, CP-EXP-01 es \textbf{Conforme}.

\subsection{Pruebas de interfaz y usabilidad}

La revisión abarcó controles de año y mes, listas dependientes, mensajes, tablas y botones de calcular, limpiar, detalle, editar, eliminar y exportar. Los manejadores y validaciones existen y evitan varios fallos silenciosos. Sin embargo, la inspección de código no sustituye la interacción real con foco, teclado y ratón; por ello CP-UI-01 se mantiene como conforme por inspección y debe cerrarse con evidencia manual en el ejecutable final.

\subsection{Resultados de las pruebas}

\begin{landscape}
\scriptsize
\begin{longtable}{@{}>{\raggedright\arraybackslash}p{0.08\linewidth}>{\raggedright\arraybackslash}p{0.14\linewidth}>{\raggedright\arraybackslash}p{0.22\linewidth}>{\raggedright\arraybackslash}p{0.25\linewidth}>{\raggedright\arraybackslash}p{0.13\linewidth}>{\raggedright\arraybackslash}p{0.08\linewidth}@{}}
\caption{Matriz de pruebas funcionales y estadísticas}\label{tab:matriz_pruebas}\\
\toprule
\textbf{ID} & \textbf{Módulo} & \textbf{Entrada o acción} & \textbf{Resultado esperado} & \textbf{Evidencia revisada} & \textbf{Estado} \\
\midrule
\endfirsthead
\toprule
\textbf{ID} & \textbf{Módulo} & \textbf{Entrada o acción} & \textbf{Resultado esperado} & \textbf{Evidencia revisada} & \textbf{Estado} \\
\midrule
\endhead
CP-DAT-01 & Carga \ICOCIV{} & Anexo de series del \ICOCIV{} publicado por el \DANE. & Doce tablas, periodos y valores numéricos detectados. & Cargador y serie real de 61 datos. & Conforme \\
CP-DAT-02 & Selector jerárquico & Ruta 530201/53211/53211\_07. & Ruta, fila fuente y serie coincidentes; reinicio descendente. & Controlador y ventana principal. & Conforme por inspección \\
CP-SER-01 & Serie histórica & Vías urbanas. & Enero 2021--enero 2026; 100,8010--140,2000. & Tablas reproducibles. & Conforme \\
CP-EST-01 & Validación estadística & Serie real y casos adversos definidos. & Continuidad, longitud, faltantes, duplicados y atípicos controlados. & Servicio y archivo de pruebas. & Conforme por inspección \\
CP-EST-02 & Backtesting & Drift, horizontes 1, 3, 6, 12 y 18. & Ventanas sin fuga temporal y errores separados por horizonte. & Tabla \ref{tab:prueba_walkforward_real}. & Conforme \\
CP-EST-03 & Modelos & Naive, Drift, Holt y regresiones condicionadas. & Pronósticos trazables y comparación fuera de muestra. & Ejemplos numéricos reales. & Conforme \\
CP-EST-04 & Métricas & Errores walk-forward. & MAE, RMSE, MAPE, sMAPE y MASE por horizonte. & Tabla \ref{tab:prueba_metricas_reales}. & Conforme \\
CP-EST-05 & Intervalos (evaluados, no adoptados) & Errores por horizonte. & La banda se construye con los errores del horizonte exacto y contiene al pronóstico. & Comprobación interna; no se publica en ninguna salida. & Conforme; procedimiento retirado \\
CP-HOR-01 & Horizonte viable & Solicitudes de 1 a 18 meses. & Estado y advertencia conservados; sin índices inventados. & Tabla \ref{tab:prueba_horizonte_viable}. & Conforme \\
CP-REP-01 & Informe DOCX & Resultado de proyección válido. & Ruta, serie, modelo, métricas, horizonte y metodología. & Generador de reportes. & Conforme por inspección \\
CP-EMP-01 & Empalme general & Caso que cruza diciembre de 2021. & \(R_1=72.000\), \(R_2=116.640\), \(R=188.640\). & Sustitución numérica. & Conforme \\
CP-EMP-02 & Series \ICCP{} & Cambio entre los tres tipos. & Listas separadas y selección previa invalidada. & Widget de empalme. & Conforme por inspección \\
CP-EMP-03 & Unidad & ``Sin selección''. & Bloqueo y mensaje de unidad válida. & Widget de empalme. & Conforme por inspección \\
CP-EMP-04 & Acero & \(P_0\), \(I_x\), \(q\) e índices de prueba. & \(R=R_1+R_2\) y cálculo de \(Z\). & Sustitución numérica. & Conforme \\
CP-EMP-05 & Proyección en empalme & Fecha posterior al último dato real. & Reutilización del servicio y marca de proyectado. & Widget y servicio. & Conforme por inspección \\
CP-TAB-01 & Tablas acumulativas & Dos cálculos válidos. & Conteo \(1\rightarrow2\) en ambas tablas. & Historial y reconstrucción de tablas. & Conforme por inspección \\
CP-EXP-01 & Excel & Historial con cálculos válidos. & Una sola pestaña según RF-10. & Exportador genera una sola hoja con fórmulas verificadas por prueba automática. & Conforme \\
CP-UI-01 & Interfaz & Botones y controles numéricos. & Respuesta visible, rangos válidos y mensajes claros. & Manejadores de eventos. & Conforme por inspección \\
\bottomrule
\end{longtable}
\end{landscape}

\subsection{Síntesis de reproducibilidad de las pruebas}

Los resultados estadísticos pueden reconstruirse con el Anexo de series del \ICOCIV{} publicado por el \DANE, la ruta 530201/53211/53211\_07, el corte enero de 2026 y las tablas de este capítulo. Los cálculos de empalme pueden repetirse con las entradas y sustituciones indicadas. Para una evidencia institucional completa deben archivarse, además, versión del software, archivo fuente, captura de ruta, salida del modelo y exportable generado.

\subsection{Errores encontrados y correcciones}

Durante el desarrollo se corrigieron errores de listas \ICCP, unidad, formato de resultados, encabezados de tablas, controles de fechas y acumulación de resultados. También se verificó, exclusivamente con el lineamiento de empalme adjunto, que \(R_1\) y \(R_2\) son ajustes parciales monetarios y que \(R=R_1+R_2\). La revisión inicial encontró que el exportador generaba cinco hojas; se corrigió para consolidar toda la información en una sola hoja, conforme al requisito RF-10, y la prueba automática de exportación verifica esa condición.

\subsection{Limitaciones de las pruebas}

Las pruebas dependen de los archivos disponibles y de escenarios representativos. No todos los insumos contractuales tienen una equivalencia única entre \ICCP{} e \ICOCIV, y la validez de una proyección depende de la longitud y estabilidad de la ruta seleccionada. La batería Python no se reejecutó en el entorno aislado usado para compilar este documento porque allí falta \texttt{scikit-learn}; por ello se diferenciaron los resultados numéricos reproducibles de las verificaciones por inspección. Quedan como cierres obligatorios la ejecución integral en el entorno instalado y la evidencia visual de dos cálculos consecutivos.

\section{MANUAL DE USUARIO}

\subsection{Requisitos previos}

El usuario debe contar con la aplicación instalada y con los anexos oficiales \ICOCIV{} requeridos para el análisis. También debe disponer de la información contractual básica del insumo que desea revisar: descripción, unidad, precio base, fecha inicial, fecha final y criterio técnico de equivalencia con \ICCP{} e \ICOCIV.

\subsection{Inicio de la aplicación}

El usuario abre la aplicación y accede a los módulos disponibles desde la interfaz principal. La pantalla organiza el flujo en secciones de carga, selección, análisis, resultados, reportes y empalme. Antes de ejecutar cálculos se recomienda verificar que el archivo \ICOCIV{} cargado corresponda a la versión de datos que se desea usar.

\subsection{Carga de archivos ICOCIV}

El usuario selecciona el archivo oficial \ICOCIV{} desde el botón de carga. El sistema lee el archivo, detecta las tablas y organiza los periodos mensuales. Si el archivo no tiene la estructura esperada o no contiene datos suficientes, la aplicación debe mostrar un mensaje de error o advertencia.

\subsection{Uso del selector jerárquico ICOCIV}

El usuario navega por los niveles disponibles hasta seleccionar la ruta de análisis. La selección debe hacerse con cuidado, porque esa ruta define la serie histórica, la proyección y la equivalencia usada en el empalme. Después de elegir una ruta, se recomienda revisar la fila fuente o la serie histórica para confirmar que corresponde al insumo o grupo esperado.

\subsection{Consulta de serie histórica}

La aplicación muestra la serie mensual asociada a la ruta seleccionada. Esta consulta permite revisar periodos disponibles, valores faltantes y comportamiento general antes de proyectar. Si la serie no tiene datos suficientes, el usuario debe seleccionar otra ruta o limitar el análisis a periodos reales disponibles.

\subsection{Ejecución de análisis y proyección}

El usuario define el periodo objetivo y ejecuta el análisis. La aplicación ordena y valida la serie, calcula descriptivos y variaciones, evalúa los modelos candidatos habilitados, ejecuta validación walk-forward por horizonte y muestra el modelo seleccionado. El flujo recomendado es:

\begin{enumerate}
\item confirmar la ruta completa y la fila fuente;
\item revisar que la serie tenga el corte y la continuidad esperados;
\item indicar el periodo futuro solicitado;
\item ejecutar la proyección y leer primero el estado del horizonte;
\item leer el pronóstico puntual junto con las métricas y el estado metodológico, teniendo presente que esta versión no publica intervalo de predicción y que, por tanto, la incertidumbre del resultado no viene acotada;
\item revisar advertencias, diagnóstico, explicación y metodología antes de usar o exportar el valor.
\end{enumerate}

Si la serie es demasiado corta, tiene duplicados críticos o el horizonte resulta no viable, el usuario no debe forzar un valor manual. Debe corregir la fuente, reducir el horizonte o documentar que no existe una proyección defendible.

\subsection{Interpretación de resultados de proyección}

Los resultados deben interpretarse con base en métricas, horizonte viable, evidencia fuera de muestra y observaciones metodológicas. Un índice proyectado no equivale a un índice oficial publicado por el \DANE. Por ello, el usuario debe conservar la evidencia del modelo, la fecha proyectada, el estado del horizonte y las advertencias.

Cuando el panel indique que un horizonte es recomendable, de cautela, escenario o no recomendable, el usuario debe leer esa clasificación junto con las métricas y el estado metodológico. Debe tenerse presente que esta versión no publica intervalo de predicción: la incertidumbre del valor mostrado no viene acotada, y esa ausencia es una limitación declarada, no un descuido de la salida. Un valor puntual puede parecer razonable, pero si el horizonte aparece como escenario o como no recomendable, la proyección no debe usarse como resultado técnico principal sin justificación adicional.

\textbf{Sobre los identificadores ``P0-C'' y ``P0-E'' en los mensajes.} Algunos mensajes de advertencia citan estos dos identificadores internos, y conviene que el lector sepa qué representan sin necesidad de consultar el repositorio del proyecto. No son métricas ni criterios estadísticos nuevos: son los nombres con que este trabajo registra dos limitaciones metodológicas que se declaran explícitamente, en lugar de resolverse con una regla sin sustento. \textbf{P0-C} identifica la ausencia de una construcción de intervalo de predicción adoptable bajo la configuración vigente: la banda que la aplicación calcula internamente no tiene un método completo y sustentado, y por eso no se publica como intervalo institucional (véase la discusión completa en la sección \ref{sec:diagnostico_residual_intervalos}). \textbf{P0-E} identifica que el primer origen de la validación temporal, \(N_0\), no está determinado por ninguna fuente ni derivación suficiente —es una configuración provisional, no un mínimo de identificabilidad demostrado— (véase la sección de limitaciones, \ref{sec:limitaciones_encontradas}). Que un mensaje mencione P0-C o P0-E no cambia el modelo, el pronóstico ni el horizonte entregados: es la aplicación siendo transparente sobre qué parte de su propio procedimiento sigue sin cerrar.

Si la serie presenta un patrón recurrente de cambio de año, el panel lo informa explícitamente. El panel separa dos hechos que no son el mismo. El primero es que la serie presente el patrón, propiedad que se determina con las transiciones diciembre--enero observadas y que no depende de lo que se pida proyectar. El segundo es que ese patrón tenga efecto dentro del horizonte solicitado, lo cual ocurre únicamente cuando alguno de los pasos proyectados cae en enero. A estos dos se suma un tercero, que es el \emph{efecto del tratamiento} y exige simultáneamente que el tratamiento se aplique y que el horizonte contenga al menos un enero. Como esta versión no aplica el tratamiento en ningún caso, ese tercer hecho es siempre negativo, incluso cuando el horizonte sí cruza enero; el panel lo dice explícitamente para que no se lea una cosa por otra. La distinción evita además interpretar como ausencia de patrón lo que solo es ausencia de efecto en un horizonte corto. En consecuencia la trayectoria mostrada es la del modelo base: los valores de enero no se elevan ni los del resto del año se reducen. El usuario debe conservar esta información junto con el resto de la evidencia.

\begin{table}[H]
\centering
\caption{Guía de lectura de las salidas estadísticas}
\label{tab:manual_lectura_metricas}
\small
\begin{tabularx}{\textwidth}{@{}p{2.5cm}p{4.2cm}Y@{}}
\toprule
\textbf{Salida} & \textbf{Qué mide} & \textbf{Cómo interpretarla} \\
\midrule
MAE & Error absoluto promedio en puntos de índice. & Menor es mejor dentro de la misma ruta y horizonte. Conserva la unidad del índice. \\
RMSE & Raíz del error cuadrático promedio. & Penaliza con mayor intensidad los errores grandes; compárese con MAE para detectar sensibilidad a fallos extremos. \\
MAPE & Error porcentual respecto del observado. & Facilita una lectura relativa, pero no debe usarse con valores reales cero o casi cero. \\
sMAPE & Error porcentual simétrico. & Reduce parte de la asimetría de MAPE; se interpreta junto con las demás métricas. \\
MASE & Error escalado frente a una referencia naive interna. & Menor que 1 indica mejora frente a esa referencia; mayor que 1 exige cautela, incluso con MAPE bajo. \\
Sesgo & Promedio de errores con signo. & Un valor alejado de cero indica tendencia sistemática a sobrestimar o subestimar. \\
Intervalo de predicción & \emph{No se entrega en esta versión.} & La aplicación publica el pronóstico puntual sin una banda que lo acote. La incertidumbre se lee con las métricas fuera de muestra de esta misma tabla y con el número de ventanas que las sostiene. Es una limitación declarada, no una omisión de la interfaz. \\
Estado del horizonte & Si el horizonte se entrega y con qué evidencia lo hace. & Un horizonte se niega solo por imposibilidad: pronóstico no finito o \(W=0\). Cualquier otra magnitud se publica como advertencia y no impide el uso en empalme. \\
\bottomrule
\end{tabularx}
\end{table}

\subsection{Lectura de las pestañas de resultados}

El panel de resultados resume el periodo, valor, modelo y estado del horizonte. \textit{Serie histórica} permite comprobar los valores observados; \textit{Gráfica} muestra la relación entre historia, proyección e intervalos; \textit{Fila fuente} identifica el registro exacto del anexo; \textit{Explicación} traduce las métricas y advertencias; \textit{Diagnóstico} presenta controles de datos y residuos; y \textit{Metodología} documenta la preparación, comparación y selección. Estas pestañas deben revisarse como una unidad: ninguna cifra aislada reemplaza la trazabilidad completa.

Cuando el empalme solicita una fecha posterior al último dato real, estas mismas pestañas se actualizan con la proyección activada desde el empalme. El usuario debe confirmar que la ruta y el periodo proyectado coincidan con el cálculo antes de aceptar el registro.

\subsection{Generación de informe DOCX}

El usuario puede generar un informe cuando requiera documentar el análisis. El reporte conserva la ruta seleccionada, fila fuente, rango de datos, modelo, métricas por horizonte, resultado puntual, intervalos, estado del horizonte, advertencias, gráfica y metodología. Antes de compartirlo se deben verificar tres coincidencias: ruta del informe frente a la interfaz, último periodo real frente al archivo y fecha proyectada frente a la solicitud. También se recomienda registrar el nombre y versión del anexo usado.

\subsection{Uso del módulo de empalme ICCP--ICOCIV}

\subsubsection{Selección del periodo de análisis}

El usuario define fecha inicial y fecha final mediante los campos de año y mes. La aplicación identifica si el cálculo cae antes de diciembre de 2021, después de diciembre de 2021 o cruza el periodo de transición. Si la fecha final supera el último dato \ICOCIV{} real, el sistema debe usar el módulo de proyección existente para estimar el índice final, siempre que el horizonte sea viable.

\subsubsection{Selección de serie ICCP}

El usuario selecciona el tipo de serie \ICCP{} y luego la serie correspondiente. Las opciones válidas de tipo son Total \ICCP, Canasta general y Grupo de obra. La lista de series cambia según el tipo seleccionado. Si el usuario cambia el tipo, debe seleccionar nuevamente la serie para evitar cálculos con una combinación inválida.

\subsubsection{Selección de ruta ICOCIV equivalente}

El usuario selecciona la ruta \ICOCIV{} equivalente desde el selector jerárquico. Esta equivalencia debe responder al criterio técnico del caso. La aplicación registra la ruta, pero no decide automáticamente si una equivalencia contractual es procedente.

\subsubsection{Ingreso de unidad y valores P y A}

El usuario ingresa el insumo contractual, la unidad, \(P\), \(A\) y los demás datos solicitados por el formulario. La unidad se selecciona desde la lista cerrada: m lineal, m2, m3, kg, Unidad o Global. La aplicación bloquea el cálculo cuando la opción visible es ``Sin selección''.

\subsubsection{Cálculo general de actualización}

La aplicación calcula \(R_1\), \(R_2\), \(R\) y valor actualizado. En el cálculo general, \(R_1\) corresponde al ajuste parcial con \ICCP, \(R_2\) al ajuste parcial con \ICOCIV{} y \(R\) a la suma de ambos. El resultado principal debe revisarse junto con las tablas inferiores, porque allí se muestran índices, fechas, ruta y fórmulas.

El usuario debe verificar que \(R=R_1+R_2\) y que el valor actualizado corresponda a la base ajustable más el ajuste total. Si el índice \ICOCIV{} final fue proyectado, la lectura debe conservar esa advertencia porque el valor proviene del módulo estadístico y no de una publicación oficial nueva.

\subsubsection{Cálculo especial para acero}

El usuario ingresa \(P_0\), \(I_x\) y \(q\) cuando aplique el caso especial para acero. Este flujo calcula los ajustes parciales y el valor adicional por fluctuación, manteniendo separadas las variables del acero frente al cálculo general.

En este caso, \(Z\) representa la comparación entre el valor facturado del acero \((I_x \times q)\) y el valor ofertado ajustado \((P_0+R)\). Por eso el usuario debe revisar \(P_0\), \(I_x\), \(q\), \(R_1\), \(R_2\), \(R\) y \(Z\) en el detalle técnico antes de exportar.

\subsubsection{Interpretación de tablas de equivalencia y cálculo}

Las tablas muestran la equivalencia seleccionada, índices usados y resultados. La primera tabla conserva el insumo contractual, el grupo o serie \ICCP{} equivalente y el insumo o ruta \ICOCIV{} equivalente. La segunda tabla conserva unidad, precio base, índices iniciales y finales, \(R_1\), \(R_2\), \(R\) y valor actualizado. Si el índice \ICOCIV{} final fue proyectado, debe quedar identificado como proyectado.

\subsubsection{Exportación a Excel}

El usuario exporta el resultado para revisión y archivo. Antes de exportar, debe comprobar que existan cálculos válidos y que no haya campos incompletos. El exportable organiza información general, trazabilidad, equivalencias, cálculo, metodología/fórmulas y detalle de acero como secciones consecutivas dentro de una sola pestaña, conforme al requisito RF-10.

\subsection{Errores frecuentes y recomendaciones de uso}

Errores frecuentes incluyen no seleccionar unidad, no seleccionar serie \ICCP, usar una ruta \ICOCIV{} sin índices, solicitar periodos sin datos, pedir una proyección fuera del horizonte viable, confundir una serie de Canasta general con un Grupo de obra o interpretar los factores de índice como valores monetarios. Para reducirlos se recomienda revisar la ruta y la fila fuente, confirmar unidad y fechas, verificar que \(R=R_1+R_2\), leer las advertencias del módulo de proyección y archivar el exportable junto con el anexo y el soporte contractual. Mientras CP-EXP-01 permanezca abierto, también debe registrarse que el libro tiene cinco pestañas y no la única hoja exigida.

\section{DISCUSIÓN Y CONCLUSIONES}

\subsection{Discusión de resultados}

Los resultados del desarrollo muestran que la herramienta logró integrar en una sola aplicación procesos que antes podían quedar distribuidos entre archivos, hojas de cálculo y revisiones manuales. La carga de anexos \ICOCIV, el selector jerárquico, la construcción de series, el análisis estadístico, la proyección y el empalme \ICCP--\ICOCIV{} quedaron organizados como flujos conectados, pero separados por responsabilidad. Esta separación facilita mantenimiento y permite revisar cada resultado desde su fuente.

El módulo de proyección aporta valor porque no entrega únicamente un índice estimado; también muestra el comportamiento histórico, el modelo seleccionado, las métricas, las advertencias y el horizonte viable. El tratamiento del patrón de cambio de año ilustra el criterio general adoptado, y lo ilustra por la vía más exigente: detectar un comportamiento con estadística robusta no basta para tratarlo. El patrón está medido y es contundente, pero ningún componente del ajuste ---ni el estimador del salto futuro, ni la forma funcional del factor, ni las puertas de activación calibradas sobre el propio anexo--- resistió la verificación de sustento, de modo que el tratamiento se retiró y el fenómeno se publica como diagnóstico. Declarar la limitación fue preferible a conservar un ajuste que mejoraba el error en algunas series sin un método publicado que lo respaldara: la mejora empírica no es, por sí sola, fundamento metodológico. El módulo de empalme complementa ese flujo al aplicar las fórmulas de transición revisadas en los lineamientos, conservar la selección de series \ICCP{} e \ICOCIV{} y presentar tablas acumulativas para varios insumos.

\subsection{Cumplimiento de objetivos}

El objetivo general se cumplió mediante el desarrollo de una aplicación orientada al análisis de variaciones y proyecciones de precios en obras civiles con base en información oficial del \ICCP{} y del \ICOCIV. Los objetivos específicos se evidencian en la implementación de consulta, filtrado, visualización y procesamiento de datos; en las funcionalidades de análisis y proyección de series históricas; en la evaluación de enfoques estadísticos y métricas de validación; en las pruebas funcionales realizadas con información representativa; y en la elaboración del presente manual de usuario junto con la documentación técnica del sistema.

\subsection{Aportes del software desarrollado}

El principal aporte es integrar consulta, análisis, proyección, empalme y trazabilidad en una sola aplicación de apoyo técnico.

\subsection{Utilidad para la Secretaría de Infraestructura de la Gobernación del Cauca}

La herramienta puede apoyar revisiones técnicas de variaciones de precios, organización de índices y generación de reportes.

\subsection{Limitaciones encontradas}
\label{sec:limitaciones_encontradas}

\subsubsection*{Limitación metodológica central: el primer origen de la validación temporal}

La limitación más importante que este trabajo identificó no está en los datos
sino en el propio diseño de la validación. Conviene separar dos cosas que suelen
leerse juntas.

\textbf{Lo que sí está sustentado.} El \emph{esquema} de validación empleado
---origen móvil con ventana expansiva, reestimación del modelo en cada origen y
ausencia de fuga de información futura--- está respaldado por
\textcite{hyndman2021fpp3} §5.10 y por \textcite{tashman2000}. También lo están la
restricción a la muestra común de pares \((t,h)\), que impide que un candidato
obtenga ventaja por disponer de menos observaciones, y la selección por menor
error cuadrático fuera de muestra.

\textbf{Lo que no lo está.} \emph{Desde qué origen} debe comenzar esa
validación. Ninguna de las fuentes consultadas prescribe un tamaño inicial de
entrenamiento para un esquema de origen móvil: \textcite{hyndman2021fpp3} §5.10 y
\textcite{tashman2000} describen el procedimiento sin fijar su inicio, y la única
proporción del libro ---§5.8, «about 20\,\%» de conjunto de prueba--- corresponde
a una partición única y no a un origen móvil. La literatura econométrica de
evaluación de pronósticos formaliza la ventana inicial \(R\) y la razón
\(\pi = P/R\), pero las \emph{asume} dadas: no las prescribe, y los valores
críticos asintóticos dependen de \(\pi\), lo que confirma que la elección afecta
a la inferencia sin estar determinada por la teoría.

\textbf{El valor vigente es provisional.} La aplicación usa \(N_0 = 6\), que
corresponde al máximo de los mínimos operativos codificados para intentar el
ajuste de cada candidato ---seis por Holt amortiguado, que tiene cinco
parámetros---. Ese número satisface una condición \emph{necesaria} de
comparabilidad, pero no constituye un mínimo de identificabilidad estadística
demostrado: contar parámetros no garantiza que la muestra los identifique. Una
inspección empírica del ajuste de ese modelo muestra además soluciones
frecuentes en la frontera de la caja de parámetros incluso al aumentar el tamaño
de entrenamiento, de modo que esa vía tampoco aporta una base para derivar un
\(N_0\) por identificabilidad. La observación es diagnóstica y compatible con la
limitación; no constituye una demostración de no identificabilidad.

\textbf{La elección es material, y está medida.} Sobre diez series y cincuenta y
nueve combinaciones de series y primeros orígenes, el modelo que la aplicación
selecciona \textbf{cambia en seis de las diez series y en once de las cincuenta y
nueve combinaciones}. En una de ellas la variación no es monótona, de modo que
tampoco existe un \(N_0\) «suficientemente grande» a partir del cual la selección
se estabilice. Como la selección determina el modelo que produce toda la
trayectoria publicada, el primer origen no es un detalle de implementación.

\textbf{Qué se hizo con esa evidencia.} Nada más que declararla. La medición de
sensibilidad documenta que la decisión \emph{importa}; no puede usarse para
\emph{tomarla}, porque elegir el valor que produce el resultado preferido es
precisamente la práctica que la literatura advierte como \emph{data snooping} y
que este trabajo se prohibió desde el principio. Por eso la aplicación reporta la
evidencia fuera de muestra como \textbf{provisional} en cada resultado, el
criterio correspondiente figura como \emph{pendiente de decisión} en la tabla
auditable, y no se impuso ninguna regla no sustentada para aparentar un cierre.

Resolver esta cuestión exigiría una fuente que prescriba el tamaño inicial ---que
no se localizó--- o un rediseño metodológico mayor hacia procedimientos robustos
a la elección de la ventana de estimación, que redefiniría el criterio de
selección y excede el alcance declarado de este trabajo.

\subsubsection*{Otras limitaciones}


Las principales limitaciones se relacionan con la disponibilidad y calidad de los datos de entrada. Una serie corta, discontinua o con cambios abruptos reduce la confiabilidad de cualquier proyección. Además, la equivalencia entre un insumo contractual, una serie \ICCP{} y una ruta \ICOCIV{} no siempre es única, por lo que debe ser revisada por el profesional responsable. En términos operativos, la aplicación depende de que los anexos cargados conserven una estructura compatible con los formatos esperados por el sistema. A ello se añade una limitación propia de la trayectoria mensual. La versión final identifica de forma descriptiva el patrón recurrente que algunas series \ICOCIV{} presentan en la transición de diciembre a enero, pero \textbf{no incorpora ninguna corrección sobre los valores proyectados}: el tratamiento evaluado durante el desarrollo fue retirado porque su formulación no alcanzó el sustento metodológico exigido. En consecuencia, en las series donde ese patrón sea material, la trayectoria puntual puede no reproducir por completo la distribución temporal histórica alrededor del cambio de año: el crecimiento se reparte según el modelo seleccionado y no se reconcentra en enero. El fenómeno se mide y se informa; corregirlo exige un procedimiento sustentado que esta versión no adopta.

\subsection{Conclusiones}

El proyecto permitió construir una herramienta que organiza información oficial, reduce operaciones manuales y mejora la trazabilidad de cálculos asociados a variación y proyección de precios en obras civiles. La aplicación no reemplaza el criterio técnico del profesional, pero sí ofrece un entorno más controlado para consultar series, evaluar modelos, calcular empalmes y generar evidencias.

La separación entre análisis/proyección \ICOCIV{} y empalme \ICCP--\ICOCIV{} resultó adecuada porque cada proceso tiene entradas, validaciones y salidas diferentes. Al mismo tiempo, la integración entre ambos módulos permite que el empalme use el índice \ICOCIV{} real o proyectado según corresponda, sin duplicar metodologías.

\textbf{Sobre el rigor metodológico alcanzado y sus límites.} El trabajo sometió
la capa estadística a una revisión sistemática que identificó y retiró reglas
decisorias sin respaldo: umbrales y proporciones sin fuente en el diseño del
backtesting, una escalera de confianza sin sustento, un tope de horizonte
arbitrario y una atribución bibliográfica que no resistía la verificación. Dos
resultados merecen enunciarse por separado, porque son de naturaleza distinta.

El primero es que \textbf{el intervalo de predicción se retiró de todas las
salidas}. Ninguno de los métodos evaluados alcanzó una construcción completa
sustentada, y publicar sus límites habría afirmado una precisión que la
aplicación no puede defender. El pronóstico puntual se conserva porque una
deficiencia del intervalo no invalida por sí sola el valor puntual, y la
aplicación declara expresamente que la incertidumbre no viene acotada. No se
validó el intervalo: se retiró.

El segundo es que \textbf{el primer origen de la validación temporal quedó sin
determinar}, según se detalla en las limitaciones. El trabajo optó por declarar
esa limitación y medir su efecto en lugar de imponer un valor arbitrario que le
diera apariencia de cierre. Esa decisión es coherente con el alcance declarado,
que excluye expresamente crear una metodología estadística oficial nueva.

En consecuencia, y aunque la herramienta cumple los objetivos planteados y su
comportamiento es reproducible y auditable, \textbf{no corresponde afirmar que su
metodología esté cerrada}. Lo que este trabajo entrega es un procedimiento
implementado con rigor, evaluado de forma reproducible y con sus limitaciones
identificadas y publicadas.

\subsection{Recomendaciones}

Se recomienda mantener actualizados los anexos oficiales, validar técnicamente cada equivalencia \ICCP--\ICOCIV{}, y ejecutar la batería completa en el entorno de entrega antes de la socialización institucional.

\subsection{Trabajos futuros}

En el plano metodológico, el trabajo futuro más definido es el \textbf{tratamiento del patrón de cambio de año}. Esta versión lo identifica y lo describe, pero no lo corrige, porque el procedimiento evaluado durante el desarrollo no alcanzó el sustento exigido. Queda pendiente estudiar un tratamiento respaldado por literatura verificable, integrado \emph{dentro} del esquema de validación con ventana expansiva ---de modo que su aporte se mida con evidencia fuera de muestra y no sobre el mismo anexo con que se calibraría---, y comparado contra la trayectoria sin corregir. No debe abordarse reactivando el factor retirado ni introduciendo cortes de activación ad hoc: la exigencia es que la formulación, sus parámetros y su regla de uso estén sustentados antes de que ninguna corrección alcance un valor publicado.

Con el mismo criterio queda abierta la construcción de un \textbf{intervalo de predicción} cuya formulación completa ---supuestos, nivel nominal, grados de libertad, dispersión, tratamiento por horizonte y evaluación de cobertura--- esté sustentada, y la resolución de la elección del primer origen de la validación temporal, hoy declarada como limitación.

Como trabajos futuros se recomienda ampliar la batería de pruebas automatizadas, incorporar evidencias visuales finales al manual, fortalecer la gestión de versiones de anexos oficiales y mejorar la trazabilidad de exportaciones. También podría evaluarse la incorporación de perfiles de usuario o mecanismos de auditoría cuando la herramienta se use en flujos institucionales con varios revisores.

\newpage
\section{REFERENCIAS}

\bibliography{referencias}

\newpage
\section{ANEXOS}

\subsection{Anexo A. Fuentes de datos}

Incluye anexos \ICOCIV{} del \DANE, Anexo 10 \ICCP{} histórico y documentos técnicos disponibles en el proyecto.

\end{document}
